\documentclass[11pt]{article}
\usepackage[margin=.85in]{geometry}
\usepackage{amsmath,amssymb,mathtools,mathrsfs,hyperref}
\hypersetup{colorlinks=true,linkcolor=blue,urlcolor=blue}
\newcommand{\scri}{\mathscr I^+}
\newcommand{\TF}{\mathrm{TF}}
\newcommand{\dd}{\partial}
\newcommand{\Tr}{\operatorname{tr}}
\title{Regular scri closure audit for the elliptic-shift hyperboloidal Z4c hybrid}
\author{Linear derivation, closure theorem and verification record}
\date{25 October 2026}
\begin{document}\maketitle
\noindent\textbf{Current result.} Section~\ref{sec:assembled-linear-closure} gives the full coupled linear scri closure and fixed-phase evolution for the repaired formulation: original reference lapse, elliptic metric-Dirac shift, an added boundary clock $\delta\alpha|_{\scri}=0$, the disclosed Appendix-C scalar equation, and the stated damping parameters. The theorem uses an explicit compatible nonlocal quotient graph norm, with preserved constraint falloffs and regular pole traces. It is not an arbitrary-raw-Sobolev, nonlinear or discretization theorem. The user selected the full linearized scope. Earlier candidates, failed closures and superseded progress statements are retained as a research record; the assembled theorem and final certificates specify the final result.
\tableofcontents

\section{Objective and completion requirements}
The objective is a regular actual-scri closure for free hyperboloidal Z4c with a metric-derived elliptic shift and a reference-reduced live lapse, preserving constraint falloffs and satisfying a full coupled boundary estimate. The target is to stay as close as possible to arXiv:1412.4801 apart from shift and possibly lapse conditions. Completion requires all of:
\begin{enumerate}
\item An explicit formulation and mapping to the target's final equations, with every bulk constraint addition disclosed.
\item A state space and regular limiting equations including freely evolving differential constraint errors.
\item Invariance of the asymptotic state space, not just vanishing pole numerators initially.
\item All scalar, vector, tensor and angular sectors, including lapse and elliptic boundary reconstruction.
\item A frequency-uniform coupled estimate in a stated norm; no-growing-mode tests alone do not suffice.
\item An explicit statement whether the theorem is linearized or nonlinear. A linear theorem must not be described as nonlinear well-posedness.
\end{enumerate}

\section{What the reference paper establishes}
The six-page black-hole paper \url{https://arxiv.org/html/1412.4801} uses the final equations in Appendix C of \url{https://arxiv.org/html/1412.3827}. It demonstrates stable spherical evolutions and convergence, using a staggered grid that excludes scri and extrapolates values there. Section 7 of the longer paper expressly says no rigorous continuum well-posedness or numerical stability proof is supplied. The requested coupled theorem is therefore stronger than the evidence established in the reference.

The final lapse's $\Delta K$ maps to the physical $\hat K$ deviation:
\[
\Delta K=K_{\rm phys}-2\Theta_{\rm phys}-K_*=p-p_*.
\]
It must not be silently reinterpreted as a physical-$K$ driver, which has the interior harmonic-speed Jordan defect found in the preceding derivation.

There is a specific difference between the repository tensor equation and Appendix C's scalar-constraint equation. Let
\[
T=\Theta_{\rm phys},\quad w=\bar n(\Omega),\quad
d_\kappa=\kappa_1(2+\kappa_2),\quad z_i=(\Lambda_i-V_i)/2.
\]
At $C_{\rm Z4c}=0$, the repository equation is
\[
T_t=\beta^i\dd_iT+\frac{\alpha H}{2\Omega}
+\Omega\alpha\chi\widetilde D_i z^i
-\frac{\alpha(3w+d_\kappa)}{\Omega}T .
\]
The corresponding Appendix C treatment has no $3w$ contribution:
\[
T_t=\beta^i\dd_iT+\frac{\alpha H}{2\Omega}
+\Omega\alpha\chi\widetilde D_i z^i
-\frac{\alpha d_\kappa}{\Omega}T .
\]
Matching this term requires the explicit constraint addition
\[
\boxed{\Delta T_t=3\alpha wT/\Omega.}
\]
It vanishes at $T=0$ but changes the off-constraint system. The paper is not explicit about the intent behind this particular source-equation difference; it is not attributed here to an undocumented stabilization decision.

\section{Background and exact first-jet obstruction}
Use
\[
\Omega=\frac{S^2-r^2}{2aS},\quad
A=\frac{S^2+r^2}{2aS},\quad B=-x/a,\quad
G=\nabla\Omega=-x/(aS),\quad K_*=-3/a.
\]
The exact first variations and their pole split $U_t=R_U+S_U/\Omega$ are in the companion linearization note.

Choose a smooth toroidal vector $\zeta=f(r)(e\times x)$, with $f$ nonzero at scri. Set every metric, curvature, lapse, shift, trace and scalar-constraint perturbation to zero initially, but choose
\[
z_i=\Omega\zeta_i,\qquad \delta\Lambda_i=2\Omega\zeta_i.
\]
Then $V=H=M=T=0$, the elliptic shift is zero, scri nullness holds, and every initial pole numerator vanishes. Let
\[
E_{ij}=G_i\zeta_j+G_j\zeta_i,\qquad \Tr E=0.
\]
Direct substitution gives at scri
\[
a_{ij,t}=3AE_{ij},\qquad
z_{i,t}=A(3/a-\kappa_1)\zeta_i.
\]
The curvature pole's time derivative is
\begin{equation}\boxed{
(S_{a,ij})_t=-2A^2\kappa_1 E_{ij}.
}\end{equation}
For constant $\kappa_1>0$ this is nonzero. Thus initially residue-compatible smooth data do not define an invariant classical smooth state space. This rotational example is unaffected by the scalar-constraint correction above and by a lapse modification preserving the scalar-free invariant sector.

This is an obstruction to a classical conformally smooth closure with those allowed initial jets. It is not a theorem excluding weaker weighted semigroup solutions with corner layers or a more restrictive domain.

\section{Scaling damping alone still fails: exact rotational sector}
Let $\mathcal R=e\times x$ and $\mathcal T_{ij}=x_i\mathcal R_j+x_j\mathcal R_i$. Consider
\[
h_{ij}=H(r)\mathcal T_{ij},\quad
a_{ij}=P(r)\mathcal T_{ij},\quad z_i=Z(r)\mathcal R_i,\quad
v_i=W(r)\mathcal R_i .
\]
All scalar perturbations vanish. One finds
\[
V_i=(rH'+5H)\mathcal R_i.
\]
With $V=0$ and regularity at the center, coordinate-time freezing implies $H=0$. The elliptic shift is reconstructed from
\[
W'=2ArP,\qquad W(S)=0.
\]
There is no lapse/Theta freedom to repair the remaining rotational subsystem.

Set $\kappa_1=\Omega\bar\kappa$, where $\bar\kappa\geq0$ is smooth in a scri collar. The exact reduced equations are
\begin{align}
P_t={}&-\frac{rP'+2P}{a}+\frac A r Z'
+\frac A\Omega\left[-\frac{2Z}{aS}+(K_*-w_*)P\right],\\
Z_t={}&-\frac{rZ'+3Z}{a}+A(rP'+5P)
+\frac A\Omega\left[\frac{2r^2P}{aS}+\frac{2Z}{a}\right]-A\bar\kappa Z .
\end{align}
The necessary scri relation is $C=P+Z/S=0$. Its exact evolution is
\[
C_t=-\frac{r\Omega}{S}C'
+\frac{r^2-5S^2}{2aS^2}C
+\frac{2aA\Omega}{rS}Z'
-\frac A S(1/S+\bar\kappa)Z .
\]
Hence the relation is not invariant. In particular, $P=0$, $Z=\Omega f(r)$ satisfies $C=C_t=0$ at scri initially, but
\[
Z_t(S)=3Af(S)/a,\qquad
C_{tt}(S)=-\frac{3A}{aS}(1/a+A\bar\kappa)f(S)\ne0.
\]
This shows why merely replacing each $1/\Omega$ quotient by its first limiting derivative is not a closure.

\section{A concrete finite constraint addition repairs this sector}
Keep $\kappa_1=\Omega\bar\kappa$ in the collar and add to the \emph{regular} trace-free curvature equation
\begin{equation}\boxed{
\Delta R_{a,ij}=-(1+aA\bar\kappa)
(z_iG_j+z_jG_i)^{\TF}.
}\end{equation}
This is lower differential order in the interior and vanishes when the Z4 connection constraint vanishes. It is a bulk constraint addition beyond a gauge change and must be disclosed as such.

In the rotational sector it changes the residual equation to
\[
C_t=-\frac{r\Omega}{S}C'
+\frac{r^2-5S^2}{2aS^2}C
+\frac{2aA\Omega}{rS}Z'
+\frac{\Omega}{S^2}Z.
\]
Now set $C=\Omega Y$. The exact system is
\begin{align}
Y_t={}&-\frac{r\Omega}{S}Y'
+\frac{2aA}{rS}Z'
+\frac{3r^2-5S^2}{2aS^2}Y+\frac{Z}{S^2},\\
Z_t={}&Ar\Omega Y'
-r\left(\frac1a+\frac AS\right)Z'
+\frac{A(5S^2-3r^2)}{2aS}Y
-\left(\frac3a+\frac AS+A\bar\kappa\right)Z .
\end{align}
All coefficients are smooth to scri on an outer collar $r>0$. The principal matrix has
\[
\Tr M=-2r/a,\qquad\det M=-\Omega^2,\qquad
\lambda_\pm=-r/a\pm A.
\]
The eigenvalue gap is $2A>0$, so a smooth uniformly positive symmetrizer exists. At scri the coordinate propagation speeds are $0,2A$; there is no incoming mode. The boundary equations are ordinary limits of the regular system:
\begin{align}
Y_t|_{\scri}&=\frac2S Z'-\frac Ya+\frac Z{S^2},\\
Z_t|_{\scri}&=-2A Z'+A^2Y-(4/a+A\bar\kappa)Z .
\end{align}
This gives a regular variable-coefficient linear sector, rather than just a frozen symbol test. It preserves $C=O(\Omega)$ and the normal-jet compatibility hierarchy through the regular equations.

The rotational Hamiltonian and scalar constraints stay zero. The momentum condition is not $M_i=0$ everywhere. Its coefficient satisfies
\[
M_i+\frac2a z_i=\Omega\,\mathcal N(r)\mathcal R_i,
\]
\[
\mathcal N=\Omega rY'-\frac rS Z'
+\left[\frac{r^2}{aS}+5\Omega\right]Y-\frac ZS .
\]
Thus a regular combination of momentum and connection constraints is preserved while allowing independent constraint errors.

\textbf{Scope:} no result for higher odd harmonics, scalar/even sectors, or the full coupled lapse/elliptic boundary estimate is inferred from this sector. The proposed tensor addition remains a candidate in those sectors.

\section{Extension to every linear odd-parity harmonic}
An independent audit reproduces the rotational repair and extends it to
all odd harmonics on a collar. Let $L=\ell(\ell+1)$, $\Lambda=L-2$,
$D^AX_A=0$, $X_{AB}=D_{(A}X_{B)}$ and
$D^BX_{AB}=-\Lambda X_A/2$. Use spherical covariant components
\[
h_{rA}=fX_A,\ h_{AB}=gX_{AB},\
a_{rA}=pX_A,\ a_{AB}=qX_{AB},\ z_A=zX_A .
\]
Define
\[
F=rf+g-rg'/2,\qquad
Y_1=(p+rz/S)/\Omega,\qquad Y_2=(q-F/S)/\Omega .
\]
The metric connection memory satisfies
$f'+2f/r-\Lambda g/(2r^2)=V$, with $V_t=-\eta V$.
The exact repaired equations, independent of this prescribed memory, are
\begin{align}
(Y_1)_t={}&-\frac{r\Omega}{S}Y_1'+\frac{2aA}{S}z'
-\frac{Y_1}{a}-\frac{2S^2+r^2}{rS^2}z
-\frac{\Lambda A}{2rS}Y_2+\frac{\Lambda aA}{r^3S}F,\\
z_t={}&A\Omega Y_1'-r(1/a+A/S)z'
+\frac{AS}{ar}Y_1-\frac{\Lambda A\Omega}{2r^2}Y_2
-\frac{\Lambda A}{2r^2S}F+(-1/a+A/S-A\bar\kappa)z,\\
F_t={}&-\frac{r\Omega}{S}F'+rA\Omega Y_2'
+F/a-(A^2+\Omega^2)Y_2-2Ar\Omega Y_1+2Ar^2z/S,\\
(Y_2)_t={}&-r(1/a+A/S)Y_2'+\frac{2aA}{rS}F'
+(1/a-2A/S)Y_2+\frac{2Ar}{S}Y_1+\frac{4aA}{S}z-F/r^2 .
\end{align}
For $\ell=1$, $Y_1=r^3Y$, $z=r^2Z$, the first pair reproduces
the rotational system above.

Let $U=(Y_1,z)$, $W=(F,Y_2)$. Their radial matrices are
\[
M_U=\begin{pmatrix}-r\Omega/S&2aA/S\\ A\Omega&-r(1/a+A/S)\end{pmatrix},
\quad
M_W=\begin{pmatrix}-r\Omega/S&rA\Omega\\
2aA/(rS)&-r(1/a+A/S)\end{pmatrix}.
\]
Angular couplings are $U_t=\Lambda BW+\cdots$, $W_t=CU+\cdots$:
\[
B=A\begin{pmatrix}a/(r^3S)&-1/(2rS)\\
-1/(2r^2S)&-\Omega/(2r^2)\end{pmatrix},\quad
C=A\begin{pmatrix}-2r\Omega&2r^2/S\\2r/S&4a/S\end{pmatrix}.
\]
They satisfy $BC=CB=-A^2I/r^2$ and
$(M_U+rI/a)B+B(M_W+rI/a)=0$.
After $\widehat W=i\sqrt{\Lambda}W$, define
\[
R=(\operatorname{diag}(M_U,M_W)+rI/a)/A,\qquad
T=\frac rA\begin{pmatrix}0&-B\\ C&0\end{pmatrix}.
\]
Then $R^2=T^2=I$, $RT+TR=0$. The positive symmetrizer
\[
H=I+R^TR+T^TT+(RT)^T(RT)
\]
is smooth and uniformly bounded on a collar through scri. It symmetrizes
all radial/angular principal directions. The eigenvalues are
\[
-\frac ra k_r\ \pm A\sqrt{k_r^2+k_{\rm ang}^2/r^2}
\]
with multiplicity two. All remaining coefficients are bounded uniformly
in $\ell$ in the norm of
$(Y_1,z,\sqrt\Lambda F,\sqrt\Lambda Y_2)$. The scri radial flux is
nonpositive; inner-collar influx must be controlled by the interior.

The odd shift is reconstructed from a homogeneous-Dirichlet vector
Laplacian. Metric reconstruction satisfies
\[
g''-\frac{L}{r^2}g=2V-\frac{2F'}r-\frac{2F}{r^2}.
\]
The apparent boundary forcing $F(S)$ cancels after integrating this
equation against $r^{\ell+1}$. The dynamic boundary trace obeys
\[
(g_B)_t+\frac{\ell-1}{a}g_B
=-\frac2a S^{-\ell}
\left[\int_0^S r^{\ell+1}V\,dr+
(\ell-1)\int_0^S r^{\ell-1}F\,dr\right].
\]
The functionals have bounds $C\ell^{-1/2}\|V\|$ and
$C\ell^{1/2}\|F\|$, and the boundary damping is $O(\ell)$.
This avoids a pointwise $F$ trace loss in odd reconstruction.
The $\ell=1$ rotation is neutral and determined by initial gauge data.
This is a uniform odd-sector collar result; a full global
origin-to-scri matching theorem and the even sector are not proved.

\section{The same candidate fails the even scalar test}
For spherical Dirac data $h=0$, put
\[
a_{ij}=P(n_in_j-\delta_{ij}/3),\quad z_i=Zn_i,\quad v_i=Vn_i .
\]
The shift obeys $V'-V/r=AP$ and nullness
$V(S)=-u-Ac/2$. The radial curvature pole condition is
\[
E=P+c'+2Z=0\quad\hbox{at scri}.
\]
Use the Appendix-C scalar correction and the odd repair above.
Choose smooth initial perturbations
\[
c=u=V=P=Z=0,\qquad T=\Omega^2\tau(r),\qquad q=-2T .
\]
Then physical $k=q+2T=0$, $H=M=0$, and all pole numerators vanish.
The initial lapse RHS is regular, $u_t=2F_{\rm lapse}\Omega\tau$;
$c_t=P_t=V_t=0$. But the exact connection constraint equation gives
\[
Z_t=A T'/\Omega\ \longrightarrow\ -2A\tau/a.
\]
Consequently
\begin{equation}\boxed{E_t(S)=-4A\tau(S)/a\ne0.}\end{equation}
The odd repair is therefore \emph{not} an all-sector closure. Its
success must not be promoted to a full-system result.
A repair must couple the weighted normal scalar-constraint jets to
curvature/connection regularity. Merely increasing the initial power
of $\Omega$ moves this obstruction to a later compatibility order.

\section{Boundary symbol versus a full estimate}
In a normalized constant half-space model $x<0$, take $A=1$, $B^x=-1$, $f=1$. Wave modes satisfy
\[
(s+\lambda)^2=\lambda^2-k^2,\qquad
\lambda=-\frac{s^2+k^2}{2s}.
\]
For $\Re s>0$, these do not decay toward $x=-\infty$. The remaining decaying spatial gauge modes satisfy $L\xi=0$.

For tangential frequency $k>0$ in the $y$ direction, their Dirichlet-to-Neumann relation is
\[
\xi_x'=\frac{k}{7}(6\xi_x-i\xi_y),\quad
\xi_y'=\frac{k}{7}(8\xi_y-i\xi_x),\quad \xi_z'=k\xi_z .
\]
Null-normal shift data and zero tangential shift give boundary matrix
\[
\begin{pmatrix}
s&0&0\\-ik/7&s+8k/7&0\\0&0&s+k
\end{pmatrix},
\qquad
\det=s(s+8k/7)(s+k).
\]
There is no exponentially growing boundary eigenmode. The factor $s=0$ is normal boundary-displacement memory.

However, grazing constraint waves can force this neutral displacement and violate an ordinary no-loss energy estimate. Take scalar constraint eigenvectors
\[
(Q,m,\theta,z)=(1,1/3,1,0),\quad p=\partial_nu=h=\Gamma=0.
\]
For wave vector $(K_x,k,0)$, let
\[
K=\sqrt{K_x^2+k^2},\quad \mu=K_x/K,\quad \omega=K-K_x.
\]
Then $c=1/(iK)$ and $v_{\rm wave}=n/(2iK)$ times the oscillatory phase. The normal gauge displacement satisfies
\[
q_t=-\frac{1+\mu}{2iK}e^{i\omega t}.
\]
For $k=m\to\infty$, $K_x=m^2\rho$, $1<\rho<2$, a smooth wavepacket has initial normalized volume norm $O(m^{-1})$. Its boundary displacement at fixed positive time is $O(m^{-2})$, but the elliptic extension has width $m^{-1}$ and metric-gradient amplitude $O(1)$. Its norm is therefore $O(m^{-1/2})$, an amplification of order $\sqrt m$.

This construction concerns the \emph{normalized model}. Its data are not a solution satisfying the full conformal pole hierarchy. It does not disprove every regularized weighted closure. It does show that a no-growing-boundary-mode calculation cannot replace the requested coupled estimate and that the constraint/gauge trace interaction must be handled explicitly.

\section{Closed subsidiary identities and further work}
A useful exact linear constraint route uses the physical background metric $\gamma^p=\Omega^{-2}\delta$, lapse $N=A/\Omega$, mean curvature $K_*=-3/a$, and its covariant derivative $D^p$. Let $\epsilon=0$ denote the repository scalar equation and $\epsilon=1$ the Appendix-C scalar correction. Define
\[
d_K=6(1-\epsilon)w_*+3\kappa_1(1+\kappa_2),
\]
\[
\mathcal Q_{ij}=2N D^p_{(i}z_{j)}+\frac{\gamma^p_{ij}}3 r_Q,
\qquad
r_Q=2A G\cdot z-Nd_KT .
\]
Then $\Tr\mathcal Q=2\Omega A\,\partial_i z_i-Nd_KT$.
The linear physical ADM constraint identities give the closed candidate subsidiary system
\begin{align}
(\partial_t-\mathcal L_B)H={}&-2N D_p^iM_i-4M^iD_i^pN+2NK_*H
+\frac{4K_*}{3}\Tr\mathcal Q,\\
(\partial_t-\mathcal L_B)M_i={}&-\frac N2 D_i^pH-H D_i^pN
+NK_*M_i+D_p^j\mathcal Q_{ij}-D_i^p\Tr\mathcal Q,\\
T_t={}&B\cdot\nabla T+\frac{AH}{2\Omega}
+\Omega A\,\operatorname{div}z
-\frac{A[3(1-\epsilon)w_*+\kappa_1(2+\kappa_2)]}{\Omega}T,\\
z_{i,t}={}&B\cdot\nabla z_i-2z_i/a
+\frac A\Omega[M_i+\partial_iT+(2/a-\kappa_1)z_i].
\end{align}
The ADM derivation and trace matching have been checked analytically, but a direct independent substitution into every full tensor equation remains a verification item. The finite curvature addition of Section 5 changes $\mathcal Q$ and must be included before using these identities for the repaired candidate.

The main unfinished work is an all-sector regular reduction or a suitable Fuchsian energy estimate, followed by the full elliptic/hyperbolic boundary estimate. Large singular damping can help a weighted estimate, but does not automatically supply a smooth finite-dimensional residue manifold. Scalar tuning alone is insufficient.

\section{A directly audited singular-damping residue}
For the unmodified tensor curvature equation and the Appendix-C scalar
correction, consider
\[
H=\Omega^2h,\quad M_i=\Omega m_i,\quad
T=\Omega^2\tau,\quad z_i=\Omega\zeta_i.
\]
Set $l=1/a$, $x=a\kappa_1$, $y=\kappa_2$. The leading value-variable
residue in scalar variables $(h/l,m_n/l,\tau,\zeta_n)$ is
\[
J_s=l
\begin{pmatrix}
-4&-4&12x(1+y)&8\\
0&-2&-2x(1+y)&-28/3\\
1/2&0&2-x(2+y)&-1\\
0&1&-2&3-x
\end{pmatrix}.
\]
Each transverse block on $(m_A/l,\zeta_A)$ is
\[
J_A=l\begin{pmatrix}-2&-6\\1&3-x\end{pmatrix}.
\]
The $\tau,\zeta,m$ columns were directly checked against the original
linearized tensor equations using respectively pure $T=\Omega^2\tau$,
pure $z=\Omega\zeta$, and $a=\Omega C$ perturbations, in addition to the
physical ADM subsidiary derivation.

For $x=10$, $y=-1/2$ the dimensionless scalar characteristic polynomial is
\[
\lambda^4+26\lambda^3+\frac{589}{3}\lambda^2
+\frac{1682}{3}\lambda+\frac{1232}{3}.
\]
Its Routh minors are $4544$ and $6810176/3$, both positive, and all
coefficients are positive. Thus this residue is Hurwitz; transverse
eigenvalues are $-4/a,-5/a$. Small direct matrix arithmetic independently
confirmed these polynomial coefficients and minors.

This does not prove a Fuchsian estimate. The principal wave energy sees
derivative variables. For instance, a physical orthonormal principal
reduction on $(H,M_n,\partial_nT,\partial_nz_n)$ has matrix
\[
B_c=\begin{pmatrix}
0&-2&0&0\\
-1/2&0&0&0\\
1/2&0&0&1\\
0&1&1&0
\end{pmatrix},\qquad B_c^2=I .
\]
Its positive symmetrizer is $I+B_c^TB_c$. In the weighted compactified
variables the natural fields are
\[
h,\quad m_i,\quad 2G_i\tau+\Omega\partial_i\tau,\quad
G_i\zeta_j+\Omega\partial_i\zeta_j .
\]
Their residues and reduction compatibilities differ from those of the
value-variable $J$ above. A common energy, derivative commutators and
boundary trace control remain to be established. Strong physical
singular damping is kept as a separate route because it stays closer to
the reference equations; the smooth-sector curvature repair is not
silently combined with this residue calculation.

\section{Earlier finite-order constraint candidate}
The closest-to-reference route currently under study keeps the original
tensor curvature equation (without the odd-sector finite addition), uses
the Appendix-C scalar correction, and takes in an outer collar
\[
\boxed{a\kappa_1=10,\qquad\kappa_2=-9/10.}
\]
These physical damping coefficients retain the reference's singular
$\kappa_1/\Omega$ mechanism. They can be smoothly blended to other
interior damping parameters away from scri without changing the interior
principal differential order. This is a candidate, not an adopted or
proved full formulation.

For a compatible common normal power $\Omega^s$ in
$(h/l,m_n/l,\tau,\zeta_n)$, where $H=\Omega^2h$,
$M=\Omega m$, $T=\Omega^2\tau$, $z=\Omega\zeta$, the scalar
indicial equations at these parameters are
\[
I(s)\begin{pmatrix}h/l\\m_n/l\\\tau\\\zeta_n\end{pmatrix}=0,\qquad
I(s)=
\begin{pmatrix}
1/2&0&s-9&-(s+1)\\
0&1&-(s+2)&s-7\\
s-4&2s-4&12&8(s+1)\\
s/2&s-2&-2(s+1)&-(4s+28)/3
\end{pmatrix}.
\]
For this explicitly ordered matrix,
\[
\det I(s)=\frac83(s-2)(16s-133).
\]
Eliminating the first two equations gives $(sL+K)(\tau,\zeta_n)^T=0$:
\[
L=\begin{pmatrix}26&20\\7&26/3\end{pmatrix},\qquad
K=\begin{pmatrix}-68&-28\\-6&-70/3\end{pmatrix},
\quad
\det(sL+K)=\frac{16}{3}(s-2)(16s-133).
\]
The sign here is fixed by the displayed column ordering. An earlier
opposite-order determinant was quoted with the opposite sign; the roots
and subsequent Fuchsian matrix were unchanged.

The root $s=2$ has scalar kernel
$(22,32/3,1,4/3)$ and two transverse kernels
$m_A=5l\zeta_A$. Thus the natural leading constraint traces are
\[
H=\Omega^4H_4,\quad T=\Omega^4T_4,\quad
M=\Omega^3M_3,\quad z=\Omega^3Z_3,
\]
\[
\boxed{
H_4=22lT_4,\quad Z_{3n}=\tfrac43T_4,\quad
M_{3n}=\tfrac{32}{3}lT_4,\quad M_{3A}=5lZ_{3A}.
}
\]
These are three outgoing leading amplitudes, not freely supplied
incoming boundary data. The second scalar branch has exponent $133/16$
in the base weighted variables. Its raw powers are
$z,M\sim\Omega^{149/16}$ and $H,T\sim\Omega^{165/16}$.
These powers permit finite normal Sobolev regularity; they do not imply
infinite conformal smoothness.

The leading radial scalar matrix is
\[
F=-L^{-1}K=
\begin{pmatrix}11/2&-21/8\\-15/4&77/16\end{pmatrix}.
\]
After factoring $\Omega^2$, $F_0=F-2I$ has eigenvalues
$0,101/16$ and the exact symmetrizer
\[
S=\operatorname{diag}(10,7),\qquad
SF_0=\begin{pmatrix}35&-105/4\\-105/4&315/16\end{pmatrix}\geq0 .
\]
The natural modal amplitudes
\[
A_0=(45T_4+42Z_{3n})/101,\qquad
B_\mu=(56T_4-42Z_{3n})/101
\]
have leading radial equations
$\Omega\partial_\Omega A_0=0$,
$\Omega\partial_\Omega B_\mu=(101/16)B_\mu$.
At integer radial recursion orders $j\geq1$,
$I(2+j)$ is invertible. In particular there is no integer indicial
resonance through order six. A normal Frobenius recursion is therefore
consistent with log-free integer jets, conditional on the full
regular-singular reduction. Uniform time/angular existence is not
established by this observation.

Subtracting the analytic natural-field jet through order six leaves a
fractional exponent $101/16$, separated from weight six by $5/16$.
This suggests a trace-compatible remainder estimate rather than an
arbitrary positive weight admitting singular conformal traces.
The required Hardy estimates, evolving outgoing jets and angular
frequency uniformity remain open.

The independent exact-arithmetic script \texttt{indicial-audit.mjs}
checks both displayed determinants, the root kernel, the radial modal
decomposition, the value-residue polynomial and Routh minors, and the
strict glancing dissipation. Its saved JSON output explicitly excludes
a PDE well-posedness claim.

\subsection{Constraint-free normal gauge response}
In the frozen normalized model, nullness forces a normal harmonic
coordinate displacement $q_B$ through $(q_B)_t=F_B$.
For a lapse mode with frequency
$\omega=K(\epsilon-\mu)$, $\mu=K_x/K$,
\[
F_B=\omega p/(iK^2),\qquad F_B/(i\omega)=-p/K^2.
\]
For TT polarization amplitude $\mathcal H$,
\[
F_B=(1-\mu^2)\mathcal H/2
=\omega(\epsilon+\mu)\mathcal H/(2K).
\]
Thus the apparent grazing denominator cancels on constraint-free modes.
The resulting displacement is the trace of an order-minus-two potential
of curvature/lapse momentum, or an order-minus-one potential of the
metric. The ordinary $H^2$ bulk to $H^{3/2}$ trace estimate then controls
the constant-coefficient harmonic extension without the earlier scalar
constraint packet's loss. Low-frequency gauge normalization is needed.
No variable-coefficient sphere reconstruction theorem is inferred from
these Fourier identities.

\subsection{Why a published weighted theorem cannot simply be cited}
Chrusciel--Lengard, \emph{Radiation Fields},
\url{https://numdam.org/item/10.24033/bsmf.2478.pdf},
provide a useful weighted hyperboloidal framework, but their
Proposition 3.1 hypotheses (3.17)--(3.18) do not cover the current
nonzero glancing $1/\Omega$ residue. They require bounded characteristic
tangent lower-order coefficients and weaker crossblock singularities.
Their Theorem 3.4 also assumes already controlled space/time derivatives
and a bounded tangent block. Jet subtraction changes forcing, not that
homogeneous singular coefficient. A new dissipative Fuchsian estimate,
or another applicable theorem with all hypotheses verified, is needed.

\section{A common full subsidiary energy: improved candidate}
This section supersedes the parameters of the preceding finite-order
candidate. It retains the original tensor curvature equation and uses
the Appendix-C scalar equation ($\epsilon=1$), but sets
\[
\boxed{a\kappa_1=4,\qquad \kappa_2=-99/100}
\]
on an outer collar. These are a mathematically checked candidate
choice, not parameters adopted in production or validated by evolution.
Unlike a Hurwitz test on values alone, the following calculation uses
one energy for both the complete wave principal part and its singular
residue, including constraint-violating perturbations.

Write $l=1/a$, $d_0=1/(aS)$, $D_0=\partial_t-B\cdot\nabla$ and
\[
D_K=3\kappa_1(1+\kappa_2),\quad
d_\theta=\kappa_1(2+\kappa_2),\quad C_K=2l-\kappa_1.
\]
The base weighted fields and compatible derivative fields are
\[
h_c=H/\Omega^2,\quad m=M/\Omega,\quad
\tau=T/\Omega^2,\quad \zeta=z/\Omega,\quad
Q_i=\partial_iT/\Omega,\quad Z_{ij}=\partial_i z_j,
\]
\[
Q_i=2G_i\tau+\Omega\partial_i\tau,\qquad
Z_{ij}=G_i\zeta_j+\Omega\partial_i\zeta_j,\qquad
\pi=h_c/2+\Tr Z,\quad W=m+Q .
\]
The principal equations are four scalar waves and four advected values:
\[
D_0\pi\simeq A\partial_iQ_i,\quad
D_0Q_i\simeq A\partial_i\pi,\quad
D_0W_j\simeq A\partial_iZ_{ij},\quad
D_0Z_{ij}\simeq A\partial_iW_j.
\]
The values $\tau,\zeta$ have no principal derivatives in this reduction.
The equality of mixed derivatives of $z$ is used; if all derivative
fields are independently evolved, this amounts to a reduction-curl
constraint addition, and its compatibility must be preserved.

For completeness, one explicit first-order form, before the
$(\pi,W)$ change of variables, is the following. Put $g=G\cdot\zeta$,
\[
e_{ij}=Z_{ij}+Z_{ji}+2(G_i\zeta_j+G_j\zeta_i)
-\frac43\delta_{ij}g-\frac{D_K}{3}\delta_{ij}\tau,\qquad
e=2\Tr Z-D_K\tau .
\]
Then
\begin{align}
D_0h_c={}&-2A\operatorname{div}m+4G\cdot m
+\frac A\Omega[4G\cdot m-8l\Tr Z+4lD_K\tau
-(6l+2w_*)h_c],\\
D_0\tau={}&\frac A\Omega[h_c/2+\Tr Z-(d_\theta+2w_*)\tau],\\
D_0\zeta_i={}&-2l\zeta_i+
\frac A\Omega[m_i+Q_i+(C_K-w_*)\zeta_i],\\
D_0Q_i={}&A\partial_i(h_c/2+\Tr Z)
+\frac A\Omega[-(w_*+d_\theta)Q_i
+G_i(h_c/2+\Tr Z+d_\theta\tau)]\notag\\
&-lQ_i-G_i(h_c/2+\Tr Z-d_\theta\tau),\\
D_0Z_{ij}={}&A\partial_i(m_j+Q_j)
+\frac A\Omega C_K(Z_{ij}-G_i\zeta_j)
-3lZ_{ij}-G_i(m_j+Q_j+C_K\zeta_j),\\
D_0m_i={}&-\frac A2\partial_i h_c
+A(\partial_jZ_{ji}-\partial_i\Tr Z)\notag\\
&+\frac A\Omega[(P-w_*)m_i+2G_i\Tr Z-\tfrac{10}3G_jZ_{ij}
-2G_i g-6G^2\zeta_i+\tfrac{2D_K}3Q_i-\tfrac{2D_K}3G_i\tau]\notag\\
&+G_i h_c-lm_i-\tfrac{20}3Ad_0\zeta_i-G_je_{ij}+G_i e.
\end{align}
These formulas, together with
\[
w_*=-l+l\Omega/A,\quad G^2=l^2-2d_0\Omega,\quad
G=-ln+\frac{2\Omega}{S+r}n,
\]
exhibit a bounded remainder: in consistently dimensioned variables
$U=(\pi,W,Q,Z,\tau/a,\zeta/a)$,
\[
U_t=(B^iI+AP_i)\partial_iU+
\frac{A}{a\Omega}J_0(n)U+J_{\rm reg}(x)U,
\]
where $J_{\rm reg}$ is smooth and bounded on a collar. There are no
unaccounted $\Omega^{-2}$ coefficients in this reduction.

The principal symmetrizer is the wave Euclidean metric, together with
\[
\boxed{
\mathcal E(U)=|\pi|^2+|W|^2+|Q|^2+|Z|^2+
a^{-2}(4\tau^2+4\tau\zeta_n+8|\zeta|^2).
}
\]
It is uniformly positive. In a frame with $n$ as normal,
the dimensionless scalar residue on
$(\pi,W_n,Q_n,Z_{nn},Z_{11}+Z_{22},\tau/a,\zeta_n/a)$ is
\[
J_s=
\begin{pmatrix}
-4&-2&2&-2&-2&6/25&-2\\
-1&-2&-24/25&4/3&-2&-99/25&-8\\
-1&0&-76/25&0&0&-101/25&0\\
0&0&0&-2&0&0&-2\\
0&0&0&0&-2&0&0\\
1&0&0&0&0&-51/25&0\\
0&1&0&0&0&0&-1
\end{pmatrix}.
\]
The scalar energy matrix is
$\operatorname{diag}(1,1,1,1,1/2)\oplus
\left(\begin{smallmatrix}4&2\\2&8\end{smallmatrix}\right)$.
The exact $LDL^T$ pivots of $-\operatorname{sym}(H_sJ_s)$ are
\[
4,\quad\frac{23}{16},\quad\frac{38346}{14375},\quad
\frac{7036}{8217},\quad\frac{59165}{147756},\quad
\frac{37690637}{73956250},\quad\frac{211573056}{37690637}.
\]
Each of the two transverse five-dimensional blocks has energy
$\operatorname{diag}(1,1,1,1,8)$ and dissipation pivots
\[
2,\quad\frac{1828}{625},\quad2,\quad
\frac{4577}{8226},\quad\frac{51555}{9154}.
\]
The remaining three tangential $Z$ components have dissipation $2I$.
This accounts for all twenty components. Every pivot is positive.
The exact rational program \texttt{common-energy-audit.mjs} checks
these identities and independently checks the compatible radial
pencil below.

Consequently, every sufficiently regular, reduction-compatible solution
obeys on a collar
\[
\boxed{
\frac{dE}{dt}\leq C E
-c\int\frac{A}{a\Omega}|U|^2\,dx
+\hbox{controlled inner-interface flux}.
}
\]
At scri the principal energy flux has wave eigenvalues $0,-2A$,
and advected eigenvalues $-A$; it is nonpositive. No incoming
constraint boundary data are prescribed there. The value metric can
be smoothly patched to a nonsingular interior metric; the wave metric
is unchanged, and all interior lower-order coefficients are bounded.
Stationarity and rotational covariance give the same a priori
estimate for any fixed finite collection of time derivatives and
rotational Lie derivatives, using their tensor actions.

This is an actual a priori subsidiary energy estimate, stronger than
the earlier residue-eigenvalue tests. It does \emph{not} by itself
prove existence of a closed-domain evolution, stronger normal
falloffs, or any coupled geometric/lapse/elliptic estimate.

\subsection{Compatible normal powers and a finite-graph-norm route}
For the improved parameters the reduced scalar pencil is
\[
L=\begin{pmatrix}302/25&8\\49/25&8/3\end{pmatrix},\qquad
K=\begin{pmatrix}-596/25&-4\\-102/25&-34/3\end{pmatrix},
\]
\[
\det(sL+K)=\frac8{15}(s-2)(31s-238),\qquad
F=-L^{-1}K=
\begin{pmatrix}58/31&-150/31\\24/155&242/31\end{pmatrix}.
\]
Thus the two compatible scalar powers are $2$ and $238/31$.
The root-two natural fields satisfy
\[
H=\Omega^4H_4,\quad T=\Omega^4T_4,\quad
M=\Omega^3M_3,\quad z=\Omega^3Z_3,
\]
\[
\boxed{
H_4=-\frac{2}{25a}T_4,\quad
M_{3n}=\frac{302}{75a}T_4,\quad
Z_{3n}=-\frac2{75}T_4,\quad
M_{3A}=-Z_{3A}/a.
}
\]
The other scalar natural exponent is $\mu=176/31$.
The modal rows $(6,5)$ and $(2,75)$ diagonalize $F-2I$,
and
\[
S=\begin{pmatrix}40&180\\180&5650\end{pmatrix}>0,\qquad
S(F-2I)=\frac{176}{31}
\begin{pmatrix}4&150\\150&5625\end{pmatrix}\geq0 .
\]
There is no integer indicial resonance.

A finite normal regularity theorem is given in the included proof sections.
Let $\mathcal G_{40}$ include actual time jets and all words in the
conormal dilation, covariant rotations and cutoff Cartesian interior
derivatives through total order forty. Conormal regularity is present
initially and propagated directly. The earlier time/rotation-only
inference is withdrawn. Time jets must be the actual PDE-generator
jets, not independently prescribed functions.

The exact radial reduction is irregular singular in a fast block:
the outgoing glancing coefficient vanishes quadratically. With
$v=(\tau,\zeta)$, $N=\Omega\partial_\Omega v$ and
$R_g=N-Fv$, its schematic but order-explicit form is
\[
\Omega\partial_\Omega v=Fv+R_g,\qquad
\Omega^3\gamma^2\partial_\Omega R_g-L_g R_g
+q_g\Omega\partial_tR_g
=\Omega\mathcal D_1 v+\Omega^2\mathcal D_2v+
\Omega^2 C R_g ,
\]
where $\gamma=r/(aS)$, $-L_g$ is coercive, $q_g>0$,
$\mathcal D_j$ has time/angular order at most $j$, and there
are no angular derivatives of $N$ in the second-order normal equation.
Ordinary Frobenius theory does not apply to this full fast block.

At fixed time, a coercive radial Volterra inverse gains one
$\Omega$ power at a cost of at most two additional time/angular
generators. The slow Fuchsian equation then gains the same power
modulo crossed indicial jets. Eight gains from base $b$-weight
$-1/2$ to $15/2$, six reserved conormal levels, six angular derivatives
and two time derivatives have conservative generator count
$16+6+6+2=30<40$. The simultaneous-weight embedding is
\[
v=\sum_{m=2}^7\Omega^m a_m(t,\vartheta)+R,\qquad
\Omega^{-15/2}R\in H_b^6.
\]
The root-two coefficient is in the three-dimensional outgoing kernel;
higher integer jets are fixed recursively. Both algebraic branches
remain allowed. After dividing $v$ by $\Omega^2$, the remainder
weight is $11/2$, sufficient for ordinary normal $H^6$; its distance
from $\mu$ is $11/62$.

This paragraph records the route; the subsequently included sections
give the weak construction, exact wave coefficient audit and finite
graph-domain embedding argument for this linear subsidiary system.
A spacetime fast-energy argument has an additional initial half-weight
trace term; the fixed-time Volterra route avoids assuming that term.
The capped construction and exact homogeneous reduction-error transport
give a compatible limiting realization. Finite graph regularity remains
a strong domain requirement rather than an arbitrary ordinary-$H^6$
initial-data theorem. These subsidiary results do not supply the missing
geometric/gauge boundary estimate.

\section{A physical elliptic normal-trace obstruction}
There is a separate obstruction to a simple all-angular extension of
a scalar lapse correction. Constraint falloffs, even exact constraints,
do not determine the elliptic shift's normal derivative.

The scalar York pieces, with $\phi=\Omega c$, are
\[
a_c=\frac a2(\operatorname{Hess}\phi)^\TF,\quad
k_c=\frac a2[\Omega\Delta\phi-3G\cdot\nabla\phi-3d_0\phi],
\quad
v_c=\frac a2[A\nabla\phi+\phi G].
\]
They satisfy exactly
\[
\Omega\operatorname{div}a_c-2a_cG=\frac23\nabla k_c,\qquad
Lv_c=2A\operatorname{div}a_c-2a_cG .
\]
Under metric Dirac gauge the Hamiltonian identity is
$k=k_c+(3a/2)hGG-aH/4$.
A potential determined only by $h$ can reproduce this additional
Hamiltonian term, but it cannot remove physical TT curvature from
the shift normal derivative.

Indeed choose initially
\[
h=c=u=k=T=z=0,\qquad a_{ij}=\Omega^2t_{ij},
\]
where $t$ is a smooth compactly supported Euclidean TT tensor.
Then $H=M=0$, while
\[
Lv=2A\operatorname{div}a-2aG,\qquad v|_{\scri}=0.
\]
Generically $\partial_nv_n|_{\scri}\ne0$, so $v_n=O(\Omega)$
rather than $O(\Omega^2)$. One exact demonstration takes $a=S=1$,
an $\ell=2$ scalar harmonic $Y$, and a nonnegative nonzero radial
$P$ compactly supported in $(0,1)$:
\[
t_{rr}=PY,\quad t_{rA}=QD_AY,\quad
t_{AB}=r^2[-Pq_{AB}Y/2+RY_{AB}],
\]
\[
Q=\frac{r^2}{6}(P'+3P/r),\qquad R=(2Q'+4Q/r-P)/4 .
\]
Use the origin-regular solution of $L\xi=0$ with $\xi|_{\scri}=nY$,
\[
\xi_r=(11r+6r^3)Y/17,\qquad
\xi_A=(11/34)r^2(1-r^2)D_AY .
\]
Green's identity reduces the relevant volume pairing, with
$\psi=A\Omega^2$, to a positive multiple of
\[
\int_0^1 r^2P(r)
\underbrace{[\psi''\xi_A/Y_A+
\psi'((\xi_A/Y_A)'-\xi_r/Y-(\xi_A/Y_A)/r)]}
_{(3/34)r^4(1-r^2)(1+25r^2)>0}\,dr .
\]
Thus $\partial_nv_n$ is nonzero. This disproves a proposed lapse
source $G\cdot(v_c-v+v_h)/\Omega^2$ with an $h$-only potential $v_h$:
all its other scalar sources vanish on these physical data.

A repair must include the TT volume response, its
Dirichlet-to-Neumann functional, or a boundary gauge formulation that
controls that response. This is a counterexample to that particular
closure, not to every elliptic-shift hybrid.

\section{Conformal-harmonic candidate: pole cancellation is insufficient}
A reference reduction of a physical Bona--Masso condition does not
automatically yield a reference reduction of a conformal harmonic
condition. This distinction supplies a simpler candidate than the
failed scalar-potential extension above.

Let $\bar K=(K_{\rm phys}-3w)/\Omega$ be the trace of the
Penrose spatial extrinsic curvature, and use
$\bar{\hat K}=(K_{\rm phys}-2\Theta_{\rm phys}-3w)/\Omega$
in a conformal harmonic lapse. Its background value is
$\bar{\hat K}_*=-3/(aA)$. Adding a fixed reference source to
$(\partial_t-\beta\cdot\nabla)\alpha=-\alpha^2\bar{\hat K}$
makes the background stationary. Its exact first variation is
\[
\boxed{
u_t=B\cdot\nabla u-G\cdot v+\frac6a u
-\frac{A^2}{\Omega}(q-3\delta w),\qquad
\delta w=-\frac{G\cdot v+w_*u}{A}.
}
\]
The fixed reference source is
$-Aw_*-3A/a$; it has zero variation.
The pole numerator is therefore
$-A^2q-3AG\cdot v-3Aw_*u$.
Geometric chi regularity requires $k=3\delta w$ at scri, so
this numerator reduces to $2A^2T$ there. Natural
$T=O(\Omega^4)$ removes it, without a new algebraic condition
relating $u_{\scri}$ to the metric trace.
This changes gauge sources in addition to choosing the harmonic
principal coefficient. It does not impose CMC.

An exact identity is especially useful. For
$Y=u+(3A/2)c$, the coupled lapse and chi equations give
\[
\boxed{
Y_t=B\cdot\nabla Y-A\operatorname{div}v-G\cdot v
+\frac3a u-\frac{3Aw_*}{2}c+\frac{2A^2}{\Omega}T .
}
\]
All other pole terms cancel. No Hamiltonian, momentum or connection
constraint was set to zero in obtaining this identity.
It is a regular equation for this combination under the scalar
constraint falloff. It does not alone produce a closed regular
evolution for $\bar{\hat K}$, since differentiating
$\delta w$ introduces the time derivative of the elliptically
reconstructed shift.

At each fixed $\Omega>0$ the lapse has the usual $\hat K$ harmonic
principal part, rather than the defective physical-$K$ driver.
An interior $1+\log$-like modification can be blended with a coefficient
$f-1=O(\Omega^2)$ in the collar. The absence of a new lapse residue is
a useful necessary repair. It is not a proof of an all-angular
elliptic/hyperbolic boundary estimate.

A second, explicitly nonlocal, gauge option fixes $u|_{\scri}=0$.
For example let
\[
D=cG^2-hGG,\qquad
N_s=-A^2q-AG\cdot v+Aw_*u-aA^2D,
\qquad F_s=B\cdot\nabla u+N_s/\Omega .
\]
On compatible data with $H=T=0$ at scri, nullness gives
$N_s|_{\scri}=-2Au/a=0$. Its finite limiting RHS includes the
normal derivative of the elliptic shift, including the TT volume
response. A boundary projection can remove exactly that trace:
\[
u_t=P_AF_s,\qquad
P_A=I-E_A\operatorname{Tr},\qquad
\operatorname{div}(A^{-2}\nabla E_Ag)=0,\quad E_Ag|_{\scri}=g.
\]
This projects a gauge RHS, not the Einstein constraints. The
existence and regularity of $\operatorname{Tr}F_s$ must be proved
in the proposed full state space.

There is a genuine model estimate supporting this option. For
\[
u_t=P_A(B\cdot\nabla u-A^2p),\qquad
p_t=B\cdot\nabla p-\Delta u,\qquad u|_{\scri}=0,
\]
the energy
$\int(A^{-2}|\nabla u|^2+p^2)\,dx$ satisfies an estimate
without derivative loss. Indeed
$\int A^{-2}\nabla u\cdot\nabla E_Ag=0$ by the zero trace
of $u$, while the null boundary flux is
\[
-A(p+u_n/A)^2\leq0.
\]
Smooth coefficient gradients contribute bounded bulk terms.
This estimate includes the elliptic gauge projection; it is more than
an interior-symbol calculation. The full $F_s$ has configuration and
shift terms whose gradient-energy pairing is not yet controlled.
Their treatment requires a complete York-coupled estimate, rather
than declaring their bounded $L^2$ sizes sufficient.

The exact counterexamples below invalidate the proposed conformal-harmonic
boundary-model stability claim despite its pole cancellation. A strictly
subluminal coefficient alone also has an explicit singular-outflow problem.
The projected fixed-trace choice remains a different candidate with only
an isolated gauge-model energy argument. No full three-dimensional theorem
is inferred from any of these partial calculations.

\section{A weak realization of the linear subsidiary system}
\label{sec:constraint-existence}

This section concerns only the linear, first-order subsidiary system. It does
not prove existence for the geometric evolution or for its elliptic gauge.
The parameters are $a\kappa _1=4$, $\kappa _2=-99/100$, and the
Appendix-C scalar correction $\epsilon=1$. We use nondimensional $a=1$
when displaying its positive matrices. For general $a$, the value variables
in the energy are $\widehat\tau=\tau/a$ and
$\widehat\zeta=\zeta/a$; the wave variables already have the same dimensions.

Let $\mathcal B=\{r<S\}$ and let $\Omega=(S^2-r^2)/(2aS)$.
After the exact differential reduction, write the twenty-component equation as
\begin{equation}
 \partial_t U=\mathcal A U
 =V^i(x)\partial_iU+
 \frac{A(x)}{a\Omega(x)}J_0(n)U+R(x)U
 \qquad\hbox{in a collar of }\partial\mathcal B .
 \label{eq:subsidiary-abstract}
\end{equation}
Here $R$ is smooth and bounded up to the boundary. On the sixteen wave
components $V^i=B^iI+AP_i$, where the $P_i$ are the constant first-order
four-scalar-wave matrices; on the four value components $V^i=B^iI$.
The exact reduction, not merely an indicial calculation, supplies this
decomposition. The matrix $J_0$ is its boundary residue, with all bounded
differences absorbed into $R$.

The common metric is the identity on wave velocities and gradients, and on
the values has quadratic form
\begin{equation}
 4\widehat\tau^2+4\widehat\tau\widehat\zeta_n
       +8|\widehat\zeta|^2 .
 \label{eq:subsidiary-value-metric}
\end{equation}
Denote the resulting matrix by $H(n)$. The exact rational block calculation
establishes constants $h_*,h^*,\sigma>0$, independent of $n$, such that
\begin{equation}
 h_*I\le H\le h^*I,\qquad HV^i=(HV^i)^T,
 \qquad -\operatorname{sym}(HJ_0)\ge\sigma I .
 \label{eq:subsidiary-common-metric}
\end{equation}
The uniformity follows also from rotational covariance. No asserted
positivity of an uncomputed interior residue is needed.

Choose a smooth radial cutoff $\chi$, equal to one in a smaller boundary
collar and zero near the center. Extend $H$ smoothly and positively into the
interior, retaining the identity on the wave sector and a positive metric on
the value sector. This keeps $HV^i$ symmetric: value principal matrices are
scalar multiples of the identity. Rewrite the global operator as
\begin{equation}
 \mathcal A=V^i\partial_i+
       \chi\frac{A}{a\Omega}J_0+R_{\rm b},
 \label{eq:subsidiary-global-decomposition}
\end{equation}
where $R_{\rm b}$ is bounded. The cutoff transition is a positive distance
from scri, so this change introduces no additional singularity. One may
interpolate the value metric rotationally covariantly; the resulting metric
is smooth at the center and the global operator remains rotationally
covariant.

\paragraph{Proposition (linear weak subsidiary evolution).}
Assume the exact coefficients and positive matrix identities above. For every
$U_0\in L^2_H(\mathcal B)$ there is a unique
$U\in C([0,\infty);L^2_H)$ solving
\eqref{eq:subsidiary-abstract} distributionally on $\Omega>0$ with initial
value $U_0$. It obeys, on every bounded time interval,
\begin{equation}
 \|U(t)\|_H^2+
 c\int_0^t\!\int_{\mathcal B}
      \frac{\chi A}{a\Omega}|U|^2\,dx\,ds
 \le C_T\|U_0\|_H^2 .
 \label{eq:subsidiary-weak-energy}
\end{equation}
The solution operators form a strongly continuous semigroup. Its generator
is the maximal distributional realization
\begin{equation}
 D(\mathcal A)=\{U\in L^2_H:
          \mathcal A_{\rm formal}U\in L^2_H
          \text{ distributionally on }\Omega>0\}.
 \label{eq:subsidiary-maximal-domain}
\end{equation}
In particular, no independently chosen incoming boundary trace is added at
scri. This statement is about the first-order subsidiary extension; the
reduction-compatible subspace is invariant as proved below.

\paragraph{Proof: bounded approximants.}
Replace only the explicitly certified singular multiplier in
\eqref{eq:subsidiary-global-decomposition}, setting
\[
 q_N(\Omega)=\frac{1}{\sqrt{\Omega^2+N^{-2}}},\qquad
 \mathcal A_N=V^i\partial_i+
          \chi\frac Aa q_NJ_0+R_{\rm b}.
\]
For $\Omega>0$, $0\le q_N\le\Omega^{-1}$ and $q_N\uparrow\Omega^{-1}$.
Every approximant preserves precisely the same pointwise dissipativity
certificate. We do not cap an entire unverified-sign residue.

For each finite $N$, extend the smooth bounded coefficients and a positive
symmetrizer to a neighborhood of the closed ball, then to a bounded-coefficient
whole-space symmetric hyperbolic system. Ordinary Friedrichs regularization
and the bounded-coefficient energy estimate construct its $L^2$ Cauchy
solution. Restriction to the ball is independent of the exterior extension.
Indeed, at $r=S$ the wave normal coefficient $B_nI+AP_n$ has eigenvalues
$0,-2A$, while the value coefficient is $-A$. Thus its energy boundary flux
is nonpositive. A difference with zero data inside the ball has zero energy
there. The same argument first for smooth solutions and then by $L^2$
approximation justifies the restriction without prescribing extra boundary
data. Coefficients may be cut off farther outside the ball to ensure all
whole-space coefficients are bounded; this does not affect the restriction.

Integration by parts and \eqref{eq:subsidiary-common-metric} give
\begin{equation}
 \|U_N(t)\|_H^2+
 2\sigma\int_0^t\!\int\chi\frac Aa q_N|U_N|^2
 \le\|U_0\|_H^2+C\int_0^t\|U_N(s)\|_H^2\,ds .
 \label{eq:subsidiary-cap-energy}
\end{equation}
Here $C$ depends on bounded $R_{\rm b}$ and derivatives of $HV^i$, and is
independent of $N$. The potentially large derivatives of $q_N$ never enter
this zeroth-order energy estimate.

\paragraph{Proof: passage to the weak limit.}
On every finite time interval the approximants are bounded in
$L^\infty_tL^2_x$. On every compact subset of $\Omega>0$, their time
derivatives are bounded in $L^\infty_tH^{-1}_x$, uniformly in $N$.
A diagonal weak extraction therefore gives a limit, weakly continuous in
$L^2$ and solving the original equation distributionally in the open ball.
One may use a countable dense family of compactly supported smooth tests and
their uniformly equicontinuous pairings to obtain convergence at every time.
Density and the uniform $L^2$ bound give global weak continuity and the
correct weak initial trace. No strong convergence of a product of unknowns
is required: the equation is linear, and every coefficient is bounded on
each compact subset of the open ball.

For each fixed interior compact set, weak lower semicontinuity passes the
weighted part of \eqref{eq:subsidiary-cap-energy} to the limit. Exhaustion
by compact sets and monotone convergence then give the weighted integral in
\eqref{eq:subsidiary-weak-energy}. The same lower semicontinuity gives the
endpoint energy bound. This proves weak existence with the asserted energy.

\paragraph{Proof: weak uniqueness without assuming boundary traces.}
Let $\eta_\varepsilon(\Omega)$ vanish for $\Omega<\varepsilon$ and equal one
for $\Omega>2\varepsilon$, with
$0\le\eta_\varepsilon'\le C/\varepsilon$. Local Friedrichs mollification
justifies the energy identity with multiplier $\eta_\varepsilon H$ for an
$L^2$ weak solution: on its support all reaction coefficients are bounded,
the principal coefficients are Lipschitz, and their Friedrichs commutators
converge in $L^2$. This is a local statement; it does not differentiate the
singular multiplier across scri.

The only potentially positive cutoff error comes from the glancing normal
subspace. Indeed,
\[
 V^i\partial_i\Omega=-\gamma(B_nI+AP_n),\qquad
 \gamma=\frac r{aS},\qquad
 B_n+A=\frac{(S-r)^2}{2aS}=O(\Omega^2).
\]
Every other wave/value eigenvalue of $V^i\partial_i\Omega$ is nonnegative
in a sufficiently small collar. Hence the positive part of the cutoff error
is at most
\[
 C\varepsilon\int_{\{\varepsilon<\Omega<2\varepsilon\}}|U|^2\,dx.
\]
Its time integral tends to zero for $L^\infty_tL^2_x$ solutions. All negative
cutoff fluxes may be discarded. The certified reaction is dissipative and
its nonnegative integral passes by Fatou. Letting $\varepsilon\downarrow0$
therefore gives the same global energy inequality for a weak difference.
Zero initial data imply zero solution. No pointwise trace, selected sequence
of boundary spheres, or prior weighted trace estimate is required.

The weak initial continuity and upper energy bound imply strong continuity
at zero. Uniqueness and time independence give the semigroup property;
boundedness on finite time intervals then gives strong continuity at every
time. Thus the weak construction produces a $C_0$ semigroup, rather than
assuming one in advance.

\paragraph{Proof: identification of the maximal domain.}
For $U\in D_{\max}$ as defined in \eqref{eq:subsidiary-maximal-domain}, the
stationary local Friedrichs/cutoff argument gives
\begin{equation}
 \operatorname{Re}\langle\mathcal A_{\rm formal}U,U\rangle_H
 \le C\|U\|_H^2-
 \sigma\int\chi\frac A{a\Omega}|U|^2 .
 \label{eq:subsidiary-stationary-domain-energy}
\end{equation}
In particular the weighted integral is finite; it need not be imposed as a
separate domain assumption. The cutoff term has the same favorable sign
and $O(\varepsilon)$ positive bound as above.

The semigroup generator is contained in $D_{\max}$ by differentiating the
distributional equation against compactly supported tests. Conversely,
take $U\in D_{\max}$, choose real $\lambda>C$, and set
$F=(\lambda-\mathcal A_{\rm formal})U$. The Bochner integral
\[
 V=\int_0^\infty e^{-\lambda t}S(t)F\,dt
\]
belongs to the semigroup generator domain and satisfies
$(\lambda-\mathcal A_{\rm formal})V=F$. Therefore $W=U-V\in D_{\max}$
satisfies $(\lambda-\mathcal A_{\rm formal})W=0$.
Equation \eqref{eq:subsidiary-stationary-domain-energy} forces $W=0$.
This proves equality of the generator and maximal distributional domains.
It also makes the higher graph spaces precise: generator powers are the
corresponding recursively defined formal powers in $L^2$.

\paragraph{Reduction compatibility.}
Return to the unscaled value variables and define, distributionally in the
open ball,
\[
 E^Q_i=Q_i-2G_i\tau-\Omega\partial_i\tau
       =Q_i-\frac{\partial_iT}{\Omega},\qquad
 E^Z_{ij}=Z_{ij}-G_i\zeta_j-\Omega\partial_i\zeta_j
       =Z_{ij}-\partial_i z_j .
\]
The exact first-order equations, with no use of a Hamiltonian or momentum
projection, give
\begin{align}
 (\partial_t-B\cdot\partial)E^Q_i
  &=-\left[l+\frac{A(w+d_\theta)}\Omega\right]E^Q_i,\\
 (\partial_t-B\cdot\partial)E^Z_{ij}
  &=\left[\frac{Ac}\Omega-3l\right]E^Z_{ij},
 \qquad d_\theta=\kappa_1(2+\kappa_2),\quad c=2l-\kappa_1.
 \label{eq:subsidiary-reduction-transport}
\end{align}
All forcing cancels exactly. Backward characteristics of these equations
remain inside the ball: forward characteristics solve $\dot x=-B=x/a$.
On every compact portion of a backward characteristic the coefficients are
smooth and bounded. Local distributional uniqueness for transport thus
preserves zero initial reduction errors, without any data at scri.
The compatible subspace is closed in $L^2$ with its distributional
interpretation and invariant under $S(t)$. Consequently
$Z_{ij}=\partial_i z_j$ and
$\partial_jZ_{ij}=\partial_iZ_{jj}$ hold for compatible limits even though
the capped approximants need not preserve these identities.

\paragraph{The explicit joint phase and its interior generators.}
There is no need to assume that $\mathcal A_N^jU_0$ is bounded uniformly in
$N$; capping can break cancellations.  Maximal-domain identification instead
makes time-generator graph domains of the limiting semigroup invariant.
For the normal embedding we use a \emph{stronger} joint phase from the outset,
not just this time graph domain.  Set
\[
 \mathcal D=b(r)\partial_r,\qquad b=\Omega/\Omega_r
\]
on an outer collar, and cut it off smoothly away from the centre.  Let
$\chi\partial_i$ be three Cartesian generators supported away from scri,
with $\chi=1$ near the centre.  Choose overlapping cutoffs so that these
and $\mathcal D$, together with the three covariant rotations, span the
ordinary spatial derivatives uniformly on every compact interior region.
Write $\mathcal Z$ for this finite family.  At order $N$ the phase norm is
\[
 \mathcal G_N(U)=\sum_{p+|I|\le N}
  \|\mathcal Z^I\partial_t^pU\|_{L^2_H},
\]
where all words $I$ in the generators are included and time jets are the
actual limiting-generator jets.  Thus the phase explicitly includes
conormal and ordinary interior regularity.  For example $N=128$ controls
ordinary interior derivatives through order $128$.  It does not assert that
time/rotation graph regularity alone generates the conormal derivatives.

\paragraph{Propagation of the joint phase without a capped commutator.}
Write $\mathcal A=V^r\partial_r+V^A\partial_A+R$ in the collar, with
$R=A_*J_0(n)/(a\Omega)+R_b$.  The exact identity is
\[
 [\mathcal A,\mathcal D]=b'\mathcal A
 -(b'V^A+b\partial_rV^A)\partial_A
 -(\partial_rV^r)\mathcal D-b'R-b\partial_rR.
\]
Its apparent pole is removable, since $J_0(n)$ is radially constant and
\[
 b'\frac{A_*J_0}{a\Omega}
 +b\partial_r\frac{A_*J_0}{a\Omega}
 =\left[\frac{(b'-1)A_*}{a\Omega}
       +\frac{\partial_rA_*}{a\Omega_r}\right]J_0,
 \qquad b'-1=-\frac{\Omega\Omega_{rr}}{\Omega_r^2}.
\]
All remaining coefficients are smooth and bounded.  Replacing
$\mathcal AU$ by $U_t$ makes the conormal hierarchy triangular: a current
conormal derivative has the original generator plus bounded multiplication,
and its time/angular forcing has lower conormal order at the same total
joint order.  Rotational covariance gives the exact rotational commutations.
The interior Cartesian commutators have only bounded first-order
coefficients, supported where $\Omega$ is bounded below; expressing them in
the same overlapping generator family closes the ordinary interior jet
energy at the same total order.

Construct each conormal level using the limiting semigroup and bounded
perturbations, with its prescribed joint initial value.  On interior domains
the distributional derivative satisfies the same equation and initial data;
local uniqueness identifies it with the constructed $L^2$ field.  Backward
domains remain inside the ball at finite time, so no unknown derivative
trace at scri is required.  The finite augmented energy and completion in
$\mathcal G_N$ give an invariant strongly continuous joint phase.  This uses
no uniform commutator bound for the capped multipliers.

The previous unqualified claim that a time/rotation-only $\mathcal G_{40}$
supplied the six conormal derivatives is withdrawn.  The normal embedding
below uses the explicitly strengthened joint phase just defined.  It still
requires the fast/slow weight estimates and finite indicial-jet recursion;
no coupled geometric/gauge conclusion follows from the subsidiary
semigroup alone.


\section{Exact four-wave reduction and audit of normal powers}
\label{sec:constraint-wave-reduction}

This is a coefficient-level audit of the normal structure used in
Section~\ref{sec:constraint-normal-bootstrap}. Put
\[
 l=1/a,\quad d_0=(aS)^{-1},\quad
 d=\kappa_1(2+\kappa_2),\quad D=3\kappa_1(1+\kappa_2),
 \quad c=2l-\kappa_1,\quad C=A+\Omega,
 \quad D_0=\partial_t-B\cdot\partial,
\]
and write $\mathcal P=D_0^2-A^2\Delta$,
$s_z=\partial_i z_i$, $g_z=G\cdot z$. All vector components in the following
two equations are Cartesian. The exact corrected subsidiary equations give
\begin{align}
H&=(2\Omega/A)D_0T-2\Omega^2s_z+2dT,\\
M_i&=(\Omega/A)(D_0z_i+2lz_i)-\partial_iT-cz_i.
\end{align}
Substitution into the exact Hamiltonian and momentum equations, using
$[D_0,\partial_i]=-l\partial_i$, gives the scalar equation
\begin{align}
0={}&\mathcal P T+
 \frac{wC+A(d+6l)}\Omega D_0T
 -(2A+\Omega)G\cdot D_0z\\
&+[-2A^2w+3lA\Omega-A^2(c+2l)]s_z
 +\frac{A(3A+2\Omega)}\Omega G\cdot\partial T\\
&+\left[-2l(2A+\Omega)+\frac{Ac(3A+2\Omega)}\Omega\right]g_z
 -\frac{2lA^2(D-3d)}{\Omega^2}T .
\label{eq:constraint-scalar-wave-exact}
\end{align}
For the vector equation define the scalar coefficient
\[
 Z_c=\frac{2lwC}{A}+2l^2(\Omega/A+3)
 -lc(1+3A/\Omega)+2G^2+\frac{20}{3}Ad_0+
 \frac{6AG^2}\Omega.
\]
Then the exact vector equations are
\begin{align}
0={}&\mathcal P z_i+
 \left[\frac{wC}\Omega+3l+\frac{A(3l-c)}\Omega\right]D_0z_i
 -\frac{CG_i}{\Omega^2}D_0T-\frac{2A^2G_i}\Omega s_z\\
&+A G\cdot\partial z_i+
 \frac{A(10A/3+\Omega)}\Omega G_j\partial_i z_j
 +\frac A\Omega Z_c z_i\\
&+\left[\frac{2A}{3\Omega}+\frac{2A^2}{\Omega^2}\right]G_i g_z
 -\frac{A^2(3l-d+2D/3)}{\Omega^2}\partial_iT
 -\frac{2\kappa_1AC}{\Omega^3}G_iT .
\label{eq:constraint-vector-wave-exact}
\end{align}
The $\Omega^{-3}T$ term is real; it must not be omitted. It is balanced by
the different powers in $T=\Omega^2\tau$ and $z=\Omega\zeta$.

One derivation of the vector identity first uses
\begin{align}
 D_p^jQ^p_{ij}-D_i\operatorname{tr}_pQ^p
 ={}&A\Omega(\Delta z_i-\partial_i s_z)
 -\Omega G\cdot\partial z_i
 -(10A/3+\Omega)G_j\partial_i z_j+2CG_is_z\\
&-\left[2G^2+\tfrac{20}{3}Ad_0+
                  \tfrac{6AG^2}\Omega\right]z_i
 -\left[\tfrac23+\tfrac{2A}\Omega\right]G_i g_z\\
&+\frac{2DA}{3\Omega}\partial_iT
 -\frac{2DC}{3\Omega^2}G_iT .
\end{align}
The mixed terms $D_0\partial_iT$ and $\partial_iD_0T$ cancel, as do
$\partial_i s_z$ from the Hamiltonian derivative and the tensor divergence.
Likewise $D_0s_z$ cancels $\partial_iD_0z_i$ in the scalar equation. Thus no
second-order cross principal operator survives.
The commutator in $D_0M_i$ also contributes $+l\partial_iT$ on the left;
it cancels the corresponding momentum-reaction contribution. There is
therefore no extra $-Al\partial_iT/\Omega$ in
\eqref{eq:constraint-vector-wave-exact}. A lightweight independent
automatic-differentiation evaluation of the nested and expanded identities
at a generic nonradial point detected and corrected precisely this term.

\paragraph{Audit of the normal eight-system.}
In spherical coordinates, with Cartesian dependent components,
\[
 \mathcal P=\partial_t^2-2B_r\partial_t\partial_r
 -\Omega^2\partial_r^2+
 (B_r\partial_rB_r-2A^2/r)\partial_r
 -A^2r^{-2}\Delta_{\mathbb S^2}.
\]
Since $\partial_r=-\gamma\partial_\Omega$,
$\gamma=r/(aS)$, substituting the balanced powers and multiplying the
balanced equations by $\Omega^2$ has exactly the following consequences:
\begin{itemize}
\item The highest radial term is $-\gamma^2\Omega^3\partial_\Omega N$,
where $N=\Omega\partial_\Omega(\tau,\zeta)$. The term obtained by
differentiating $1/\Omega$ is a smooth $O(\Omega^2)N$ coefficient.
\item The only derivative of $N$ in time/angular directions is
$2B_r\gamma\Omega\partial_tN$. There is no angular derivative of $N$.
\item All other first time/angular derivatives of $(\tau,\zeta)$ carry
at least one power of $\Omega$. Second time/angular derivatives carry
$\Omega^2$. Every zero-order and $N$ coefficient is smooth up to scri.
\item Passing to normal/tangential vector components changes only smooth
angular lower-order coefficients: the frame is constant along each radial
ray. It does not introduce a mixed normal/angular principal term.
\end{itemize}
These statements follow term by term from
\eqref{eq:constraint-scalar-wave-exact}--\eqref{eq:constraint-vector-wave-exact};
there is no residual $\Omega^{-1}$ in a derivative coefficient of the
eight-system after the displayed multiplication.

At scri, in $a=1$ units and for $\kappa_1=4$, $\kappa_2=-99/100$, the
scalar/normal $N$ and value matrices in
$-\gamma^2\Omega^3N_\Omega+LN+Kv+\cdots=0$ are
\[
 L_0=A_{\scri}^2
 \begin{pmatrix}-151/25&-4\\-49/25&-8/3\end{pmatrix},\qquad
 K_0=A_{\scri}^2
 \begin{pmatrix}298/25&2\\102/25&34/3\end{pmatrix}.
\]
Each transverse block is $L_0=-4A_{\scri}^2$,
$K_0=8A_{\scri}^2$. Hence $-L_0$ has positive symmetric part and
\[
 -L_0^{-1}K_0=
 \begin{pmatrix}58/31&-150/31\\24/155&242/31\end{pmatrix},
\]
with transverse entries $2$. Equivalently $L_0=2A_{\scri}^2lJ_{gg}$ with
the appropriately dimensioned glancing matrix. The fast coercivity therefore
uses the same wave metric as the previously certified energy, without an
unverified change of sign or target metric.

Finally, setting $F=-L^{-1}K$ and $R_g=N-Fv$ introduces only
$\Omega^3F_\Omega v$ and $\Omega^2F(R_g+Fv)$ from the highest normal term.
The first is higher order; the $R_g$ term can be absorbed into a smooth
coercive fast zero-order matrix in a smaller collar. The mixed time term
gives $\Omega\partial_tR_g$ and $\Omega F\partial_tv$. Thus the fast forcing
order/power count required by the finite graph bootstrap is verified by the
exact equations, not merely assumed from the indicial roots.


\section{Finite graph regularity at scri: a normal-bootstrap lemma}
\label{sec:constraint-normal-bootstrap}

This section records the precise normal-bootstrap argument separately from
weak existence. It is a linear regularity statement, not a sharp
$H^6\longrightarrow H^6$ geometric evolution theorem. The hypotheses below
include the exact normal reduction; a formal indicial polynomial alone would
not suffice.

\paragraph{Normal variables and exact structure.}
Work in a collar where $r$ is bounded away from zero. Put
\[
 v=(\tau,\zeta),\qquad N=\Omega\partial_\Omega v,
 \qquad\gamma=-\partial_r\Omega=\frac r{aS}>0.
\]
The derivative-reduction identities recover $N$ without a new differentiability
assumption:
\[
 N_\tau=Q_n/G_n-2\tau,\qquad
 N_{\zeta_j}=Z_{nj}/G_n-\zeta_j .
\]
All angular differentiations below are rotation Lie derivatives on the
appropriate scalar/vector/tensor fields; they span tangential differentiations
uniformly on the collar. They commute with radial dilation and with the
rotationally covariant equations.

Eliminate $H,M$ using the exact equations
\begin{align}
 H[T,z]&=\frac{2\Omega}{A}D_0T-2\Omega^2\partial_i z_i
             +2d_\theta T,\\
 M_i[T,z]&=\frac\Omega A(D_0z_i+2lz_i)
             -\partial_iT-cz_i,
 \qquad D_0=\partial_t-B\cdot\partial .
\end{align}
Substituting these expressions in the exact Hamiltonian and momentum
subsidiary equations gives four second-order equations with common principal
part $D_0^2-A^2\Delta$. Its radial second-order coefficient is
$B_r^2-A^2=-\Omega^2$ exactly. After $T=\Omega^2\tau$,
$z=\Omega\zeta$, multiplication by $\Omega^2$, and introduction of $N$,
the equations have the structural form
\begin{equation}
 -\gamma^2\Omega^3\partial_\Omega N+
 L(\Omega,n)N+K(\Omega,n)v+
 2B_r\gamma\Omega\partial_tN
 +\Omega\mathcal T_1v+\Omega^2\mathcal T_2v=0 .
 \label{eq:normal-eight-structural}
\end{equation}
The matrix coefficients are smooth and bounded up to scri. The operators
$\mathcal T_j$ involve only time and angular differentiations and have order
at most $j$; zero-order terms can be included in them. There are no hidden
negative powers or angular derivatives of $N$. This follows from the exact
second-order equations, not from freezing their principal symbols. The
unscaled vector equations contain $T/\Omega^3$, but the displayed balanced
weights turn that term into a bounded coefficient of $v$ in
\eqref{eq:normal-eight-structural}.

At scri the normal matrix is
\[
 L_0=2A_{\scri}^2lJ_{gg},
\]
where the scalar/normal block of $J_{gg}$ is
\[
 \begin{pmatrix}
 -l-d_\theta/2&-\kappa_1/2\\
 l+\kappa_1\kappa_2/2&2l/3-\kappa_1/2
 \end{pmatrix}
\]
and each transverse entry is $-\kappa_1/2$. The common wave metric makes
$-L_0$ strictly coercive. Thus, in a smaller collar, $L$ is invertible and
$-L$ remains coercive. Define
\[
 F=-L^{-1}K,
 \qquad R_g=N-Fv .
\]
The exact slow equation is
\begin{equation}
 (\Omega\partial_\Omega-F_0)v=R_g+\Omega F_1v,
 \qquad F=F_0+\Omega F_1 .
 \label{eq:normal-slow}
\end{equation}
Equation \eqref{eq:normal-eight-structural} gives a fast equation with leading
operator $\Omega^3\gamma^2\partial_\Omega-L$ on $R_g$. Its remaining terms
are sums of $\Omega$ times time/angular operators of order at most one and
$\Omega^2$ times operators of order at most two on $v,R_g$.
For a fixed-time estimate, the $\Omega\partial_tR_g$ term is placed on the
right. Differentiating $Fv$ produces $\Omega^2FN$ and smooth higher-order
terms, because $\partial_\Omega v=N/\Omega$; these introduce no unrecorded
normal derivative and no worse power.

\paragraph{Indicial data.}
For $a\kappa_1=4$, $\kappa_2=-99/100$, the scalar/normal block of $F_0$ is
\[
 \begin{pmatrix}58/31&-150/31\\24/155&242/31\end{pmatrix};
\]
its eigenvalues are $2,238/31$. The transverse blocks are $2I$.
The root $2$ is semisimple with multiplicity three. Spectral projectors are
smooth rotationally covariant functions of $n$. There is no integer
resonance between the two roots. The leading root-two coefficients satisfy
\[
 (h/l,m_n/l,\tau,\zeta_n)
   =(-2/25,302/75,1,-2/75)\tau,
 \qquad m_A=-l\zeta_A.
\]

\paragraph{Two one-dimensional estimates used below.}
First, the fast operator has a fixed-time inverse on every fixed power
weight. For
\[
 \Omega^3a(\Omega)\partial_\Omega u+D(\Omega)u=f,
 \qquad a>0,\quad\operatorname{sym}_HD\ge d_*H>0,
\]
integration by parts against $\Omega^p u^*H$ absorbs the coefficient-derivative
error, bounded by $C_p\Omega_0^2\|u\|_p^2$, in a sufficiently small collar.
The outer endpoint term is nonnegative. The inner growing homogeneous branch
is excluded by membership in a power-weighted $L^2$ space. Hence the inverse
is bounded in $L^2(\Omega^p d\Omega)$ for every fixed real $p$. The collar
may be chosen once for the finite set of weights used below. For constant
scalar $D$ the formula is
\[
 u(\Omega)=e^{D/(2\Omega^2)}
 \int_0^\Omega e^{-D/(2s^2)}f(s)\,s^{-3}\,ds,
\]
whose kernel mass is $1/D$. This fixes the sign of the admissible inverse.

Second, diagonalizing $F_0$ reduces
$(\Omega\partial_\Omega-F_0)v=f$ to scalar Euler equations.
Their Green operators are bounded between equal $b$-weights whenever the
weight differs from an indicial root, with constants controlled by the
inverse distance to that root. When improving a weight across an indicial
root, the corresponding $\Omega^\lambda$ term is extracted. This is obtained
directly by integrating
$\partial_\Omega(\Omega^{-\lambda}v)=\Omega^{-\lambda-1}f$ and applying the
one-dimensional weighted Hardy inequality, or equivalently Young's inequality
after $s=-\log\Omega$. The outer-collar trace is controlled by the local
$L^2$ norms of $v,N$ on a fixed overlap annulus. No independent data at scri
are introduced.

\paragraph{Finite normal-bootstrap statement.}
Let a reduction-compatible linear subsidiary solution have finite, continuous
$L^2$ \emph{joint} graph norms
\[
 \mathcal G_{40}(t)=
 \sum_{p+|I|\le40}
 \|\mathcal Z^I\partial_t^pU(t)\|_{L^2_H}.
\]
Here $\mathcal Z$ contains $\mathcal D=\Omega\partial_\Omega$ on the collar,
the three covariant rotations, and cutoff Cartesian derivatives spanning
the interior, with the overlapping cutoffs defined in the existence section.
All words in these generators are included.  The initial phase therefore
includes conormal regularity \emph{from the outset}; it is propagated by the
first-order commutator proof in that section.  Time derivatives are actual
generator derivatives, not freely prescribed corner jets.  This strengthens
the initial domain.  The earlier assertion that time/rotation-only graph
regularity supplies the conormal levels is withdrawn.
Assuming the exact structural reduction
\eqref{eq:normal-eight-structural}, it has an expansion
\begin{equation}
 v=\sum_{m=2}^7\Omega^m a_m(t,n)+R,
 \qquad \Omega^{-15/2}R\in H_b^6
 \label{eq:normal-finite-expansion}
\end{equation}
locally in the collar, with estimates by $C\mathcal G_{40}(t)$ and the fixed
overlap-annulus norms. The coefficient $a_2$ lies in the root-two kernel;
the integer coefficients $a_3,\ldots,a_7$ are fixed by recursion and the
time/angular derivatives of lower coefficients. The second algebraic branch
$\Omega^{238/31}$ is allowed in $R$.

\paragraph{Proof and derivative bookkeeping.}
Use $b$-measure $d\Omega/\Omega\,d\omega$, and let $\alpha$ denote the
normal weight, so that a weighted norm is the norm of
$\Omega^{-\alpha}u$. Ordinary collar $L^2$ is the starting $b$-weight
$\alpha=-1/2$. Both $v$ and $N$, hence also $R_g$, are bounded linear
functions of $U$ in this norm. The same holds after every finite joint
time/rotation/conormal word in the available budget.  The cutoff Cartesian
generators provide the corresponding ordinary estimates in the interior
and on the overlap annulus.

If the finite joint family of $v,R_g$ is controlled at weight $\alpha$ with
two additional time/rotation generators,
the fast right hand side is controlled at weight $\alpha+1$: its terms have
at least one factor of $\Omega$, and at most two time/angular differentiations.
The fixed-time fast estimate therefore controls $R_g$ at weight
$\alpha+1$.  This is applied \emph{simultaneously} to the required conormal
levels: the source $\Omega\partial_t\mathcal D^jR_g$ is known from the
previous weight at the same conormal level and one higher time level.
It is not used to construct an unknown conormal derivative.
The commutator of $\mathcal D^j$ with $\Omega^3\partial_\Omega$ contributes
$\Omega^2$ times the current $\mathcal D^jR_g$ and lower levels.  The current
term is absorbed into a small smooth change of the coercive fast matrix.
Thus the same fixed-time inverse applies at all these levels.
Equation \eqref{eq:normal-slow} then controls $v$ at the same
improved weight, modulo crossed indicial terms. It also controls $N=Fv+R_g$.
The other fields are reconstructed by
\begin{align}
 Q&=G(2\tau+N_\tau)+\Omega\nabla_{\rm tan}\tau,\\
 Z_{ij}&=G_i(\zeta_j+N_{\zeta_j})
                 +\Omega\partial_{{\rm tan},i}\zeta_j,\\
 h&=\frac{2\Omega}{A}\partial_t\tau+2wN_\tau
                  +(4w+2d_\theta)\tau-2\operatorname{tr}Z,\\
 m&=\frac\Omega A\partial_t\zeta+wN_\zeta
                  +(w-c+2l\Omega/A)\zeta-Q.
\end{align}
Thus the weight improvement returns the full field, not only $v$.

When a known polynomial is present, subtract its recursively determined
fast and slow polynomials before applying this estimate. Smooth coefficients
and time/angular differentiation preserve these finite expansions. At each
new integer order $m\ne2$, $(mI-F_0)$ is invertible; at $m=2$ its kernel is
the outgoing coefficient already identified. The coefficient is bounded by
the Euler Green estimate from the preceding weight and forcing norm.
There are no root-zero or root-one terms because neither is an indicial
root and the initial weight improvements exclude them. Eight improvements
\[
 -\tfrac12\longrightarrow\tfrac12\longrightarrow\cdots
 \longrightarrow\tfrac{15}{2}
\]
extract precisely the powers $\Omega^2,\ldots,\Omega^7$ and cost at most
sixteen time/angular generators. Endpoints never hit an indicial root;
$238/31-15/2=11/62>0$.

Six conormal derivatives are reserved in the starting joint norm and are
carried through the simultaneous weight induction.  The identity
$\mathcal Dv=Fv+R_g$ is compatible with this hierarchy, but is not used to
infer missing conormal norms.  Six angular differentiations and two spare
time differentiations give the conservative budget
$16+6+6+2=30<40$.  At each weight gain the remaining time/rotation budget
decreases by at most two while the reserved conormal levels remain present.
This removes the circular step that would result from putting
$\Omega\partial_t\mathcal DR_g$ on the right before knowing
$\mathcal DR_g$ at the preceding weight.  No infinite regularity or
unrecorded negative-weight commutator is used.

The use of fixed-time fast inversion is essential to this bookkeeping.
If instead the $\Omega\partial_tR_g$ term were kept on the left in a
space-time estimate, an initial half-weight trace would appear and would
require a separate estimate. It is not silently included here.

\paragraph{Ordinary finite regularity and scope.}
After factoring $v=\Omega^2v_{\rm nat}$, the remainder in
\eqref{eq:normal-finite-expansion} has $b$-weight $11/2$. For
$V=\Omega^{-11/2}R_{\rm nat}$,
\[
 \partial_\Omega^6R_{\rm nat}
 =\Omega^{-1/2}\prod_{j=0}^5
       (\Omega\partial_\Omega+11/2-j)V.
\]
The squared ordinary measure $d\Omega$ is exactly the $b$-measure
$d\Omega/\Omega$ on the right. Mixed derivatives are controlled with the
remaining angular budget. Consequently the natural constraint fields
$(H/\Omega^4,T/\Omega^4,M/\Omega^3,z/\Omega^3)$ have ordinary collar
$H^6$ regularity and finite traces. Their leading trace relations are the
root-two kernel relations above. The generic fractional natural exponent
$176/31$ need not be $H^7$; no infinite smoothness is asserted.

This is a nonsharp graph-domain embedding. If the subsidiary semigroup and
its maximal graph domain have been established separately, its strengthened
joint graph norms
propagate without derivative loss in that chosen topology, and the embedding
gives preservation of these natural falloffs. It does not imply an ordinary
$H^6$-data estimate for arbitrary data, does not replace the exact normal
reduction by its indicial polynomial, and does not establish the full
geometric/elliptic-gauge evolution.


% Independent analytical audit.  This is an input fragment, not a main file.
\section{Exact null-halfspace gauge counterexamples}
\label{sec:gauge-counterexamples}

These calculations concern the conformal-harmonic lapse candidate, including
its constant Bona--Masso variant.  They do not establish an instability of the
original reference lapse, whose pole coefficients are different.  They are
exact for the null-halfspace Minkowski model specified below.  Transferring a
boundary-model obstruction to the original CMC ball requires the localization
argument described at the end of this section.

\subsection{Exact flat background and a constraint-free scalar sector}

Let $x>0$ be the inward normal coordinate and set
\[
 \bar g=dx^2+2\,dt\,dx+dy^2+dz^2,\qquad \Omega=x,
 \qquad \alpha_*=1,\quad \beta_*^x=1.
\]
The inverse conformal metric has $\bar g^{tt}=-1$,
$\bar g^{tx}=1$, and $\bar g^{xx}=0$.  The conformal gradient of $\Omega$
is null and its Hessian vanishes.  The conformal transformation formula
therefore shows that $g_{\rm phys}=x^{-2}\bar g$ is an exact flat vacuum metric.
Its physical mean curvature is $K_*=-3$, whereas $\bar K_*=0$.

Consider the physical infinitesimal diffeomorphism
\[
 \xi^t=-x\psi,\qquad \xi^x=x\psi,\qquad \xi^A=0.
\]
Keeping the coordinate conformal factor $\Omega=x$ fixed gives
\[
 \delta\bar g=-\mathcal L_\xi\bar g+2\psi\bar g.
\]
Write $D=\partial_t-\partial_x$, $G=dx$, and use the flat spatial metric
for all spatial derivatives and trace-free projections.  Direct ADM
decomposition gives
\begin{align}
 h_{ij}&=0,& c&=-2\psi,& u&=xD\psi,&
 v_i&=-\partial_i(x\psi),\label{eq:halfspace-gauge-fields}\\
 a_{ij}&=-\big[x\partial_i\partial_j\psi+
 G_i\partial_j\psi+G_j\partial_i\psi\big]^{\rm TF},
 & k&=-x^2\Delta\psi+x\psi_x+3\psi.
 \label{eq:halfspace-gauge-curvature}
\end{align}
Here $h$ is the unit-determinant spatial metric perturbation,
$c=\delta\chi$, $u=\delta\bar\alpha$, $v=\delta\beta$,
$a=\delta\bar A$, and $k=\delta K_{\rm phys}$.  Set
$\Theta=Z_i=0$, so the physical $\hat K$ perturbation $q$ equals $k$.

These fields solve the linear vacuum Einstein equations because they are a
physical diffeomorphism of an exact flat metric.  The Hamiltonian and momentum
constraints and the spacetime Weyl perturbation vanish identically.  They also
satisfy the metric-derived Dirac freezing condition: with
\[
 (L_{\rm conf}v)_{ij}=\partial_i v_j+\partial_j v_i
 -\frac23\delta_{ij}\partial_kv^k,
\]
one has $L_{\rm conf}v=2a$, hence $\partial_t h=0$ and
$V_i=\partial_jh_{ij}=0$.  The corresponding elliptic shift equation is
$(\Delta v_i+\partial_i\operatorname{div}v/3)=2\partial^ja_{ij}$.

For a regular $\psi$, the boundary values are
\[
 u|_{x=0}=0,\qquad v_A|_{x=0}=0,\qquad v_x|_{x=0}=-\psi.
\]
The inward-normal linear nullness relation is
$v_x=u-(h_{xx}-c)/2$ and is satisfied.  Thus these are genuine coupled
metric/lapse/elliptic-shift perturbations, not a prescribed shift trace treated
in isolation.

The normal conformal-factor perturbation and conformal trace are
\[
 \delta w=u-v_x=x\psi_t+\psi,\qquad
 \delta\bar K_{\hat{}}=\frac{q-3\delta w}{x}
 =-x\Delta\psi+\psi_x-3\psi_t.
\]
Consequently a lapse with constant conformal Bona--Masso coefficient $f>0$,
\[
 Du=-f\,\delta\bar K_{\hat{}},
\]
is equivalent in this sector to
\begin{equation}
 x\big(D^2\psi-f\Delta\psi\big)
 -(1+3f)\psi_t+(1+f)\psi_x=0.
 \label{eq:halfspace-general-f}
\end{equation}
Equivalently,
$D^2\psi-f\Delta\psi=[(1+3f)D\psi+2f\psi_x]/x$.

\subsection{Conformal-harmonic $f=1$: unbounded tangential-frequency growth}

For $f=1$, take a tangential Fourier mode
$\psi=e^{st+iky}F(kx)$, with $k>0$, $s=\sigma k$, and $X=kx$.
Equation~\eqref{eq:halfspace-general-f} becomes
\begin{equation}
 (2\sigma X-2)F'(X)
 +\big[4\sigma-(\sigma^2+1)X\big]F(X)=0.
 \label{eq:harmonic-turning-ode}
\end{equation}
For any positive integer $m$, choose
\[
 \sigma=(2m+3)^{-1/2},\qquad
 F(X)=\exp\!\left(\frac{\sigma^2+1}{2\sigma}X\right)
 (1-\sigma X)_+^m.
\]
For $0<X<1/\sigma$ its logarithmic derivative is
\[
 \frac{F'}F=\frac{\sigma^2+1}{2\sigma}
 -\frac{m\sigma}{1-\sigma X}.
\]
Substitution in~\eqref{eq:harmonic-turning-ode} leaves exactly
$\sigma(2m+3)-1/\sigma=0$.  The function is zero for
$X\geq1/\sigma$ and is $C^{m-1}$ across the turning point.  Its cutoff
therefore introduces no distributional error in the first-order equation;
choosing $m$ sufficiently large also makes all metric and curvature equations
classical to any prescribed finite order.  In particular $F(0)=1$ and
$F'(0)=2\sigma$.  Every boundary condition stated above is satisfied for all
times.

On a tangential torus, normalize the initial amplitude in any fixed finite
Sobolev norm containing the required metric, lapse, and curvature fields.
Choose $m$ once, larger than the derivative count required by that norm.
The same state at time $t$ is multiplied by $e^{\sigma kt}$.  Thus no
frequency-uniform estimate with a growth constant independent of $k$ can hold
for this exact halfspace model in that finite Sobolev class.  Constraint or
Weyl estimates cannot remove the example: all those fields are identically
zero.  Requiring analyticity or imposing an additional gauge domain would
change the problem; neither is part of the proposed free hybrid.

\subsection{Strictly subluminal $f=1/3$: compatible data with singular outflow}

Absence of a decaying growing normal mode for $0<f<1$ is not sufficient to
prove preservation of regular scri.  At the particular subluminal value
$f=1/3$, equation~\eqref{eq:halfspace-general-f} has the exact solutions
\begin{equation}
 \psi(t,x)=\frac{F(t+x/v)}{x},\qquad
 v=1\pm\frac1{\sqrt3},
 \label{eq:subluminal-singular-outflow}
\end{equation}
for arbitrary smooth $F$.  To check this, set
$\lambda=1/v$ and substitute $x^{-1}F(t+\lambda x)$.
The coefficients of $F/x^2$ and $F'/x$ cancel identically.  The coefficient
of $F''$ is $1-2\lambda+(2/3)\lambda^2=0$.

Choose $F\in C^\infty_c((T_0,T_1))$ with $T_0>0$.
Both $\psi(0,x)$ and $\psi_t(0,x)$ are smooth and compactly supported in
$[vT_0,vT_1]$, away from scri.  Every boundary compatibility jet initially
vanishes.  The solution remains zero near scri until the outgoing front
arrives.  At a later time for which $F(t)\ne0$,
$\psi\sim F(t)/x$, so $c=-2\psi$ and the conformal lapse perturbation
are unbounded.  The constraints and Weyl perturbation remain zero in the
interior.

This example is a regularity obstruction, not an all-time regular solution of
the mixed boundary problem: the singular continuation has
$\xi^x=x\psi\to F(t)$ and therefore ceases to be scri-fixing.
In general nullness also fails once the front reaches the boundary.
A proof of nonexistence for the complete elliptic boundary-coupled problem
must address the possible $S\ne0$ spatial York response.  A bounded smooth
York potential cannot cancel the $1/x$ metric trace; in the tangentially
constant halfspace sector its normal component satisfies $S_x''=0$.
The counterexample nonetheless rules out recommending subluminal slicing
merely on a characteristic-count or decaying-normal-mode test.

\subsection{Relation to the original CMC ball}

The halfspace background is an exact flat physical geometry and is the leading
normal model obtained by magnifying the compactified CMC ball near scri.
It is not the exact global CMC reference, whose conformal lapse, shift, and
factor vary with radius.  A full transfer requires tangential localization,
rescaling $x=X/k$, preservation of the elliptic boundary conditions, and a
uniform estimate on the coefficient and localization remainders.  One useful
route is to start with physical diffeomorphisms of the actual Minkowski
reference, preserving zero constraints and zero Weyl exactly, and correct the
metric-derived connection by the elliptic York solve.  The correction must be
estimated rather than assumed negligible.  Bounded reference-source terms
such as those proportional to the reference conformal mean curvature are
lower order in the magnified problem, but this observation alone is not the
completed quasimode argument.

The conclusions justified here are therefore: an exact frequency-uniform
estimate fails for the conformal-harmonic null-halfspace model; the $f=1/3$
model has exact singular outflow from compatible compact initial data; and
neither counterexample is cured by subsidiary constraints or physical Weyl
control.  An exact global CMC-ball instability theorem is not claimed.

% Independent analytical audit; input fragment.
\section{A finite graph hierarchy for the original reference potential wave}
\label{sec:original-gauge-graph}

This section concerns the original-reference pure-gauge normal model, not the
conformal-harmonic counterexample operator.  It provides an a priori graph
estimate and a strong-normal embedding for its homogeneous scalar sector.
Existence by closure of a common invariant graph domain, the full $S\ne0$
elliptic boundary coupling, and transfer to variable CMC coefficients remain
separate tasks.

\subsection{Exact forced scalar equation}

In the null-halfspace notation of Section~\ref{sec:gauge-counterexamples},
let $S_i$ be the spatial covariant York generator, satisfying
$\Delta S_i+\partial_i\operatorname{div}S/3=0$, and write
$\xi^t=S_x-x\psi$, $\xi^x=x\psi$.  Direct physical diffeomorphism gives
\begin{align*}
 h&=-L_{\rm conf}S,&
 c&=\tfrac23\operatorname{div}S-2\psi,\\
 u&=-D S_x+xD\psi,&
 v&=-DS+\nabla(S_x-x\psi),\\
 k&=x\Delta S_x-3\partial_xS_x-x^2\Delta\psi+x\psi_x+3\psi.
\end{align*}
For the original reference-lapse pole $(-q-v_x-u)/x$, with $q=k$ in
the constraint-free sector, substitution gives the exact equation
\[
 P\psi=\frac{\mathcal L S_x}{x}+\frac{2\partial_tS_x}{x^2},\qquad
 \mathcal L=\partial_t^2-2\partial_t\partial_x-\Delta_\perp,
 \qquad P=\mathcal L+\frac2{x^2}.
\]
The singular source has an explicit regular lift.  Since
$[\mathcal L,x]=-2\partial_t$, the substitution
\[
 \psi=\partial_t S_x+\frac{x}{2}\mathcal L S_x+\rho
\]
gives
\begin{equation}
 P\rho=-\frac{x}{2}\mathcal L^2 S_x.
 \label{eq:original-reference-forced-lift}
\end{equation}
Thus smooth York data produce an $O(x)$ bulk source after the two displayed
boundary jets are lifted.  Estimating the York generator through its coupled
boundary evolution is still required; this identity does not assume that
source is independently prescribed.

\subsection{Null energy and spatial dilation}

First consider $P\rho=0$.  On a finite normal collar, or with suitable decay at
the inner end, define
\[
 E[\phi]=\int\left(|\phi_t|^2+|\nabla_\perp\phi|^2+
 \frac{2|\phi|^2}{x^2}\right)\,dx\,dy\,dz.
\]
For smooth regular fields with justified traces, multiplication of
$P\phi=F$ by $2\phi_t$ gives
\[
 \frac{d}{dt}E[\phi]
 =-2\int_{x=0}|\phi_t|^2\,dy\,dz
 +2\operatorname{Re}\int\overline{\phi_t}F,
\]
up to the separately controlled inner-boundary flux.  Consequently
$d\sqrt E/dt\leq\|F\|_{L^2}$.  The energy is degenerate in the normal
derivative; positive potential does not itself supply $\phi_x$.

Use the spatial-time dilation
\[
 B=t\partial_t+x\partial_x.
\]
Its exact commutator is
\[
 [P,B]=2P+2\Delta_\perp.
\]
In particular $P(B\rho)=2\Delta_\perp\rho$.  More generally,
\[
 P B^j\rho=\Delta_\perp\big[(B+2)^j-B^j\big]\rho.
\]
Time and tangential derivatives commute with $P$.  The finite hierarchy
\begin{equation}
 \mathcal G_N(t)=
 \sum_{j+p+|\alpha|\leq N}
 E[B^j\partial_t^p\partial_\perp^\alpha\rho](t)
 \label{eq:original-reference-graph-norm}
\end{equation}
therefore closes: every source contains at most $j-1$ dilation operators and
two tangential derivatives, and its $L^2$ norm is controlled by an energy at
the same index using that energy's tangential derivative.  A finite sum and
Cauchy--Schwarz give
\[
 \mathcal G_N(t)\leq C_{N,T}\mathcal G_N(0),\qquad 0\leq t\leq T,
\]
for this a priori regular class.  The same argument permits independently
bounded commuted sources.  It does not hide a normal derivative in the source.

\subsection{Strong-normal embedding}

The elementary identity
\[
 \phi_x=\frac{B\phi-t\phi_t}{x}
\]
implies
\[
 \|\phi_x\|_{L^2}\leq\frac1{\sqrt2}
 \left(\sqrt{E[B\phi]}+|t|\sqrt{E[\phi_t]}\right).
\]
Thus the finite hierarchy controls a full, unweighted normal derivative, not
merely a conormal derivative.

For any commuted field $P\phi=F$, the equation also gives
\[
 \frac{2\phi}{x^2}=F-\phi_{tt}+2\phi_{tx}+\Delta_\perp\phi.
\]
Apply the previous normal estimate to $\phi_t$.  This controls
$\|\phi/x^2\|_{L^2}$ by finitely many energies with two additional
time/dilation/tangential indices and the known source.  No assertion that a
Hardy inequality alone controls the inverse-square potential is needed.

For example,
\[
 \frac{B^2\rho}{x^2}=
 \rho_{xx}+\frac{\rho_x}{x}
 +\frac{2t\rho_{tx}}x
 +\frac{t^2\rho_{tt}}{x^2}+\frac{t\rho_t}{x^2},
\]
while
\[
 \frac{\rho_x}{x}=\frac{B\rho-t\rho_t}{x^2},\qquad
 \frac{\rho_{tx}}x=\frac{B\rho_t-t\rho_{tt}}{x^2}.
\]
The inverse-square estimates applied to these finitely many commuted fields
give an unweighted $\rho_{xx}$ bound from $\mathcal G_4$ on finite time
intervals.  Similarly $B^3\rho/x^2$ controls $x\rho_{xxx}$ once the lower
terms are estimated.  Thus $\mathcal G_5$ controls the basic pure-gauge
geometric fields
\[
 c=-2\rho,\quad u=xD\rho,\quad v=-\nabla(x\rho),\quad
 a=-[x\nabla\nabla\rho+G\otimes\nabla\rho+
 \nabla\rho\otimes G]^{\rm TF},
\]
with $c$ in $H^2$, $u,v,a$ in $H^1$, and
$k=-x^2\Delta\rho+x\rho_x+3\rho$ in $H^1$, on a bounded collar.
These are finite graph estimates with explicitly stronger initial graph
regularity; they are not estimates in an unchanged minimal NR Sobolev norm.
The corresponding estimates for a general nonzero York lift require its
commuted source in~\eqref{eq:original-reference-forced-lift} to be controlled.

\subsection{Why powers of the time generator alone do not suffice}

At a fixed interior point, suppress tangential frequencies and freeze the
positive potential at $V_0$.  The first-order time generator on $(\rho,u)$
has normal symbol
\[
 \begin{pmatrix}0&1\\-V_0&2i\xi\end{pmatrix},\qquad
 s_\pm=i\xi\pm i\sqrt{\xi^2+V_0}.
\]
For positive large $\xi$, its slow branch is
$s_-=V_0/(2i\xi)+O(\xi^{-3})$, an operator of order $-1$ rather than
a normally elliptic operator.  Thus a claim that finitely many powers of the
time generator automatically control normal Sobolev derivatives is not
justified.  Microlocal slow-branch packets expose this obstruction.  The
independent dilation commutations in~\eqref{eq:original-reference-graph-norm}
are the constructive replacement.

The estimates above are for smooth regular solutions, or limits for which
the flux and commutator calculations are established.  Proving existence on
a common invariant completed graph domain, handling its trace conditions,
and closing the coupled elliptic York boundary dynamics are necessary before
promoting this scalar a priori result to a full well-posedness theorem.

\section{Exact ball Lam\'e Dirichlet-to-Neumann boundary subsystem}
\label{sec:boundary-gauge-dn}

This section retains the original reference lapse with boundary principal
$f=1$ and restoring parameter $\rho=0$, and imposes the proposed additional
$u_{\scri}=0$ condition. It proves its necessary homogeneous boundary
subsystem. It does not, by itself, prove the interior lapse-wave
initial-boundary value problem or its compatibility with all geometric pole
conditions.

Write $l=1/a$ and $A_{\scri}=S/a$. For the plus-Lie pure physical gauge,
let $X=\tau\partial_t+\xi$, $\eta=\xi+B\tau$, and
\[
 \sigma=\frac{\eta\cdot G+Aw\tau}\Omega,
 \qquad u=A D_0\tau-(A+\Omega)\sigma,
 \qquad v=D_0\eta-l\eta-A^2\nabla\tau.
\]
Regularity of $\xi$ at the fixed sphere gives
$\eta_n=-A_{\scri}\tau$ there. Direct differentiation of the numerator of
$\sigma$ gives
\[
 \sigma_{\scri}=\partial_n\eta_n+l\tau+A_{\scri}\partial_n\tau.
\]
Here $\partial_nG_n=-1/(aS)$ and
$\partial_n(Aw)=-2/a^2$ at $r=S$; both terms are needed.
Consequently $u_{\scri}=0$ implies
$\partial_t\tau=\partial_n\eta_n+l\tau$ and hence
\begin{align}
 \partial_t\eta_n&=-A_{\scri}\partial_n\eta_n+l\eta_n,\\
 \partial_t\eta_A&=-A_{\scri}\partial_n\eta_A+l\eta_A
                         -A_{\scri}\nabla_A\eta_n.
 \label{eq:ball-boundary-eta}
\end{align}
The second equation also uses $v_A=0$ and
$\nabla_A\tau=-\nabla_A\eta_n/A_{\scri}$.

These conditions do not automatically imply the other pole compatibilities:
they give
\[
 \mu_{\scri}=-2A_{\scri}\partial_n\tau,
 \qquad v_{n,\scri}=-A_{\scri}^2\partial_n\tau.
\]
If the original lapse/geometric pole equations additionally require
$\mu_{\scri}=v_{n,\scri}=0$, the further condition is
$\partial_n\tau=0$. Its consistency with the interior time-wave equation
is a separate question. The result below establishes only the
Dirichlet-to-Neumann evolution \eqref{eq:ball-boundary-eta}.

\paragraph{Regular Lam\'e extension.}
Let $\eta$ solve
\[
 \mathcal L\eta=\Delta\eta+\tfrac13\nabla\nabla\cdot\eta=0
 \quad\hbox{in }r<S,
\]
with regularity at the center. Normalize $Y_{\ell m}$ in $L^2(\mathbb S^2)$,
put $L_\ell=\ell(\ell+1)$, and use the poloidal ansatz
\[
 \eta=U(r)Y_{\ell m}n+V(r)\nabla_{\mathbb S^2}Y_{\ell m}.
\]
Its divergence coefficient is
$U'+2U/r-L_\ell V/r$. The vector Laplacian coefficients are
\begin{align}
 (\Delta\eta)_n&=U''+2U'/r-(L_\ell+2)U/r^2+2L_\ell V/r^2,\\
 (\Delta\eta)_A&=V''+2V'/r-L_\ell V/r^2+2U/r^2.
\end{align}
For $(U,V)=r^k(U_0,V_0)$, the divergence is
$r^{k-1}[(k+2)U_0-L_\ell V_0]$. Inserting its gradient yields the two
regular branches, for $\ell\ge1$,
\[
 r^{\ell-1}(U_0,V_0)=r^{\ell-1}(\ell,1),\qquad
 r^{\ell+1}(U_0,V_0)=
 r^{\ell+1}((\ell+1)(\ell-6),\ell+9).
\]
The determinant of their two boundary vectors is $2(7\ell+3)$, so both
Dirichlet components determine a unique regular extension. The factor
$\ell-6$ causes no degeneracy of the extension.

Use normalized tangential amplitude $b=\sqrt{L_\ell}\,V(S)$, so that the
boundary $L^2$ norm is proportional to $|U(S)|^2+|b|^2$.
The ordinary normal derivative has matrix
\[
 \mathcal D_\ell=\frac1{S(7\ell+3)}
 \begin{pmatrix}
 6\ell^2+\ell+3&\sqrt{L_\ell}(\ell-6)\\
 -\sqrt{L_\ell}(\ell+9)&8\ell^2+5\ell-3
 \end{pmatrix}.
\]
In the boundary system the tangential component also contains
$\nabla_A\eta_n$. Adding $\sqrt{L_\ell}/S$ to the lower-left entry gives
\begin{equation}
 \mathcal B_\ell=\frac1{S(7\ell+3)}
 \begin{pmatrix}
 6\ell^2+\ell+3&\sqrt{L_\ell}(\ell-6)\\
 6\sqrt{L_\ell}(\ell-1)&8\ell^2+5\ell-3
 \end{pmatrix}.
 \label{eq:ball-poloidal-B}
\end{equation}
Thus $\partial_t(U,b)^T=G_\ell(U,b)^T$, with
$G_\ell=a^{-1}I-A_{\scri}\mathcal B_\ell$.

\paragraph{Spectrum, conditioning, and low modes.}
The trace is $2\ell/S$ and the eigenvalues of $\mathcal B_\ell$ are
\[
 \frac1S(\ell\pm\delta_\ell),\qquad
 \delta_\ell^2=
 \frac{(\ell-1)(\ell^2-4\ell-3)}{7\ell+3}.
\]
For $\ell=2,3,4$, $\delta_\ell$ is purely imaginary and the generator
has strictly negative real part $-(\ell-1)/a$. For $\ell=5$,
$\delta_\ell^2=4/19$; for $\ell=6$, $\delta_\ell^2=1$.
They too have strictly negative generator eigenvalues.

For every $\ell\ge7$ the uniformly positive metric
\[
 H_\ell=\operatorname{diag}
       \left(\frac{6(\ell-1)}{\ell-6},1\right),
 \qquad I\le H_\ell\le36I,
\]
symmetrizes $\mathcal B_\ell$. Moreover
$0\le\delta_\ell\le(\ell-1)/\sqrt7$. Therefore
\begin{equation}
 \|e^{tG_\ell}\|\le
 6\exp\left[-(1-1/\sqrt7)(\ell-1)t/a\right]
 \qquad(\ell\ge7).
 \label{eq:ball-high-modal-energy}
\end{equation}
This uniform estimate excludes a hidden high-frequency condition-number or
Jordan growth. The finitely many diagonalizable stable modes $2\le\ell\le6$
can be absorbed in a fixed constant; all $\ell\ge2$ consequently admit a
uniform $Ce^{-t/a}$ bound.

For $\ell=1$,
\[
 \mathcal B_1=\frac1S
 \begin{pmatrix}1&-1/\sqrt2\\0&1\end{pmatrix},\qquad
 G_1=\frac1a\begin{pmatrix}0&1/\sqrt2\\0&0\end{pmatrix}.
\]
This is a genuine neutral Jordan block, but it is a fixed finite-dimensional
one: $\|e^{tG_1}\|\le1+t/(\sqrt2a)$. It causes no frequency-dependent
ill-posedness. The $\ell=0$ radial regular solution is $\eta=C r n$ and has
zero boundary generator.

For axial modes $W(r)n\times\nabla_{\mathbb S^2}Y_{\ell m}$, the divergence
vanishes and the regular branch is $W\propto r^\ell$. Their generator is
$-(\ell-1)/a$, with neutral $\ell=1$ rotations and decay for $\ell\ge2$.

\paragraph{Boundary semigroup.}
The normalized vector spherical harmonics are orthogonal. The uniform modal
bound, together with the explicit finite-dimensional neutral blocks, defines
a strongly continuous semigroup on every $H^s$ boundary space and gives
\[
 \|\eta(t)\|_{H^s}\le C_s(1+t/a)\|\eta(0)\|_{H^s}.
\]
On the $\ell\ge2$ subspace it decays uniformly exponentially. The high-mode
generator has size comparable to $\ell$ with a uniformly invertible negative
spectrum, so its domain is $H^{s+1}$ and its stable high-mode part has the
usual order-one analytic smoothing estimates. These claims concern the
explicit boundary operator \eqref{eq:ball-poloidal-B}, not an inferred
energy from the principal symbol alone.

The result is a genuine global-ball Dirichlet-to-Neumann boundary estimate.
It establishes that this necessary repaired homogeneous trace subsystem has
no high-frequency growth. Establishing the interior time-wave problem,
the additional $\partial_n\tau=0$ compatibility when required by the pole
equations, and physical/constraint inhomogeneous forcing remains necessary
before claiming a full gauge or geometric initial-boundary value theorem.

\section{Original reference lapse: clock repair and exact ball wave}
\noindent\textbf{Scope.} These are linear pure-gauge identities and energy estimates. The condition $u|_{\scri}=0$ is an additional lapse/clock boundary gauge, not a consequence of nullness. No full physical-plus-gauge Z4c theorem is claimed. The transfer of the earlier odd estimate to the original damping branch remains conditional on an estimate for $z/\Omega^3$; a base estimate for $z/\Omega$ does not supply it.

\subsection{Background, driver, and pure-gauge variables}
Let
\[
 \Omega=\frac{S^2-r^2}{2aS},\quad A=\frac{S^2+r^2}{2aS},\quad
 B=-x/a,\quad G=\nabla\Omega=-x/(aS),\quad d=1/(aS),\quad C=A+\Omega=S/a.
\]
Write $D=\partial_t-B\cdot\nabla$ and
\[
 w=-\frac{B\cdot G}{A}=-\frac{r^2}{a^2SA}.
\]
Useful exact relations are
\[
 \nabla A=-G,\quad \nabla G=-dI,\quad D\Omega=Aw,\quad DA=-Aw,
 \quad DG=G/a,\quad A^2-|B|^2=\Omega^2,
\]
\[
 w+1/a=\Omega/(aA),\qquad D^2\Omega=2Aw/a.
\]
The original reference driver used here has harmonic principal coefficient and zero extra restoring parameter:
\begin{equation}
 Du=-G\cdot v-\frac{A^2q}{\Omega}-\frac{AG\cdot v}{\Omega}
       +\frac{Aw}{\Omega}u.\label{driver}
\end{equation}
Here $q=\delta\widehat K_{\rm phys}$; in the constraint-free pure-gauge sector it equals $\delta K_{\rm phys}$.

Use the plus-Lie physical spacetime diffeomorphism $X=\tau\partial_t+\xi^i\partial_i$ and compensate the conformal factor to keep the coordinate function $\Omega$ fixed. Define
\[
 \eta=\xi+B\tau,\qquad \sigma=\frac{\xi\cdot G}{\Omega},\qquad
 \phi=\frac{A\tau}{\Omega},\qquad H=\eta\cdot G,\qquad R=H/\Omega.
\]
The twice-conformal metric perturbation is
\[
 h_{ij}=2\partial_{(i}\eta_{j)}^{\TF},\qquad
 V^i=(\Delta\eta^i+\tfrac13\partial^i\operatorname{div}\eta)=:(L\eta)^i.
\]
Thus metric-Dirac freezing is an elliptic equation for $\eta$, with its prescribed coordinate-time memory as source. Direct metric variation gives
\begin{align}
 q&=-\left(\Delta_{\Omega^{-2}\delta}-3/a^2\right)\phi,\\
 u&=\Omega D\phi-CR,\\
 v&=D\eta-\eta/a-A\Omega\nabla\phi-CG\phi.
\end{align}
The regular normal gauge variable is
\[
 f=\phi-aR=\frac{A\tau-a\eta\cdot G}{\Omega}.
\]
There is a completely regular reconstruction, not merely cancellation of leading residues. With $W_\eta=D\eta-\eta/a$,
\begin{align}
 \tau&=(\Omega f+aH)/A,\quad \xi=\eta-B\tau,\quad \sigma=H/A+wf,\\
 v&=W_\eta-A\Omega\nabla f-CGf-aA\nabla H-aHG,\\
 q&=-J_f f-a[\Omega\Delta H-3G\cdot\nabla H-3dH],
 \quad J_f=\Omega^2\Delta-\Omega G\cdot\nabla-3/a^2,\\
 (a_g)_{ij}&=-[\Omega\partial_i\partial_jf+G_i\partial_jf+G_j\partial_if]^{\TF}
             -a(\partial_i\partial_jH)^{\TF}.
\end{align}
In the last identity, substituting $\phi=f+aH/\Omega$ into the physical normal-deformation curvature makes all inverse powers cancel; the remaining term proportional to $H\operatorname{Hess}\Omega/\Omega$ is pure trace. Thus the reconstructed metric, lapse, shift, physical trace and trace-free curvature have no hidden normal poles.
In particular,
\begin{equation}
 c=-\tfrac23\operatorname{div}\eta-2f/a,
 \qquad u=\Omega Df+aG\cdot(D\eta-\eta/a).\label{regularfields}
\end{equation}

\subsection{An extra boundary clock is necessary}
The flat local model $\bar g=dx^2+2dt\,dx+dy^2+dz^2$, $\Omega=x>0$ is itself a conformal representation of physical Minkowski spacetime: $\bar g$ is flat, $\bar\nabla\bar\nabla\Omega=0$, and $|d\Omega|_{\bar g}^2=0$.

In this paragraph use minus-Lie convention, $D=\partial_t-\partial_x$, and write $\tau=S_x-xf$, $\eta_x=S_x$. The pure-gauge fields include
\[
 u=xDf-DS_x,\quad c=\tfrac23\operatorname{div}\eta-2f,\quad
 v=-D\eta+\nabla S_x-\nabla(xf),\quad h=-2\partial_{(i}\eta_{j)}^{\TF}.
\]
The original driver reduces to
\begin{equation}
 \left(\Box-2/x^2\right)f=\frac{\Box S_x}{x}-\frac{2(S_x)_t}{x^2},
 \qquad \Box=-\partial_{tt}+2\partial_{tx}+\Delta_{\tan}.\label{flatforce}
\end{equation}
The regular boundary clock is $a=S_x|_{x=0}=\tau|_{x=0}$. Nullness alone does not set $a$ to zero.

For any real $s>0$ and tangential frequency $k>0$, zero tangential shift allows
\[
 S_x=(a+bx)e^{-kx},\qquad b=\frac{k(s+2k)}{7s+8k}a,
\]
with the common factor $e^{st+iky}$. Equation \eqref{flatforce} becomes
\[
 2s f'-(s^2+k^2+2/x^2)f=\mathcal R(x),
\]
\[
 \mathcal R=e^{-kx}\left[-(s+k)^2b-\frac{(s+k)^2a}{x}-\frac{2sa}{x^2}\right].
\]
With $\mathcal H(x)=\exp((s^2+k^2)x/(2s)-1/(sx))$, the solution
\[
 f(x)=-\mathcal H(x)\int_x^\infty\frac{\mathcal R(y)}{2s\mathcal H(y)}\,dy
\]
decays at infinity and has a smooth endpoint expansion beginning $f=sa+O(x)$. It supplies regular pure-gauge positive-growth modes unless the clock is additionally closed. This is separate from all Einstein constraints.

The implementable field condition $u|_{x=0}=0$ gives $a_t=(S_x)'|_{0}$. If $p=\eta_y|_0$ is the longitudinal tangential gauge trace, the coupled principal boundary system is
\[
 \partial_t\binom ap=\frac{k}{7}
 \begin{pmatrix}-6&-i\\6i&-8\end{pmatrix}\binom ap.
\]
Its rates are $-k\pm k/\sqrt7$. The positive symmetrizer $\operatorname{diag}(6,1)$ makes the generator strictly negative for $k>0$. The transverse rate is $-k$.

\subsection{Exact ball boundary system}
Regularity of the fixed-$\Omega$ spacetime diffeomorphism gives
\[
 \eta_n=-A\tau\quad\hbox{at }r=S.
\]
The additional condition $u|_{\scri}=0$ and zero tangential shift give
\begin{align}
 (\eta_n)_t&=-A\partial_n\eta_n+\eta_n/a,\\
 (\eta_A)_t&=-A\partial_n\eta_A+\eta_A/a-A\nabla_A\eta_n.
\end{align}
These equations couple to the elliptic Dirichlet-to-Neumann map of $L\eta=V_{\rm memory}$. The principal polar symbol is the stable matrix above. They are independent of the interior regular variable $f$.

\subsection{Exact forced ball wave and cancellation of the cubic pole}
Substitution in \eqref{driver} gives
\begin{equation}
 P\phi=\frac{2CA}{a\Omega^3}H,\qquad
 P=D^2-A^2\Delta-AG\cdot\nabla+\frac{V_0}{\Omega^2},\label{Pphi}
\end{equation}
where
\[
 V_0=3A^2/a^2-C^2|G|^2
 =\frac{3S^4+2S^2r^2+3r^4}{4a^4S^2}>0.
\]
For $f=\phi-aH/\Omega$ the exact source is
\begin{equation}
 Pf=S_f=\frac{J}{\Omega^2}-\frac{a\mathcal Q}{\Omega},\label{Sf}
\end{equation}
with
\begin{align}
 \mathcal Q&=D^2H-A^2\Delta H-AG\cdot\nabla H-2|G|^2H,\\
 J&=2aA(wDH-AG\cdot\nabla H)
 +\left[\frac{A(2A+\Omega)}S+\frac\Omega a\right]H.
\end{align}
In particular all $\Omega^{-3}$ terms cancel algebraically before imposing boundary conditions. A useful identity in this cancellation is
\[
 \frac{2CA}{a}-aV_0=\frac{\Omega(2A-\Omega)}a.
\]

\subsection{The first two boundary lifts}
Let $E$ be a fixed smooth extension from the outer sphere into a collar, constant along the radial coordinate near $r=S$. Under the additional boundary clock system,
\[
 DH|_{\scri}=2H/a,\quad
 J|_{\scri}=2A^2H_n-2AH/a,\quad
 V_0|_{\scri}=2A^2/a^2.
\]
The required leading trace of $f$ is therefore
\begin{equation}
 f_0=\frac{J|_{\scri}}{V_0|_{\scri}}
 =a^2H_n-\frac aA H=-a\partial_n\eta_n.\label{f0}
\end{equation}
This leading lift enforces the remaining original lapse pole condition. Indeed
\[
 \tau=\frac{\Omega f+aH}{A},\qquad
 \tau_n|_{\scri}=\frac{-f/a+aH_n}{A}-\frac H{A^2}=0
\]
when $f|_{\scri}=f_0$. Thus $\tau_n=0$ is derived, not silently assumed from the parabolic clock equations.

Write $\partial_\Omega=(\partial_n\Omega)^{-1}\partial_n$ near the outer sphere and choose
\begin{equation}
 f_1=\frac{\partial_\Omega[J-V_0E f_0]|_{\scri}-a\mathcal Q|_{\scri}}
 {V_0|_{\scri}}.\label{f1}
\end{equation}
Because $\partial_\Omega E f_0=0$ there and
\[
 (V_0)_n|_{\scri}=4A/a^3,
\]
this can equivalently be written
\[
 f_1=\frac{-aJ_n+4Af_0/a^2-a\mathcal Q}{2A^2/a^2}\bigg|_{\scri}.
\]
Set $f_L=E f_0+\Omega E f_1$. Then
\begin{equation}
 P(f-f_L)=F_{\rm reg}:=
 \frac{J-a\Omega\mathcal Q-V_0 f_L}{\Omega^2}
 -(D^2-A^2\Delta-AG\cdot\nabla)f_L.\label{Freg}
\end{equation}
The numerator of the quotient and its first $\Omega$ derivative vanish at scri by \eqref{f0}--\eqref{f1}. Hence $F_{\rm reg}$ is smooth for smooth compatible $\eta$, with no hidden normal poles. This is a genuine finite-jet lift, not a further assumed falloff for $f$.

For an explicit conservative regularity count on a compact collar, $\eta\in C_t^4H_x^{m+6}$ suffices to bound $F_{\rm reg}$ in $C_t^0H_x^m$: $J$ uses at most one derivative, $\mathcal Q$ two, $f_1$ two boundary derivatives, and applying the differential operator to the lifts costs two more derivatives. The spare derivatives cover traces, the fixed extension, and twice-divisible Hardy quotients. Commuting $j$ time derivatives requires the same count with $C_t^{j+4}H_x^{m+6}$. This deliberately generous count does not claim a sharp or equal-order Sobolev estimate.

\subsection{Positive wave energy after the lift}
For $f_h=f-f_L$, put $f_h=A\chi$. Direct calculation gives
\[
 \Box_{\bar g}\chi+\frac{3\Omega}{aA^3}D\chi-W\chi
 =-F_{\rm reg}/A^3,
\]
\[
 W=\frac{V_0/\Omega^2+S_A}{A^2},\qquad
 S_A=\frac{D^2A-A^2\Delta A-AG\cdot\nabla A}{A}.
\]
The function $S_A$ is smooth and $W$ has positive leading coefficient $2/(a^2\Omega^2)$. The stationary Killing energy has density
\[
 \frac1{2A}\left[\chi_t^2+\Omega^2\chi_r^2
 +A^2|\nabla_{\tan}\chi|^2+A^2W\chi^2\right].
\]
A bounded mass shift makes this positive globally if necessary. The extra first-order term is $O(\Omega)D\chi$, which is controlled by this energy, unlike an arbitrary coefficient times $D\chi$. The null outer flux is nonpositive. For compatible smooth solutions,
\[
 E_\chi(t)\le C_T\left[E_\chi(0)+\int_0^t\|F_{\rm reg}(s)\|_{L^2}^2ds\right].
\]
Stationarity permits the same estimate after commuting time derivatives. This yields a pure-gauge estimate in an explicit mixed high-$\eta$/lower-wave graph norm, provided the elliptic-DN boundary evolution is solved with its compatibility conditions. It is not yet an equal-order physical Sobolev theorem, nor does it include all physical Z4c forcing.
\subsection{Weak energy evolution and conormal--graph domains}
This section specifies the energy construction and the additional derivatives needed for regular geometric fields. It concerns the forced pure-gauge wave above, not the full physical evolution.

Choose a bounded stationary mass shift $M$ so
\[
 c_0(r)^2=V_0+\Omega^2(S_A+MA^2)>0.
\]
The added mass term is moved to a bounded remainder in the equation. On a collar set
\[
 U=(w,Q,G_{\tan},R),\qquad
 w=\chi_t,\quad Q=\Omega\chi_r,\quad
 G_{\tan}=A\nabla_{\tan}\chi,\quad R=c_0\chi/\Omega.
\]
The first-order system has symmetrizer $A^{-1}I$. Its radial principal block on $(w,Q)$ is
\[
 \begin{pmatrix}2B_n&\Omega\\\Omega&0\end{pmatrix},
\]
with eigenvalues $B_n\pm A$; the tangential blocks couple $w$ symmetrically to $G_{\tan}$. The only unbounded multiplier is a skew block on $(w,R)$:
\[
 \frac1\Omega\begin{pmatrix}0&-c_0\\c_0&0\end{pmatrix}.
\]
All other coefficients are smooth and bounded on the collar. In particular, the coefficient of $\chi_r$ in the second-order equation is $O(\Omega)$, so it is a bounded coefficient of $Q$.

For existence one may cap only this skew multiplier by $J/\max(\Omega,\varepsilon)$, keeping its two entries paired and leaving the principal matrices unchanged. Smooth approximations of the cap are equivalent. The capped symmetric systems have a common energy bound independent of the cap. One can first truncate at $\Omega=\delta$ and impose zero incoming energy there, then pass to $\delta\downarrow0$. Weak compactness gives a distributional solution on every interior compact subset and then on the full energy space. In the limit the equation $R_t=c_0w/\Omega$ recovers $\chi_t=w$ with $\chi=\Omega R/c_0$. The differences $Q-\Omega\chi_r$ and $G_{\tan}-A\nabla_{\tan}\chi$ have zero time derivative, hence compatible initial data remain compatible.

Uniqueness and the energy bound do not require integrating $1/\Omega$. Localize the energy away from scri and then let the localization approach it. The singular multiplier cancels exactly in the localized quadratic form. The positive radial principal eigenvalue is $B_n+A=O(\Omega^2)$; its cutoff error is $O(\varepsilon)$, while the negative radial part has favorable sign and can be discarded. Interior Friedrichs approximation then gives the same estimate for weak solutions. This defines the unique strongly continuous energy evolution on the closure of compatible smooth data. The argument also provides the natural maximal dissipative generator after a bounded shift.

Let $\mathcal A$ denote this stationary first-order generator and let
\[
 \mathcal D_0=\Omega\partial_\Omega=\frac{\Omega}{\Omega_r}\partial_r
\]
on the collar, extended by a smooth cutoff into the interior. Let $\mathcal R_\alpha$ be rotational derivatives, acting covariantly on the tangential components of $U$. Exact commutation of the scalar wave operator with rotations supplies the corresponding first-order commutations. On the collar,
\begin{equation}
 [\mathcal A,\mathcal D_0]
 =\mathcal A+\sum_\alpha C_\alpha(r)\mathcal R_\alpha
 +C_1(r)\mathcal D_0+C_0(r),\label{comm}
\end{equation}
where every $C$ is smooth and bounded. To see the singular cancellation directly,
\[
 [J/\Omega,\mathcal D_0]=J/\Omega-J_r/\Omega_r.
\]
For a smooth radial matrix $M_r$, the difference between
$[M_r\partial_r,\mathcal D_0]$ and $M_r\partial_r$ is
\[
 \left[-(\Omega_{rr}/\Omega_r)M_r-(M_r)_r\right]\mathcal D_0.
\]
The angular commutators are bounded angular first-order terms. Cutoff commutators are supported where $\Omega$ is bounded away from zero and are handled by ordinary interior Sobolev estimates.

Define the mixed hierarchy norm
\[
 \mathcal E_N(U)=
 \sum_{j+p+|\alpha|\le N}
 \|\mathcal D_0^j\partial_t^p\mathcal R^\alpha U\|_{L^2(A^{-1}dx)}^2
 +\text{ordinary interior Sobolev terms of order at most }N.
\]
For forced evolution the initial time jets are the affine compatible jets
\[
 U_0=U(0),\qquad U_{p+1}=\mathcal A U_p+\partial_t^pF(0).
\]
They must be used instead of silently replacing them by $\mathcal A^pU_0$. For example a regular forced solution $\chi=O(\Omega^2)$ can have $\mathcal AU_0$ with a boundary trace canceled by $F(0)$; it need not belong to the second homogeneous generator domain.

The joint domain is the completion of smooth compatible state/forcing pairs in the displayed finite hierarchy, including these affine jets. Equivalently, additional finite Taylor lifts can be computed recursively from the positive leading coefficient $V_0|_{\scri}$ so the final forcing vanishes to enough finite order and ordinary homogeneous graph domains can be used. Those are derived lifts, not an assumed disappearance of arbitrary constraint data.

Apply the base weak evolution first to the time and rotational graph variables. Then construct the $j=1,2,\ldots$ conormal levels using \eqref{comm}. Every source replaces a conormal derivative by one time or angular derivative and therefore involves lower $j$ at the same total order, or bounded multiplication at the current level. The system is triangular in $j$ and closes by induction and Gronwall. Distributional commutation is validated on interior compact subsets and extended by the common energy estimate and uniqueness; no uniform conormal commutator bound for the capped multiplier is assumed. Consequently
\[
 \mathcal E_N(U(t))\le C_{N,T}\left[
 \mathcal E_N(U(0))+
 \int_0^t\sum_{j+p+|\alpha|\le N}
 \|\mathcal D_0^j\partial_t^p\mathcal R^\alpha F(s)\|_{L^2}^2ds
 \right].
\]
This is a well-defined mixed graph/conormal energy statement, not a claim that every component has identical ordinary Sobolev order.

The independent conormal derivatives give the geometric normal derivatives needed here. Indeed
\[
 \chi_r=\frac{\Omega_r}{c_0}
 \left[\mathcal D_0R+R-(\mathcal D_0c_0/c_0)R\right].
\]
Thus $\chi_r$ is controlled without dividing an energy variable by $\Omega$. Similarly $\Omega\chi_{rr}$ is a bounded combination of $\mathcal D_0^2R,\mathcal D_0R,R$. Two total conormal/time/angular orders control the gauge metric in $H^1$ and the curvature combinations $\Omega\operatorname{Hess}f+G\otimes\nabla f$ in $L^2$, together with the corresponding lapse derivatives. The already explicit higher-$\eta$ lift norms control the remaining spatial-diffeomorphism pieces. This yields the pure-gauge mixed-regularity estimate on the ball once the elliptic-DN boundary evolution for $\eta$ is supplied.

\subsection{Exact bridge for a known regular physical particular solution}
Let $U_P$ be a fixed-$\Omega$ regular particular solution of the sourced linearized Einstein system, and write its lapse, shift and metric perturbations as $u_P,v_P,h_P,c_P$. Define
\[
 \mathcal L_u(U)=Du+G\cdot v+
 \frac{A^2q+AG\cdot v-Awu}{\Omega}.
\]
The pure-gauge identity, including normalization, is
\[
 \mathcal L_u(U_g)=\Omega\left[P\phi-\frac{2CAH}{a\Omega^3}\right].
\]
Write the known particular residual as
\[
 \mathcal L_u(U_P)=R_P+S_P/\Omega,
 \qquad S_P=A^2q_P+AG\cdot v_P-Awu_P.
\]
The sign of $S_P$ is opposite to the pole numerator on the RHS of the lapse equation. To impose the original driver on $U_P+U_g$, the regular gauge wave becomes
\[
 Pf=\frac{J-S_P}{\Omega^2}-\frac{a\mathcal Q+R_P}{\Omega}.
\]
The same two finite Taylor lifts apply with $J$ replaced by $J-S_P$ and $a\mathcal Q$ replaced by $a\mathcal Q+R_P$. No new cubic pole appears.

The total boundary gauge $u_P+u_g=0$ and zero total tangential shift give the forced elliptic-DN system
\begin{align}
 (\eta_n)_t&=-A\partial_n\eta_n+\eta_n/a+u_P,\\
 (\eta_A)_t&=-A\partial_n\eta_A+\eta_A/a-A\nabla_A\eta_n-(v_P)_A.
\end{align}
In particular $DH|_{\scri}=2H/a-u_P/a$. Suppose the particular solution satisfies its regular null condition and $H_{\rm constraint}=T=0$ at scri. Put $\mu_P=(h_P)_{nn}-c_P$. Then
\[
 S_P|_{\scri}=A^2\mu_P/a+2Au_P/a,
\]
whereas the clock forcing changes $J|_{\scri}$ by $+2Au_P/a$. These terms cancel. The leading forced lift is
\[
 f_0=f_{0,\rm homogeneous}-\frac a2\mu_P.
\]
Therefore
\[
 \tau_n|_{\scri}=-\frac{f_0-f_{0,\rm homogeneous}}{aA}
 =\frac{\mu_P}{2A}.
\]
The pure-gauge normal metric contribution is $\mu_g=-2A\tau_n$, so $\mu_P+\mu_g=0$ and total normal shift vanishes. This demonstrates the needed leading boundary compatibility rather than assuming it.

This bridge reduces physical forcing to known boundary traces and finite jets of the particular solution. It does not construct that sourced particular solution or prove its regularity; those are separate required steps before a full physical-plus-gauge linear theorem can be asserted.
\subsection{Reference $1+\log$-like principal coefficient}
The harmonic-collar calculation extends to the original reference family
\[
 F=A^2+\nu\Omega^2,\qquad \nu\ge0,
\]
where $\nu$ is smooth, radial, bounded and stationary. This restriction includes the reference profiles and the fork's background-evaluated $1+\log$-like factor. In particular the latter has $\nu_{\rm eff}=\nu_{\rm input}(2aS)^2A$. Compact support in the interior is permitted but not needed for the collar estimates.

The lapse residual changes by $\mathcal L_{u,\nu}=\mathcal L_{u,0}+\nu\Omega q$. Because the pure-gauge trace is $q=-J_f\phi$,
\[
 P_\nu=P-\nu J_f
 =D^2-(A^2+\nu\Omega^2)\Delta
 +(\nu\Omega-A)G\cdot\nabla+V+3\nu/a^2.
\]
The exact identity becomes
\[
 \mathcal L_{u,\nu}(U_g)
 =\Omega\left[P_\nu\phi-\frac{2CAH}{a\Omega^3}\right].
\]
There is no new forcing pole when passing to $f=\phi-aH/\Omega$, since
\[
 J_f(H/\Omega)=\Omega\Delta H-3G\cdot\nabla H-3dH
\]
is smooth. Thus
\[
 P_\nu f=S_f+a\nu(\Omega\Delta H-3G\cdot\nabla H-3dH).
\]
The displayed $f_0$ and $f_1$ are unchanged; the additional term in $P_\nu f_L$ is smooth as well. A known physical particular solution adds $\nu\Omega q_P$ to its lapse residual and hence only the smooth term $-\nu q_P$ to the divided wave forcing.

The new effective radial principal coercivity is
\[
 F-|B|^2=(1+\nu)\Omega^2.
\]
After the same rescaling $f_h=A\chi$, the change in the first derivative coefficient is
\[
 \nu\Omega(1+2\Omega/A)G\cdot\nabla\chi,
\]
along with the $-\nu\Omega^2\Delta\chi$ principal change and bounded zeroth-order terms. Rewriting relative to the stationary effective wave metric with lapse $\sqrt F$ and shift $B$ changes only bounded time coefficients and $O(\Omega)$ radial first-derivative coefficients. The inverse-square potential has the same positive leading coefficient at scri.

The first-order energy variables can consequently use
\[
 Q_\nu=\sqrt{1+\nu}\,\Omega\chi_r,
 \qquad G_{\tan,\nu}=\sqrt F\nabla_{\tan}\chi,
\]
with the same $w=\chi_t$ and an adjusted smooth positive $R=c_\nu\chi/\Omega$. The radial principal block is
\[
 \begin{pmatrix}2B_n&\sqrt{1+\nu}\,\Omega\\
 \sqrt{1+\nu}\,\Omega&0\end{pmatrix},
\]
whose eigenvalues are $B_n\pm\sqrt F$. Its incoming eigenvalue vanishes as $O(\Omega^2)$, the singular $(w,R)$ multiplier remains skew, and the conormal--graph commutator argument remains valid. The boundary clock and its leading lifts are unchanged. Thus the pure-gauge result is compatible with a reference $1+\log$-like interior coefficient, rather than requiring harmonic slicing everywhere. This remains a linear statement in the stated mixed norms and does not assert the missing full physical coupling theorem.
\subsection{Phase-norm audit: what the reconstruction estimates do not prove}
The high-auxiliary-regularity to lower-raw-regularity estimates above do not by themselves prove Hadamard evolution on one fixed ordinary raw-state Sobolev phase. In particular, increasing the unweighted conormal order does not automatically increase ordinary normal regularity. For $\chi=\Omega R/c_0$,
\[
 \partial_n^k\chi
 =G_n^k\Omega^{1-k}
 \prod_{j=0}^{k-1}(\mathcal D_0+1-j)(R/c_0)
 +\text{lower smooth-coefficient terms}.
\]
The estimate with two total conormal orders supplies raw metric $H^1$ and curvature $L^2$. A raw metric $H^k$/curvature $H^{k-1}$ estimate for $k\ge2$ additionally requires $\Omega^{1-k}$-weighted conormal controls (or a separately proved graph-to-weight coercivity). Such coercivity has not been established here. A branch $R\sim\Omega^\beta$ has the ordinary threshold $k<\beta+3/2$.

For example the constraint-natural fractional amplitude exponent $\mu=176/31$ allows ordinary $H^s$ only for $s<\mu+1/2=383/62$, unless its coefficient is zero or the explicit fractional term is removed before differentiation. One cannot demand arbitrarily large ordinary $H^{m+6}$ merely by increasing $m$ in the conservative lift estimate.

Two mathematically legitimate phase constructions remain available but require their own final proofs:
\begin{enumerate}
\item Let $X$ be the compatible auxiliary Banach phase, $\mathcal R:X\to\mathcal D'$ the continuous raw-state reconstruction, and $K=\ker\mathcal R$. If $K$ is closed and invariant under the auxiliary evolution, the quotient $X/K$ carries its strongly continuous evolution. The intrinsic raw norm can be defined by
\[
 \|U\|_{\rm quotient}=\inf_{\mathcal R x=U}\|x\|_X.
\]
No bounded projection onto $K$ is needed. Closedness follows from continuity into a Hausdorff distribution space, but invariance must be proved. A possible route is CEFE compatibility plus zero curvature/source implying pure-gauge rigidity, followed by the finite-dimensional Minkowski Killing-initial-data kernel and the pure-gauge uniqueness estimate. This lemma is not supplied here.
\item A fixed polyhomogeneous conormal Fr\'echet phase can tolerate finite derivative loss between seminorms. It must include the actual fractional index set and any generated logarithms, with separate boundary coefficient seminorms and conormal remainders. Continuity of the evolution requires estimates at every seminorm; one finite $H^6$ constraint estimate and one lower geometric reconstruction estimate do not suffice. The finite normal jets used by the lifts can be controlled separately rather than demanding infinitely many ordinary normal derivatives of a fractional branch.
\end{enumerate}
Accordingly, the pure-gauge and reconstruction results in this note must not be promoted to a full same-phase raw-state well-posedness theorem without one of these closing arguments.
\subsection{Finite weighted recovery resolving the phase audit at fixed order}
The missing graph-to-weight estimate in the preceding audit can be supplied at any fixed finite order, provided additional \emph{derived} forcing jets are lifted. This does not assert arbitrary ordinary regularity of a fractional branch.

Continue the Taylor lift through $f_3$. For $j=2,3$, after the lower lifts have been subtracted, choose $f_j$ to cancel the coefficient of $\Omega^{j-2}$ in the source, divided by $V_0|_{\scri}$. The mixed normal-time part contributes only at $\Omega^{j-1}$ and the pure normal principal coefficient is $O(\Omega^2)$, so the recursion is triangular with nonzero diagonal $V_0|_{\scri}$. With $f_L=\sum_{j=0}^3\Omega^jEf_j$, the final source satisfies
\[
 F_3=O(\Omega^2),\qquad F_3/\Omega^2\in L^2.
\]
This is a Hardy-four estimate on the full numerator. It uses at most four ordinary normal derivatives of the particular pole $S_P$, five of the particular metric/lapse, and six of $\eta$. Boundary $f_3$ needs only three normal derivatives of $S_P$; separate time/rotational seminorms provide the tangential regularity of its extension. Hence it does not require ordinary $H^7$ regularity across the fractional branch.

The first-order equations now recover weights algebraically. From
\[
 R_t=c_0w/\Omega
\]
one obtains $w/\Omega$ and $(\mathcal D_0w)/\Omega$ from one and two joint graph orders. The $w$ equation then gives $R/\Omega$, since
\[
 w_r=\frac{\Omega_r}{\Omega}\mathcal D_0w,
 \qquad \Omega Q_r=\Omega_r\mathcal D_0Q,
\]
and all other terms are angular derivatives, bounded coefficients, the source, or $w_t$. The exact gradient compatibilities give
\[
 Q=\Omega\,\text{(smooth linear combination of }\mathcal D_0R,R),
 \qquad G_{\tan}=\Omega\,\text{(smooth angular first-order expression in }R).
\]
Inductively, $\Omega^{-p}R$ costs at most $2p$ joint graph orders and source weights through $F_3/\Omega^{p-1}$. The same statement holds after conormal, time and rotational derivatives. Thus $p\le3$ is controlled by the four-jet lift above. There is no unproved propagation of a new weighted norm: it is recovered at each time from the already propagated joint graph variables and the known weighted source.

A conservative finite choice is:
\begin{itemize}
\item Known physical particular solution: joint time/rotation/conormal derivatives through order $32$, with metric/lapse/shift in ordinary $H^5$ and connection/curvature quantities in ordinary $H^4$.
\item Spatial gauge: joint derivatives through order $20$ in ordinary $H^6$, with its affine elliptic-DN boundary time jets. The one-order elliptic gain in the boundary parabolic forcing uses one additional time derivative of the known particular traces, still below $32$.
\item Gauge wave: affine joint graph/conormal energy through order $12$, after lifting $f_0,f_1,f_2,f_3$, and the corresponding source seminorms through order $12$, including $F_3/\Omega^2$.
\end{itemize}
A conservative anisotropic trace/Hardy audit reserves eight extra joint $\eta$ derivatives beyond the requested wave order, including the $f_1,f_2,f_3$ recursion, the two derivatives in $P$, and boundary time/angular trace margins. The order-$20$ budget therefore covers $12+8$. The Hardy-four estimate differentiates in the normal direction; separate joint norms control tangential derivatives of lift coefficients. It does not require an isotropic $H^4$ norm of every individual boundary lift. Ordinary normal requirements remain $H^5/H^4$ for the particular fields and $H^6$ for $\eta$.

These orders give the raw gauge metric and lapse in $H^4$ and curvature in $H^3$. Indeed four ordinary derivatives require at worst $\Omega^{-3}\mathcal D_0^5R$, costing $6+5=11$ joint orders; order $12$ leaves a margin. The shift receives $H^4$ from its strongly elliptic Dirichlet solve with curvature $H^3$ and the known boundary traces. The total normal shift is zero after the leading lift compatibility already proved.

The total physical trace satisfies $q|_{\scri}=0$: the weighted constraints give $H=T=0$ there, while the forced lift gives total $\mu=0$, and the chosen lapse and nullness give total normal shift zero. Consequently Hardy division gives
\[
 q/\Omega\in H^2.
\]
The regular geometric RHS has $H^2$ regularity, and the same lower-order time-graph estimates control curvature time derivatives in $H^2$. Their difference therefore defines the combined pole RHS in $H^2$, with its well-defined $H^{3/2}$ boundary trace and continuous representative. This concerns the analytically canceled pole combination, not separate numerical evaluation of divergent summands.

All affine initial time-jet conditions are continuous linear conditions in this finite joint norm on the complete upstream-plus-gauge phase. Gradient compatibility, the boundary elliptic-DN conditions, and the four lift traces form a closed subspace. The augmented hierarchy propagates them by distributional commutation and uniqueness; the weighted recovery above then propagates the needed normal jet conditions. A continuous physical reconstruction into this raw class follows.

To conclude a same-phase raw-state theorem one may now use the quotient construction described above, once its Killing-kernel invariance lemma and the upstream compatible particular-solution estimates are supplied. The finite orders in this section specify a concrete phase for that argument. They do not claim equivalence to every ordinary raw $H^4/H^3$ norm, nor admit arbitrary data violating the stated affine asymptotic compatibility.
\subsection{Explicit weighted induction, including the reference $\nu$ profile}
Here is the finite weighted-recovery argument in a form that does not hide a commuted weight. Work on a compact collar where $G_n=\Omega_r\ne0$, with smooth radial $\nu\ge0$. Put $\gamma=\sqrt{1+\nu}$ and use the positive first-order potential coefficient $c_0$ (adjusted by a bounded mass shift). Write $r_*=R/c_0=\chi/\Omega$. The compatible variables satisfy exactly
\[
 w=\Omega\partial_t r_*,\quad
 Q=\gamma\Omega G_n(\mathcal D_0+1)r_*,\quad
 G_{\tan}=\sqrt F\,\Omega\nabla_{\tan}r_*.
\]
Their $w$ equation can be written
\begin{equation}
 w_t=2B_n w_r+\gamma\Omega Q_r+
 \mathcal T G_{\tan}-c_0^2 r_*/\Omega+
 b_w w+b_QQ+b_GG_{\tan}+b_Rr_*+\mathcal F,\label{weightw}
\end{equation}
where $\mathcal T$ is angular first order, all $b$ are smooth bounded, and $\mathcal F=F_3/A$ up to an irrelevant sign convention. Thus
\[
 w_r=G_n(\mathcal D_0+1)\partial_t r_*.
\]
Also $Q_r$ is a smooth combination of $(\mathcal D_0+1)r_*$ and $\mathcal D_0(\mathcal D_0+1)r_*$; consequently $\gamma\Omega Q_r$ has an extra $\Omega$. The angular term and all lower terms except $b_Rr_*$ likewise have an extra $\Omega$ after substituting the compatible formulas.

Let $\|\cdot\|_{\mathfrak H_m}$ denote the sum of $L^2$ norms after at most $m$ joint conormal, time and rotational derivatives. Multiply \eqref{weightw} by $\Omega^{1-p}$. The unknown is $c_0^2\Omega^{-p}r_*$. The advective term is controlled by $\Omega^{-(p-1)}r_*$ with two further joint derivatives. Every spatial term with an extra $\Omega$ requires at most the lower weight $p-2$ and two derivatives. The bounded $b_R$ term requires only the already known weight $p-1$. Since
\[
 \Omega^{-q}\mathcal D_0=(\mathcal D_0+q)\Omega^{-q},
\]
all commuted versions have the same structure. Derivatives of $c_0^{-2}$ and other coefficients are bounded. Hence, with $\Omega r_*$ interpreted as a harmless bounded multiple for the base case,
\begin{equation}
 \|\Omega^{-p}r_*\|_{\mathfrak H_m}
 \le C_{p,m}\left[
 \|\Omega^{-(p-1)}r_*\|_{\mathfrak H_{m+2}}
 +\|\Omega^{1-p}\mathcal F\|_{\mathfrak H_m}\right].\label{induction}
\end{equation}
No current-level unknown is moved to the RHS: the zero-order $r_*$ term is strictly one weight lower than the target. Induction gives
\begin{equation}
 \|\Omega^{-p}r_*\|_{\mathfrak H_m}
 \le C_{p,m}\left[
 \mathcal E_{m+2p}(U)^{1/2}
 +\sum_{j=1}^p
 \|\Omega^{1-j}\mathcal F\|_{\mathfrak H_{m+2(p-j)}}
 \right].\label{inductionfull}
\end{equation}
The same inequality holds with $R$ instead of $r_*$ by smooth multiplication. For $p\le3,m\le5$, the graph index is at most $11$. Every source norm in \eqref{inductionfull} follows from the joint norms of $F_3/\Omega^2$, since $\Omega^{1-j}F_3=\Omega^{3-j}(F_3/\Omega^2)$ and $3-j\ge0$.

The first-order construction, skew multiplier, and this induction are unaffected by $\nu$: the radial coefficient is $\gamma\Omega$, the $R$ time equation remains $R_t=c_0w/\Omega$, and the compatible $Q,G$ retain their extra $\Omega$. Smooth radial profiles cover the stated reference $1+\log$-like family.

For clarity on derivative counts, an ordinary derivative of total order at most four on the scalar metric contribution requires at worst $\Omega^{-3}\mathcal D_0^4r_*$. A curvature derivative of order at most three, or the corresponding lapse derivative, requires at worst $\Omega^{-3}\mathcal D_0^5r_*$. These are bounded by \eqref{inductionfull} at graph order $11$, so order $12$ permits one additional time derivative where required. The known lifts and particular fields are controlled by their separate joint norms. Therefore the total metric/lapse/shift are $H^4$, curvature is $H^3$, and the relevant lower-order time derivatives have the regularity used for the pole-RHS statement.

The shift bound may either be read from its explicit reconstruction and \eqref{inductionfull}, or from its strongly elliptic Dirichlet equation. Its source has $H^2$ regularity from curvature $H^3$; the boundary values are the known compatible gauge traces (zero for the final total shift). Thus it is $H^4$ without an extra angular loss.

Finally the total $q$ is $H^3$ and has zero boundary trace, so $q/\Omega\in H^2$. The total normal-factor perturbation also vanishes at scri because $u$ and $v_n$ do; its quotient has at least $H^2$ regularity. All combined conformal trace and pole expressions can consequently be evaluated through these canceled combinations. Subtracting the regular $H^2$ geometric RHS from the corresponding $H^2$ time derivative gives the full pole quotient in $H^2$. The assertion concerns this combined quotient and its trace; it never permits evaluating separate divergent summands at the boundary.

% Independent analytical audit; input fragment.
\section{Regular Weyl forcing from the subsidiary fields}
\label{sec:weyl-source-energy}

This section is conditional on propagation of the stated strong subsidiary
weights.  It establishes an autonomous, regular, forced spin-two system for
physical radiation.  It does not establish the remaining geometric/gauge
reconstruction or the all-angular elliptic boundary estimate.

\subsection{Physical Ricci source and its conformal order}

Use the $C_{\rm Z4c}=0$ tensor equations with the stated Appendix-C
$\Theta$ modification.  Let $n$ be the physical unit normal, $\bar n=n/\Omega$
the conformal unit normal, and retain the physical spatial covector $Z_i=z_i$.
The linearized physical Ricci projections are
\begin{align}
 R_{nn}&=H+2\Omega^2\partial_i z^i
 -3\kappa_1(1+\kappa_2)T,\label{eq:weyl-ricci-nn}\\
 R_{ni}&=-M_i,\label{eq:weyl-ricci-ni}\\
 R_{ij}&=-2D^{\rm phys}_{(i}z_{j)}+
 \gamma^{\rm phys}_{ij}
 \left[\kappa_1(1+\kappa_2)T
 -\frac23\Omega G^kz_k\right].
 \label{eq:weyl-ricci-spatial}
\end{align}
The normal-normal expression follows from the physical Raychaudhuri equation
and the $q+2T$ evolution.  The trace-free spatial expression follows from the
$A_{ij}$ equation.  Its trace is fixed by the off-shell Gauss identity
$R_{nn}+\gamma_{\rm phys}^{ij}R_{ij}=H$.
In particular
$D^i_{\rm phys}z_i=\Omega^2\partial_i z^i-\Omega G^iz_i$;
replacing this by the conformal divergence would miss the displayed isotropic
term.  The unmodified tensor $\Theta$ equation adds $-6w_*T$ to
$R_{nn}$ and the compensating opposite spatial trace; this has the same
conformal order but is a different source.

Assume
\[
 H=\Omega^4h_4,\quad T=\Omega^4\tau_4,\quad
 M_i=\Omega^3m_i,\quad z_i=\Omega^3\zeta_i.
\]
Then $R_{ab}^{\rm phys}=\Omega^2r_{ab}$ in a regular conformal frame.
For $\kappa=\kappa_1(1+\kappa_2)$ the explicit frame components are
\begin{align*}
 r_{\bar n\bar n}&=h_4+6G\cdot\zeta+
 2\Omega\partial_i\zeta^i-3\kappa\tau_4,\\
 r_{\bar n i}&=-m_i,\\
 r_{ij}&=-2\Omega\partial_{(i}\zeta_{j)}
 -5(G_i\zeta_j+G_j\zeta_i)
 +\delta_{ij}\left[\frac43G\cdot\zeta+\kappa\tau_4\right].
\end{align*}
These expressions contain no inverse conformal factor.  To see the spatial
formula directly,
\[
 D^{\rm phys}_iz_j=\Omega^2\left[
 \Omega\partial_i\zeta_j+4G_i\zeta_j+G_j\zeta_i
 -\delta_{ij}G\cdot\zeta\right].
\]
Coordinate components are regular too, because
\[
 R_{ti}=-\frac{A}{\Omega}M_i+B^jR_{ji},\qquad
 R_{tt}=\frac{A^2}{\Omega^2}R_{nn}
 -\frac{2A}{\Omega}B^iM_i+B^iB^jR_{ij}.
\]

\subsection{Exact conformal Cotton source}

Set
\[
 \ell_{ab}=\frac12\left(r_{ab}
 -\frac16\bar g_{ab}\operatorname{tr}_{\bar g}r\right).
\]
The physical Schouten tensor is $L^{\rm phys}_{ab}=\Omega^2\ell_{ab}$,
because $R_{\rm phys}=\Omega^4\operatorname{tr}_{\bar g}r$.
Define the physical Cotton tensor without an antisymmetrization convention
ambiguity by
\[
 A^{\rm phys}_{bcd}=\nabla^{\rm phys}_c L^{\rm phys}_{db}
 -\nabla^{\rm phys}_dL^{\rm phys}_{cb}.
\]
For $d^a{}_{bcd}=\Omega^{-1}\bar C^a{}_{bcd}$ the divergence law is
\begin{equation}
 \bar\nabla_a d^a{}_{bcd}=J_{bcd},\qquad
 J_{bcd}=\Omega^{-1}A^{\rm phys}_{bcd}.
 \label{eq:weyl-divergence-source}
\end{equation}
The exact linear source is
\begin{align}
 J_{bcd}={}&\Omega(\bar\nabla_c\ell_{db}
 -\bar\nabla_d\ell_{cb})
 +3(G_c\ell_{db}-G_d\ell_{cb})\notag\\
 &-\bar g_{cb}G^e\ell_{de}
 +\bar g_{db}G^e\ell_{ce}.
 \label{eq:weyl-regular-cotton}
\end{align}
For verification, the physical/conformal connection difference gives
\[
 \nabla^{\rm phys}_c(\Omega^2\ell_{db})
 =\Omega^2\bar\nabla_c\ell_{db}
 +\Omega\big[4G_c\ell_{db}+G_d\ell_{cb}+G_b\ell_{dc}
 -\bar g_{cd}G^e\ell_{eb}-\bar g_{cb}G^e\ell_{de}\big].
\]
Antisymmetrizing in $c,d$ proves~\eqref{eq:weyl-regular-cotton}.
Every term is smooth at scri under the stated subsidiary regularity.
The one conformal derivative of $r$ requires correspondingly commuted
subsidiary estimates; the base energy alone is not claimed to provide it.

\subsection{First-order electric/magnetic evolution and null flux}

Decompose the Weyl-symmetry field $d$ into symmetric trace-free tensors
$E_{ij}=d_{i\bar n j\bar n}$ and its magnetic dual $B_{ij}$.
Choose the orientation of the magnetic dual so the flat equations read
$\partial_tE-\operatorname{curl}B=\mathcal J_E$ and
$\partial_tB+\operatorname{curl}E=\mathcal J_B$.
Here
\[
 (\operatorname{curl}Q)_{ij}
 =\epsilon_{kl(i}\partial^kQ_{j)}{}^l,
 \qquad
 (a\mathbin\times Q)_{ij}=\epsilon_{kl(i}a^kQ_{j)}{}^l.
\]
The exact sources can be specified invariantly as the STF projections of
$J_{i j\bar n}$ and of the right-dual source
$J^\star_{bcd}=\epsilon_{cd}{}^{ef}J_{bef}/2$, respectively, with the same
orientation.  This definition fixes their signs independently of coordinate
epsilon conventions.

On the CMC conformal Minkowski background, the Cartesian spatial frame is
Fermi-propagated along $\bar n$.  Its acceleration is
$a_i=\partial_i\log A=-G_i/A$ and
$\bar K_{ij}=-\delta_{ij}/(aA)$.  Consequently the evolution can be written
\begin{align}
 (\partial_t-B^k\partial_k)E
 &=A\operatorname{curl}B-\frac3aE
 +2A(a\mathbin\times B)+A\mathcal J_E,\\
 (\partial_t-B^k\partial_k)B
 &=-A\operatorname{curl}E-\frac3aB
 -2A(a\mathbin\times E)+A\mathcal J_B.
\end{align}
Equivalently these are the STF evolution projections of
\eqref{eq:weyl-divergence-source}; the latter remains the convention-independent
exact definition.  All zeroth-order frame coefficients are bounded on the
closed ball, since $A>0$ there.

The identity on the ten electric/magnetic components is a common
symmetrizer, using the Euclidean STF tensor inner product.  In direction $n$,
the symmetric curl block has eigenvalue magnitudes $1$, $1/2$, and $0$;
the vector/scalar components are retained, not projected away.  Hence its
normal operator $\mathcal M_n$ has operator norm one.  At scri,
$B^n=-A$, and therefore
\[
 B^n I+A\mathcal M_n=A(\mathcal M_n-I)\leq0.
\]
The boundary flux of the full electric/magnetic energy is nonpositive.
There is no incoming Weyl mode and no physical radiation boundary value is
freely prescribed.  The bounded connection terms and regular source give
the usual finite-time forced energy estimate.

\subsection{Divergence constraints and source integrability}

The remaining projections of~\eqref{eq:weyl-divergence-source} are the
electric and magnetic divergence constraints.  In the isotropic CMC
background their leading expressions are
\[
 \mathcal C_{E,i}=\partial^jE_{ij}-J_{\bar n i\bar n},\qquad
 \mathcal C_{B,i}=\partial^jB_{ij}-J^\star_{\bar n i\bar n},
\]
or, without frame convention dependence, the corresponding normal
projections of $\bar\nabla_a d^a{}_{bcd}-J_{bcd}$.
Their propagation has principal part
\[
 (\partial_t-B\cdot\nabla)\mathcal C_E
 =\frac A2\operatorname{curl}\mathcal C_B+\text{bounded lower terms},
\]
with the opposite curl sign for $\mathcal C_B$.
The source must satisfy its Cotton identities.  At linear order about the
flat physical background, the contracted Ricci Bianchi identity implies
$\nabla_{\rm phys}^bA^{\rm phys}_{bcd}=0$; commuting derivatives introduces
no background physical curvature.  The cyclic and trace identities follow
from the Schouten definition and the same contracted Bianchi identity.
Their conformal versions are precisely the compatibility conditions needed
for homogeneous divergence-constraint propagation.  They are consequences
of the closed subsidiary evolution and compatible initial Ricci data, not
conditions satisfied by arbitrary independently prescribed $r_{ab}$.

The divergence modes have principal speeds of magnitude at most $A/2$
relative to the normal flow, so they too are outflow at scri.  Initial Weyl
Gauss/Codazzi compatibility is still necessary.  An independently evolved
regular Bianchi field does not by itself prove that a metric reconstructed
from it solves the chosen Z4c/gauge system.

Finally, $\delta d$ is gauge invariant here because the background Weyl
tensor vanishes.  The source depends only on the propagated subsidiary
fields.  This yields the physical triangular estimate
``subsidiary source $\longrightarrow$ rescaled Weyl radiation'', independent
of lapse/shift perturbations at linear order.  The remaining geometric and
gauge reconstruction must be supplied separately.

\subsection{Exact York identity and the off-constraint electric correction}

Here set the metric-derived Dirac field $\partial^jh_{ij}$ to zero;
nonzero prescribed Dirac memory can first be separated by its own elliptic
York potential.  Put $d_0=(aS)^{-1}$, avoiding confusion with the Weyl field.
Define
\begin{align*}
 k_c&=\frac a2\Omega^2\Delta c-
 \frac a2\Omega G\cdot\nabla c-\frac{3c}{2a},&
 k_h&=\frac{3a}{2}h_{ij}G_iG_j,\\
 (a_c)_{ij}&=\frac{a\Omega}{2}(\partial_i\partial_jc)^{\rm TF}
 +\frac a2(G_i\partial_jc+G_j\partial_ic)^{\rm TF},\\
 (a_h)_{ij}&=-a C_h^k{}_{ij}G_k+ad_0h_{ij}
 -\frac{a\Omega}{2}\Delta h_{ij},
 \qquad
 C_h^k{}_{ij}=\tfrac12(\partial_i h_{jk}+\partial_jh_{ik}-\partial_kh_{ij}).
\end{align*}
The trace and momentum identities are
\[
 k=k_c+k_h-\frac a4H,\qquad
 \Omega\partial^j(a_c+a_h)_{ij}
 -2(a_c+a_h)_{ij}G^j=\frac23\partial_i(k_c+k_h).
\]
Consequently the residual $\mathcal R=a-a_c-a_h$ obeys
\[
 \partial^j(\Omega^{-2}\mathcal R_{ij})
 =\Omega^{-3}\left(M_i-\frac a6\partial_iH\right).
\]
A flat York decomposition $\mathcal R=\Omega^2(L_{\rm conf}W+T_{\rm TT})$
therefore gives
\[
 \Delta_LW=m-\frac a2G h_4-\frac a6\nabla(\Omega h_4),\qquad
 \Delta_L=\Delta+\tfrac13\nabla\operatorname{div}.
\]
The gradient term introduces no necessary derivative loss: using
$\Delta_L\nabla\phi=(4/3)\nabla\Delta\phi$, it has the particular
potential $-(a/8)\nabla\Delta^{-1}(\Omega h_4)$, with the appropriate
boundary lift.  Elliptic boundary conditions and kernel control still need
to be specified.

The free $T_{\rm TT}$ is physical information not bounded by $H,M$ alone.
Its relation to the Weyl field has an important off-constraint correction.
Let
\[
 Q_{ij}=[R^{(3)}_{ij}+K K_{ij}-K_{ik}K^k{}_j]^{\rm TF}.
\]
The Gauss decomposition gives
$E^{\rm phys}_{ij}=Q_{ij}-(1/2)(R^{\rm phys}_{ij})^{\rm TF}$.
Direct conformal spatial geometry gives
$\mathcal R=-a\Omega Q$.  Therefore, using
$(R^{\rm phys}_{ij})^{\rm TF}=-2(D^{\rm phys}_{(i}z_{j)})^{\rm TF}$,
\[
 \boxed{\quad
 \frac{\mathcal R}{\Omega^2}
 =-a E+\frac a\Omega
 (D^{\rm phys}_{(i}z_{j)})^{\rm TF} .\quad}
\]
Here $E=E^{\rm phys}/\Omega$ is the regular electric component of the
rescaled Weyl tensor used above.  The last term is $O(\Omega)$ under
$z=\Omega^3\zeta$.  Thus Weyl energy controls the missing weighted TT
curvature, with a known regular subsidiary correction.  Omitting that
correction would silently identify Gauss electric curvature with Weyl
curvature off the Einstein constraint surface.

% Stronger explicit phase repairing a potential normal-bootstrap gap.
\section{A subsidiary phase with explicit conormal control}
\label{sec:subsidiary-conormal-phase-audit}

The time/rotation-only graph norm in the earlier normal-bootstrap statement
does not, without a further argument, supply conormal derivatives of the
fast variable.  In particular commuting its fast equation produces
$\Omega\partial_t(\mathcal D R_g)$, so applying a fixed-time inverse and
then claiming an additional time derivative is not by itself a construction
of $\mathcal D R_g$.  A direct finite strengthening of the initial phase
removes this issue without changing the evolution equations.

Let $\mathcal D=b(r)\partial_r$, $b=\Omega/\Omega_r$, in a collar, cut off
smoothly away from the centre.  Let $\mathcal Z$ consist of $\mathcal D$,
the three covariant rotations, and three cutoff Cartesian derivatives
$\chi\partial_i$ supported away from scri and equal to $\partial_i$ near
the centre.  Choose overlapping cutoffs so these span ordinary derivatives
uniformly on compact interior regions.  Include all words in the norm
\[
 \sum_{p+|I|\le N}
 \|\mathcal Z^I\partial_t^pU\|_{L^2_H}.
\]
This includes ordinary interior Sobolev norms through order $N$.
Time jets are the actual compatible limiting-generator jets.  This is a
stronger explicitly stated phase, not a claim that the old time graph norm
automatically bounds the additional terms.

To prove invariance write
\[
 \mathcal A=V^r\partial_r+V^A\partial_A+R,
 \qquad R=\frac{A_*}{a\Omega}J_0(n)+R_b,
\]
where the principal coefficients, $R_b$, and $J_0$ are smooth on the collar.
The exact identity is
\[
 [\mathcal A,\mathcal D]
 =b'\mathcal A-(b'V^A+b\partial_rV^A)\partial_A
  -(\partial_rV^r)\mathcal D-b'R-b\partial_rR.
\]
Crucially its last apparent pole is removable:
\[
 b'\frac{A_*J_0}{a\Omega}
 +b\partial_r\frac{A_*J_0}{a\Omega}
 =\left(\frac{(b'-1)A_*}{a\Omega}
       +\frac{\partial_rA_*}{a\Omega_r}\right)J_0,
 \qquad b'-1=-\frac{\Omega\Omega_{rr}}{\Omega_r^2}.
\]
It is smooth and bounded.  Radial differentiation leaves $n$ fixed.
Rotations are the covariant rotations under which the entire tensor system
is equivariant.

Set $Y_j=\mathcal D^jU$.  Replace $\mathcal AU$ in the commutator by $U_t$.
At each fixed total time/rotation/conormal order the equation for $Y_j$
has the original generator plus bounded multiplication at level $j$;
its forcing consists of time and rotation derivatives at strictly lower
conormal order and the same total order.  Thus the hierarchy is triangular.
Construct its $L^2$ solutions by the subsidiary semigroup and bounded
perturbations, inducting on $j$.  On every compact interior domain the
distribution $\mathcal D^jU$ satisfies exactly this equation with the same
initial value.  Local hyperbolic uniqueness identifies it with the
constructed $Y_j$.  The backward domains remain inside $\Omega>0$ at finite
time, so this argument needs no boundary regularity of the unknown derivative.
It follows that the strengthened finite joint phase is invariant and has
the corresponding energy estimate.  Completion in this norm supplies a
complete, strongly continuous phase.  It is not necessary to obtain uniform
conormal commutator estimates for the capped approximants.

The cutoff Cartesian commutators are supported where $\Omega$ is bounded
below.  They have bounded first-order coefficients; the overlapping
generator family expresses these as the ordinary same-order interior jet
system, which closes by the standard differentiated symmetric-hyperbolic
energy.  No additional unlisted interior regularity is used.

The fast/slow weight improvement can now be applied simultaneously to the
finite set $\mathcal D^j\partial_t^p\operatorname{Rot}^\alpha(v,R_g)$.
Their starting weighted $L^2$ norms are available by construction.  For the
fast operator, the commutator of $\mathcal D^j$ with
$\Omega^3\partial_\Omega$ contributes $\Omega^2$ times the current
$\mathcal D^jR_g$ and lower levels.  These are smooth small changes in its
coercive zeroth-order matrix.  The term
$\Omega\partial_t\mathcal D^jR_g$ on the right is now genuinely known at
the previous weight from the joint hierarchy, rather than inferred
circularly.  Each weight improvement still costs at most two additional
time/rotation generators.  Reserving six conormal, six angular and two
time orders together with the sixteen weight-improvement orders uses
$6+6+2+16=30<40$.  Hence the normal $H^6$ embedding follows on this
explicitly strengthened phase; using order $128$ in the coupled construction
leaves ample finite joint budget for the subsequent sources.

The compatible core should allow the actual fractional indicial branch;
one may take interior-smooth conormal/polyhomogeneous data with the stated
finite jets and fractional index, then complete.  Density of conformally
$C^\infty$ boundary data in the same maximal high graph phase is not assumed.

% Independent linear reconstruction audit. Input fragment, not a standalone file.
\section{A regular particular metric from the sourced conformal equations}
\label{sec:sourced-cefe-reconstruction}

This section concerns the \emph{linearized} problem only.  It supplies a
particular physical metric, before imposing the final elliptic shift and
lapse gauges.  It assumes the strong subsidiary closure and the compatible
regular rescaled Weyl evolution established in the preceding sections.
It does not infer regularity of the rescaled Weyl tensor from the subsidiary
weights alone.  Initial Weyl peeling and curvature compatibility are explicit
parts of the initial-data class.

Write $g$ for the conformal metric in this section, $\widetilde g=\Omega^{-2}g$
for the physical metric, $\Sigma_a=\nabla_a\Omega$, and $L_{ab}$ for the
\emph{unphysical} Schouten tensor.  The prescribed first-order physical Ricci
source is
\[
 \delta\widetilde R_{ab}=\Omega^2r_{ab},\qquad
 \ell_{ab}=\tfrac12(r_{ab}-\tfrac16g_{ab}\operatorname{tr}_g r),
 \qquad \delta\widetilde L_{ab}=\Omega^2\ell_{ab}.
\]
Every coefficient multiplying a first-order source in this section is evaluated
on the fixed conformal Minkowski background.  Thus changes of coordinates or
conformal gauge do not change $r,\ell,J,d$ at first order: the physical Ricci
and Weyl background fields vanish.

\subsection{The source conservation identity, including its trace}

Put $S_{ab}=r_{ab}-\tfrac12g_{ab}\operatorname{tr}r$.  The connection-change
formula gives, without using a field equation,
\[
 \widetilde\nabla^a(\Omega^2S_{ab})
 =\Omega^4\nabla^aS_{ab}+\Omega^3\Sigma_b\operatorname{tr}S.
\]
Consequently the necessary and sufficient linear physical conservation law is
\begin{equation}
 \Omega\bigl(\nabla^a\ell_{ab}-\nabla_b\operatorname{tr}\ell\bigr)
 =3\Sigma_b\operatorname{tr}\ell.
 \label{eq:sourced-cefe-conservation}
\end{equation}
It is not legitimate to prescribe arbitrary smooth $r$ independently of the
subsidiary system.  Here the $H,M$ evolution equations were obtained from the
physical ADM contracted Bianchi identity with the stated $Q^{\rm phys}_{ij}$.
Substituting their Ricci projections into that identity is precisely
\eqref{eq:sourced-cefe-conservation}; it introduces no new condition on their
solutions.  This explains why the source equations must be kept, rather than
regarding $H,M,T,z$ as unrelated external functions.
For a direct normalization check set $N=A/\Omega$ and
$F_{ij}=Q^{\rm phys}_{ij}/N$.  The prescribed Einstein source has projections
\[
 \delta\widetilde G_{nn}=H/2,\qquad
 \delta\widetilde G_{ni}=-M_i,\qquad
 \delta\widetilde G_{ij}
  =\gamma^{\rm phys}_{ij}(H/2+\operatorname{tr}_{\rm phys}F)-F_{ij}.
\]
Indeed $\widetilde R_{ij}=-F_{ij}$ and
$\widetilde R_{nn}=H+\operatorname{tr}_{\rm phys}F$.
Its normal and tangential divergence projections are exactly the displayed
closed Hamiltonian and momentum subsidiary equations, including the
covector Lie derivative on the latter.  This provides an explicit algebraic
way to verify conservation rather than assuming a physical metric has
already been reconstructed.

The identity also has a nontrivial boundary consequence.  Since $d\Omega\ne0$
at scri and the source is regular, $\operatorname{tr}\ell=O(\Omega)$.  For the
explicit source in the Weyl section,
\[
 \operatorname{tr}r=-h_4-12G\cdot\zeta
       -4\Omega\partial_i\zeta^i+6\kappa\tau_4.
\]
The root-two subsidiary coefficient relation makes this trace vanish at scri.
This is consistent with the full identity, but checking this one leading
coefficient alone would not prove conservation.

The conformal Cotton tensor is the already established regular source
\[
 J_{bcd}=\Omega(\nabla_c\ell_{db}-\nabla_d\ell_{cb})
 +3(\Sigma_c\ell_{db}-\Sigma_d\ell_{cb})
 -g_{cb}\Sigma^e\ell_{de}+g_{db}\Sigma^e\ell_{ce}.
\]
For example its contracted trace vanishes exactly by
\eqref{eq:sourced-cefe-conservation}.  Its remaining integrability identities
follow by commuting the physical derivatives on the flat physical background.

\subsection{Exact regular equations for a traceful source}

The usual trace-free-matter conformal theorem cannot simply be cited for this
source.  Introduce instead
\[
 s=\frac{\Sigma_a\Sigma^a}{2\Omega}
   =s_F-\frac14\Omega^3\operatorname{tr}\ell,
 \qquad s_F=\tfrac14\Box_g\Omega+\tfrac1{24}R[g]\Omega.
\]
The quotient is an identity on $\Omega>0$; in the reduced system $s$ is an
independent regular field.  The exact sourced conformal equations are
\begin{align}
 H_{ab}:=\nabla_a\nabla_b\Omega
  &=-\Omega L_{ab}+sg_{ab}+\Omega^3\ell_{ab},
       \label{eq:sourced-cefe-hessian}\\
 \nabla_as&=-L_{ab}\Sigma^b+\Omega^2\ell_{ab}\Sigma^b,
       \label{eq:sourced-cefe-s}\\
 \nabla_cL_{ab}-\nabla_aL_{cb}
  &=\Sigma_ed^e{}_{bca}+\Omega J_{bca}=:F_{cab},
       \label{eq:sourced-cefe-curl}\\
 \nabla_ad^a{}_{bcd}&=J_{bcd},
       \label{eq:sourced-cefe-bianchi}\\
 R_{abcd}[g]&=\Omega d_{abcd}
    +2\bigl(g_{a[c}L_{d]b}-g_{b[c}L_{d]a}\bigr),
       \label{eq:sourced-cefe-curvature}\\
 2\Omega s-\Sigma^2&=0.
       \label{eq:sourced-cefe-scalar}
\end{align}
Antisymmetrization in the curvature formula has weight $1/2$.
The Schouten conformal transformation proves
\eqref{eq:sourced-cefe-hessian}; differentiation of
\eqref{eq:sourced-cefe-scalar} then proves
\eqref{eq:sourced-cefe-s}.  The Weyl conformal transformation and the two
contracted Bianchi identities prove the curl and divergence equations.
In particular all new terms are $\Omega^3\ell$, $\Omega^2\ell\Sigma$,
and $\Omega J$, with no inverse conformal factor.

\subsection{An explicit regular wave reduction}

Prescribe the smooth unphysical scalar curvature $\mathcal R(x)$ and smooth
generalized harmonic coordinate sources.  Impose $R[g]=\mathcal R$ and use
the usual reduced Ricci operator, so its principal part is
$-\tfrac12g^{cd}\partial_c\partial_dg_{ab}$.
The following equations define a regular reduced system.  In their right-hand
sides replace every curvature by the algebraic expression
\eqref{eq:sourced-cefe-curvature} and every Hessian of $\Omega$ by
\eqref{eq:sourced-cefe-hessian}:
\begin{align}
 R^{\rm red}_{ab}[g]&=2L_{ab}+\tfrac16\mathcal R g_{ab},
        \label{eq:sourced-reduced-metric}\\
 \Box_g L_{ab}&=\tfrac16\nabla_a\nabla_b\mathcal R
   +\nabla^cF_{c(ab)}
   +[\nabla^c,\nabla_{(a}]L_{|c|b)},
        \label{eq:sourced-reduced-L}\\
 \Box_g\Omega&=4s-\tfrac16\mathcal R\Omega
                    +\Omega^3\operatorname{tr}\ell,
        \label{eq:sourced-reduced-Omega}\\
 \Box_gs&=-\tfrac16\Sigma^a\nabla_a\mathcal R-L^{ab}H_{ab}
              +\nabla^a(\Omega^2\ell_{ab}\Sigma^b).
        \label{eq:sourced-reduced-s}
\end{align}
The symmetric $L$ equation follows by taking the divergence of its curl
equation and using $\nabla^aL_{ab}=\nabla_b\mathcal R/6$.
More explicitly,
\[
 \nabla^cF_{cab}
 =H^c{}_e d^e{}_{bca}+\Sigma_e\nabla^cd^e{}_{bca}
    +\Sigma^cJ_{bca}+\Omega\nabla^cJ_{bca}.
\]
Thus the largest source derivatives are first derivatives of $d,J,\ell$.
This formula also specifies the reduction without concealing a second
derivative of $\Omega$ in $\nabla F$.  All these statements concern the
linearization of the displayed system about the fixed background; no
quasilinear conclusion is needed here.

In coordinates the scalar wave operators must be reduced using the same
contracted-Christoffel gauge.  This matters for the mixed regularity below:
$\delta\Box\Omega_*$ then contains only zeroth-order metric perturbations
and prescribed gauge terms.  Leaving an uncontrolled contracted
$\delta\Gamma^a\partial_a\Omega_*$ on the right would falsely remove the
one extra derivative needed for the final conformal-factor gauge change.

\subsection{Mixed estimate and the fixed-$\Omega$ particular metric}

For integer $m\geq3$ use the phase norm
\begin{align*}
 \mathcal E_m={}&\|\delta g\|_{H^{m+1}}+
 \|\partial_t\delta g\|_{H^m}
 +\|\delta L\|_{H^m}+\|\partial_t\delta L\|_{H^{m-1}}\\
 &+\|\omega\|_{H^{m+2}}+\|\partial_t\omega\|_{H^{m+1}}
 +\|\delta s\|_{H^{m+1}}+\|\partial_t\delta s\|_{H^m},
 \qquad\omega=\delta\Omega.
\end{align*}
One sufficient source norm is
\begin{align*}
 \mathcal S_m={}&\|d\|_{H^m}+\|\partial_td\|_{H^{m-1}}
 +\|J\|_{H^m}+\|\partial_tJ\|_{H^{m-1}}\\
 &+\|\ell\|_{H^{m+1}}+\|\partial_t\ell\|_{H^m}.
\end{align*}
The independent Bianchi evolution bounds the displayed time derivative of
$d$ by its spatial derivatives and $J$.
The ordinary energy estimate for the linear reduced waves gives
\begin{equation}
 \mathcal E_m(t)\leq C_{m,T}\left(
 \mathcal E_m(0)+\int_0^t\mathcal S_m(t')\,dt'\right),\qquad 0\leq t\leq T.
 \label{eq:sourced-cefe-mixed-energy}
\end{equation}
To verify the derivative count, the $L$ forcing belongs to $H^{m-1}$;
even the tensor wave operator's second metric derivatives have this order
when $g\in H^{m+1}$.  The metric forcing belongs to $H^m$, the $s$ forcing
to $H^m$, and the $\Omega$ forcing to $H^{m+1}$.  The last assertion uses
the contracted-Christoffel gauge just specified.  The background light cone
is the common wave principal symbol.  At scri $B^n=-A$, so all reduced wave
fluxes are nonpositive.  Standard extension from the ball, or the equivalent
outflow energy construction, gives existence and uniqueness without an
independent outer boundary value.  Commuting finitely many time derivatives
and rotations gives the same estimate with the corresponding joint source
norms.  No norm count is asserted for an infinite commutator hierarchy.

Finite conormal commutations require a separate argument.  In a collar
write each regular first-order wave reduction as
$\mathcal A=V^r\partial_r+V^A\partial_A+R$ and set
$\mathcal D=b(r)\partial_r$, $b=\Omega/\Omega_r$, extended by a cutoff
away from the centre.  All matrices here are smooth and bounded.  The exact
identity is
\[
 [\mathcal A,\mathcal D]=b'\mathcal A
 -(b'V^A+b\partial_rV^A)\partial_A
 -(\partial_rV^r)\mathcal D-b'R-b\partial_rR.
\]
On a solution replace $\mathcal AU$ by $U_t$ minus its known source.
In the equation for $\mathcal D^jU$, the time and angular derivatives then
have strictly smaller conormal order and the same total joint order; the
term involving $\mathcal D(\mathcal D^{j-1}U)$ is bounded multiplication at
the current level.  Rotations span the angular derivatives on the collar.
Induction on conormal order at fixed total order and the same energy estimate
therefore close the joint time/rotation/conormal hierarchy.  The first-order
gradient compatibilities recover the stated ordinary Sobolev orders.
The $\nabla d$ term in the $L$ equation remains an upstream source
($d\in H^m$ gives $\nabla d\in H^{m-1}$); the other mixed-phase couplings
have the bounded orders already checked above.  Cutoff commutators are
supported in the interior and use ordinary interior derivatives.
Include three cutoff Cartesian generators $\chi\partial_i$, equal to the
Cartesian derivatives near the centre and supported away from scri, in the
joint family $\mathcal Z$ along with $\mathcal D$ and rotations.  The cutoffs
overlap so the family spans ordinary derivatives uniformly in the interior.
Their commutators are supported where $\Omega$ is bounded below and are the
ordinary bounded-coefficient differentiated wave system.  Thus the explicit
joint phase is
\[
 \mathcal E_{4,32}
   =\sum_{p+|I|\le32}\mathcal E_4
           (\mathcal Z^I\partial_t^p(\delta g,\delta L,\omega,\delta s)).
\]
The compatible derivative fields and their velocities are interpreted as in
the mixed phase above.  This norm includes all corresponding ordinary
interior derivatives, in particular the metric/lapse interior orders needed
by the gauge-wave source.  The source phase uses the same generator words.
Consequently this finite joint hierarchy is controlled with its listed
source and initial norms.  This trades
conormal order for time/rotation order and does not assert extra ordinary
normal differentiability across a fractional branch.

To return to the original, fixed conformal factor, choose a smooth cutoff
$\chi$ supported in a collar, equal to one near scri, and a smooth spatial
vector $N$ on its support with $N(\Omega)=1$.  Set
\[
 \xi=-\chi\omega N,\qquad
 \vartheta=-\frac{(1-\chi)\omega}{\Omega}.
\]
The quotient is supported where $\Omega$ is bounded away from zero.  The
coordinate and conformal transformation
\[
 \omega'=\omega+\xi(\Omega)+\vartheta\Omega=0,\qquad
 \delta g'=\delta g+\mathcal L_\xi g_*+2\vartheta g_*
\]
preserves the physical metric modulo a regular coordinate transformation.
It is bounded from $\mathcal E_m$ to
$\delta g'\in H^{m+1}$, $\partial_t\delta g'\in H^m$.
The extra derivative of $\omega$ is precisely what controls the Lie
derivative.  At $m=4$ this yields a fixed-$\Omega$ particular conformal metric
$g_P\in H^5$, with its extrinsic-curvature perturbation in $H^4$.
The map to $h_P,c_P,u_P,v_P,a_P,k_P$ is algebraic in $g_P$ and its first
derivatives, hence has these same stated metric/curvature orders.  Finally
$q_P=k_P-2T$ and $\Lambda_P=V(h_P)+2z$ reconstruct the independently evolved
Z4c fields; those two fields must not be discarded during reconstruction.

This reconstruction recovers the full off-constraint geometry, rather than
projecting to vacuum constraints.  For any physical metric the contracted
Gauss and Codazzi identities give
\[
 H_{\rm geom}=\widetilde R_{nn}
       +\gamma_{\rm phys}^{ij}\widetilde R_{ij},\qquad
 M_{{\rm geom},i}=-\widetilde R_{ni}.
\]
The prescribed source therefore makes these exactly the independently
propagated $H,M$.  Since $\widetilde R_{ij}=-Q^{\rm phys}_{ij}/N$, its ADM
extrinsic-curvature evolution is precisely the tensor $C_{\rm Z4c}=0$
evolution with the stated sources.  The geometric metric evolution is the
ADM kinematic identity.  The independent $T$ equation is unchanged, and
the independent $z$ equation, with $\Lambda=V+2z$, is the connection-variable
equation.  The scalar definition $q=K_{\rm phys}-2T$ recovers the trace
equation including the disclosed Appendix-C scalar correction.  Thus these
are the full linearized free Z4c equations in any subsequently chosen
regular physical coordinate gauge.  At this background the gauge changes
leave $H,M,T,z$ invariant, because their background values vanish.

\subsection{Why the reduced solution satisfies the full sourced equations}

Here is a linear proof that avoids assuming a traceful-matter extension of a
published trace-free CEFE constraint theorem.  The initial phase is required
to be a compatible sourced CEFE lift of the physical three-metric and
extrinsic curvature, with the regular rescaled Weyl data.  The harmonic and
scalar-curvature gauge data are chosen consistently.  On any compact subset
of $\Omega>0$, solve the ordinary linear sourced physical Einstein equations
in Lorenz gauge.  Since the physical background is flat, these are
\[
 \widetilde\Box\widehat h_{ab}=-2\Omega^2S_{ab}.
\]
The Lorenz constraint satisfies a homogeneous wave equation by physical
source conservation.  Its initial value is fixed to zero by the four gauge
choices; its initial normal derivative is zero by the sourced Hamiltonian
and momentum equations.  Thus this solution is a genuine physical Ricci
solution on the interior domain of dependence.

Locally put that solution into the prescribed conformal scalar-curvature
and generalized harmonic coordinate gauges.  Infinitesimal conformal changes
obey $\delta_\vartheta R=-6\Box\vartheta-2R_*\vartheta$; infinitesimal
coordinate gauge changes obey the usual vector wave operator
$\Box\xi^a+R_*{}^a{}_b\xi^b$, with a regular first-derivative coupling to
$\vartheta$.  These are ordinary coupled linear waves on each such interior
domain.  Their Cauchy values and first derivatives can be chosen to match
the reduced initial metric and conformal-factor data; the unused normal
gauge derivatives enforce the prescribed gauge sources.  The resulting
genuine rescaled Weyl field satisfies the same regular Bianchi equation with
the same $J$ and the same initial $d$.  Bianchi uniqueness on the interior
domain of dependence therefore identifies it with the independently evolved
field $d$.  The genuine fields satisfy the exact traceful equations above and therefore
the reduced equations with the same initial data.  Uniqueness for that
regular reduced system makes them coincide with its solution.

This interior argument covers every point with $\Omega>0$ at finite time.
The backward background domain of dependence of such a point remains
strictly inside the ball: the glancing radial speed is $O(\Omega^2)$ and
cannot reach $\Omega=0$ in finite backward time.  Thus no unspecified exterior
data enter the comparison.  All sourced CEFE zero quantities vanish in the
interior.  Their regular Sobolev traces then vanish at scri as well.
This proves equivalence for the compatible data class without invoking
a nonexistent arbitrary-trace source theorem.

\subsection{Initial compatibility and what is, and is not, gained}

The mixed initial phase must include a regular curvature lift.  A simple
sufficient way of obtaining $L(0)\in H^4$ is to start with the conformal
three-metric in $H^6$, its conformal extrinsic curvature in $H^5$, and
compatible $d(0)\in H^4$.  Alternatively one can include the compatible
$L(0)\in H^4$ lift explicitly in the phase.  It is incorrect to infer this
order from a bare three-metric in $H^5$, since the Gauss equation has two
spatial derivatives.  The trace and normal derivatives of $L$ are then
fixed by the prescribed scalar curvature, its curl/divergence equations,
and the initial sourced curvature identities.

There is no hidden division in constructing the initial $s$.  If
$K_{\rm phys}=\Omega\bar K+3w$, the physical Hamiltonian identity gives
\[
 s=\frac{\bar\Delta\Omega+\bar K w}{3}
 +\frac{\Omega}{12}
 (R[\bar\gamma]+\bar K^2-\bar K_{ij}\bar K^{ij})
 -\frac{H}{12\Omega}.
\]
The last term is $-\Omega^3h_4/12$ in the present data class.  Initial
$L_{ij},L_{0i}$ can likewise be obtained from the unphysical Gauss--Codazzi
curvature after subtracting $\Omega d$, and $L_{00}$ from its scalar trace;
these operations use derivatives but no inverse conformal factor.

The conclusion of this section is a bounded particular \emph{linear}
physical metric reconstruction in the mixed phase above, including its
constraint-generated Ricci source.  A choice of regular Weyl data is genuine
physical information.  The section does not replace the remaining forced
elliptic-shift/lapse mapping or its initial-data bijection.  The primary
wave-CEFE references
\url{https://arxiv.org/abs/1902.01816} and
\url{https://arxiv.org/abs/2303.07385} explicitly treat trace-free matter;
the traceful equations and the interior-equivalence argument here supply
the additional step needed for this particular prescribed linear source.

% Scope audit and an abstract realization lemma; input fragment.
\section{A same-space realization and the remaining norm audit}
\label{sec:linear-phase-quotient-audit}

A high-initial-norm estimate into a weaker norm is not, by itself, a
Hadamard well-posedness theorem in either norm.  Nor does a regular auxiliary
CEFE metric automatically give a bounded canonical inverse from every raw
Z4c state to that auxiliary gauge.  There is a precise quotient-space route
that avoids making that unsupported assertion.

\paragraph{Realization lemma.}
Let $\mathcal E$ be an invariant Banach phase of compatible extended Cauchy
states, let $T(t)$ be its strongly continuous linear evolution, and let
$\mathfrak R:\mathcal E\to\mathcal D'(B)$ be a bounded reconstruction map
to the full raw Z4c Cauchy state.  Suppose:
\begin{enumerate}
\item reconstructed solutions satisfy the original linearized bulk equations,
the chosen elliptic shift, and the chosen scri lapse/shift conditions;
\item $\ker\mathfrak R$ is invariant under $T(t)$.
\end{enumerate}
Then $\mathcal X=\mathcal E/\ker\mathfrak R$ is a Banach space of raw states,
with the explicit norm
\[
 \|U\|_{\mathcal X}
 =\inf\{\|e\|_{\mathcal E}:\mathfrak R e=U\}.
\]
It supports the strongly continuous induced evolution
$S(t)[e]=[T(t)e]$, and
$\|S(t)\|\leq\|T(t)\|$.
Indeed boundedness into distributions makes the kernel closed; invariance
makes the definition independent of the representative.  Quotienting the
strong-continuity estimate proves strong continuity.  This is a genuine
same-space assertion in a nonlocal graph norm.  It asserts no equivalence to
an unweighted raw $H^s$ norm.

In the present problem kernel invariance can be proved without assuming the
desired full raw uniqueness theorem.  A zero reconstructed initial state has
zero subsidiary and physical Weyl data.  Their independent uniqueness gives
zero subsidiary and Weyl solutions.  The total reconstructed physical metric
is therefore a pure physical gauge perturbation.  It obeys the homogeneous
final gauge problem and has zero raw initial Cauchy state.  The independent
pure-gauge uniqueness theorem then makes its raw state zero for all times.
This argument requires that the extended compatibility conditions identify
the reconstructed $H,M,T,z,d$ with their geometric definitions initially;
one may not allow unrelated extra curvature data in the kernel.

Here is an explicit rigidity argument for the possible initial gauge kernel.
The physical initial slice is hyperbolic three-space with
$\gamma_p$ of sectional curvature $-l^2$ and $K_{ij}=-l\gamma_{p,ij}$.
Decompose a pure physical gauge generator initially as
$X=F n_p+\eta$, where $F=A\tau/\Omega$.  If its raw three-metric and
extrinsic-curvature variations vanish, the initial-data variation formulae
give
\[
 D^p_{(i}\eta_{j)}=-lF\gamma_{p,ij},\qquad
 D^p_iD^p_jF=l^2F\gamma_{p,ij}.
\]
For verification, the normal extrinsic-curvature variation is
$-D^p_iD^p_jF-l^2F\gamma_{p,ij}$; the tangential term
$\mathcal L_\eta K=2l^2F\gamma_p$ supplies the other sign.  The trace is
$\delta K=-(\Delta_p-3l^2)F$ once $\delta\gamma=0$.
The Hessian equation has exactly four solutions, the restrictions of the
ambient Minkowski Cartesian coordinates to the hyperboloid.  Indeed it
determines $(F,DF)$ along every geodesic from its value at one point, and
these four independent restrictions realize all such initial values.
Therefore
\[
 \eta=-\frac1l\nabla_pF+\eta_K,
\]
where $\eta_K$ is one of the six hyperbolic Killing fields.  These ten data
are exactly the restrictions of Minkowski Killing generators.  Subtract
that global physical Killing field, which has identically zero raw metric
variation.  The remaining generator has $\tau(0)=\eta(0)=0$.
The vanishing raw shift then gives $\eta_t(0)=0$; the vanishing raw lapse,
using $u=\Omega D_0f+aG\cdot(D_0\eta-l\eta)$, gives $f_t(0)=0$ in the
interior, while $f(0)=0$.  The homogeneous DN and lapse-wave uniqueness
results make the remaining raw perturbation identically zero.  This proves
kernel invariance without assuming uniqueness of the full hybrid system.

Applying the pure-gauge estimate also requires regularity of this extracted
generator, not merely its existence on interior compact subsets.  The
trace-free Korn inequality on the Euclidean ball gives
\[
 \|\eta-\Pi_{\rm CKV}\eta\|_{H^{s+1}}
 \leq C_s\|2(\partial_{(i}\eta_{j)})^{\rm TF}\|_{H^s},
\]
with its finite-dimensional conformal-Killing kernel.  The basic estimate is
proved in \url{https://arxiv.org/abs/gr-qc/0505022}; Cartesian differentiation
and the same finite-kernel normalization give the displayed integer-order
version.  It also applies to an interior distributional potential: exhaust
the ball by concentric balls of radii bounded below, fix the kernel on one
inner ball, and use their uniform scaled estimates.  Thus $h_P\in H^5$
gives an $H^6$ spatial potential.  The remaining finite-dimensional gauge
coefficients are fixed by the lapse and shift components.
In the present variables the exact identity
\[
 c=-\tfrac23\operatorname{div}\eta-2lf
\]
recovers $f$ with no division by $\Omega$.  After homogeneous DN uniqueness
has given $\eta=0$, the regular lift gives $f|_{\scri}=0$, and Hardy's
inequality gives $f/\Omega\in L^2$.  The lapse identity then gives the
required $f_t\in L^2$ from the zero-trace lapse and the available ordinary
derivatives.  Thus the kernel solution lies in the wave energy class used
for its uniqueness, rather than relying only on an interior gauge formula.

\paragraph{How an invariant extended phase can be built.}
The elementary cascade principle is useful: if $A$ generates a $C_0$
semigroup on $P$, $B$ generates one on $F$, and
$K:D(A^j)\to F$ is bounded in its graph norm, then
\[
 p_t=Ap,\qquad f_t=Bf+Kp
\]
is a bounded triangular perturbation on $D(A^j)\times F$ and hence has a
$C_0$ semigroup there.  Its own graph powers are invariant and encode the
forced corner compatibility.  It is generally wrong to replace those
coupled graph powers by the uncoupled product $D(A^k)\times D(B^k)$:
the forcing may cancel a singular term in $Bf$.

For the DN boundary variable, one derivative of apparent forcing loss has an
explicit removal.  Write
\[
 b_t=B_{\rm DN}b+F(p),\qquad F(p)\in H^{9/2}(S^2).
\]
Remove the finitely many low modes and let
$K=(-B_{\rm DN})^{-1}$ on the remaining modes.  It maps
$H^{9/2}\to H^{11/2}$.  Set $b=\bar b+KF(p)$.  Then
\[
 \bar b_t=B_{\rm DN}\bar b+P_{\rm low}F(p)-K F_t(p).
\]
The low-mode term is finite dimensional; if $p\in D(A_p)$, the other term
is bounded in $H^{11/2}$.  Thus the lifted boundary variable has an invariant
$H^{11/2}$ phase, and its regular Lam\'e extension has $H^6$ interior
regularity.  Taking four graph powers of the resulting upstream generator
provides its time jets without demanding four additional ordinary spatial
derivatives of the particular metric.

\paragraph{What must still be checked before applying the lemma.}
The forced lapse two-jet lift must give a bounded map from the stated
upstream graph phase into the regular lapse forcing space, including the
physical source terms.  The raw reconstruction must be bounded in the chosen
realization topology, including all curvature and scri-trace variables, not
just the conformal metric.  A bounded map into distributions is enough for
the abstract quotient result, but is not enough to advertise a classical
regular-scri theorem.  If stronger traces are claimed, their embedding must
be proved for the \emph{coupled} lapse graph domain.  Finally all initial
algebraic, geometric and gauge compatibility conditions must form a closed
invariant subspace of the extended phase.  These are concrete remaining
checks; a formal triangular construction alone does not establish them.

The subsidiary fractional branch also prevents replacing this audit by an
unstated smooth-data theorem.  The allowed natural exponent
$176/31$ gives ordinary $H^6$ regularity but generally not $H^7$.
A Fr\'echet version would therefore need a conormal/polyhomogeneous scale and
a finite-seminorm estimate for every output seminorm.  A single finite
$G_{80}$ estimate does not establish that infinite family.

For the finite repaired phase, ``joint generators'' always includes
cutoff Cartesian derivatives in the interior, in addition to rotations and
$\Omega\partial_\Omega$ on the collar.  Thus subsidiary order $128$ includes
ordinary interior order $128$, the CEFE joint order $32$ includes its mixed
ordinary interior orders, the elliptic spatial gauge joint order $20$ with
base $H^6$ includes interior $H^{26}$, and wave joint order $12$ includes
ordinary interior order $12$ on its first-order energy fields.  The source
estimates must use these same generators.  Conormal collar control alone
must not be substituted for these interior derivatives.

% Assembled finite-phase linear theorem. Final audit is recorded separately.
\section{Assembled linear scri closure and fixed-phase evolution}
\label{sec:assembled-linear-closure}

This section assembles the preceding calculations for the full
\emph{linearized} problem selected by the user. It does not equate the
nonlocal phase below with all raw $H^4/H^3$ data. Auxiliary conformal
fields are used for the proof and the norm; they do not replace the
freely evolved Z4c variables in the proposed scheme.

\subsection{The formulation and boundary prescription}
Use the exact conformal Minkowski background on the ball and retain
the $C_{\rm Z4c}=0$ tensor evolution. Match the final Appendix-C scalar
equation by the disclosed addition $3\alpha w\Theta_{\rm phys}/\Omega$.
Choose
\[
 a\kappa_1=4,\qquad \kappa_2=-99/100
\]
globally, or retain these constants in the collar and account for any
smooth interior changes in bounded coefficients. The formulae recorded
here use the global constants. No mean-curvature constraint is imposed.

Keep the original reference lapse, with zero additional restoring
coefficient and the reference $1+\log$-like principal family:
\[
\boxed{
u_t=B\cdot\nabla u-G\cdot v+
\frac{-Fq-AG\cdot v+Aw_*u}{\Omega},\qquad
F=A^2+\nu(r)\Omega^2,\quad \nu\geq0 .
}
\]
The prescribed profile is smooth, radial, stationary and bounded with its
needed derivatives. The fork's background-evaluated $1+\log$ variant
is included by $\nu=\nu_{\rm input}(2aS)^2A$.
The freely evolved trace is $q=\delta(K_{\rm phys}-2\Theta_{\rm phys})$,
not a physical-$K$ lapse.

Replace the shift evolution by metric-derived freezing
$V_t=-\eta_D V$, with constant $\eta_D\geq0$:
\[
\boxed{
Lv_i=-B\cdot\nabla V_i+(a^{-1}-\eta_D)V_i
       +2A\partial^ja_{ij}-2a_{ij}G^j,\qquad
L=\Delta+\tfrac13\nabla\operatorname{div}.
}
\]
At scri prescribe
\[
\boxed{
u=0,\qquad v_A=0,\qquad
v_n=-u+\frac A2(h_{nn}-c).
}
\]
The added $u=0$ is the boundary clock condition missing from nullness alone.
The initial compatibility class below implies and preserves
$h_{nn}-c=0$, hence the total shift boundary value is zero.
These are gauge conditions, not incoming physical radiation data.
Regularity at the center supplies the interior endpoint condition.

\subsection{Explicit finite phase}
Use a fixed finite generating family $\mathcal Z$ of spatial vector fields
tangent to scri: collar dilation
$\mathcal D_\Omega=\Omega\partial_\Omega$, the rotations with their tensor
Lie action, and cutoff Cartesian derivatives in an overlapping interior
region. The collar dilation is smoothly cut off before the center.
This family spans all ordinary interior derivatives and all boundary
tangential derivatives; it does not pretend to include the ordinary
outer normal derivative at scri.

The constraint phase $\mathcal C_{128}$ includes all actual time
generator derivatives and all words in $\mathcal Z$, of total order
at most $128$, in the twenty-component common $L^2$ energy.
Conormal derivatives are part of the initial phase from the outset.
The earlier assertion that time/rotation derivatives alone create this
normal regularity is withdrawn. The direct first-order commutator
argument in the repaired subsidiary sections proves invariance.

Use the following finite compatible extended phase $\mathcal E$:
\begin{enumerate}
\item $\mathcal C_{128}$, together with the metric-connection memory
$V$ in joint order $32$ at ordinary $H^4$.
\item The compatible rescaled Weyl and sourced conformal particular
Cauchy fields in joint time/$\mathcal Z$ order $32$ at the mixed
$m=4$ phase: metric $H^5$ with velocity $H^4$,
Schouten $H^4$ with velocity $H^3$, conformal-factor perturbation
$H^6$ with velocity $H^5$, scalar $s$ in $H^5$ with velocity $H^4$,
and rescaled Weyl in $H^4$.
\item The elliptic displacement $\eta$ in joint order $20$ at $H^6$,
equivalently its compatible Lam\'e source and affine boundary-DN
Cauchy jets in $H^{11/2}$. The high-frequency inverse of the boundary
generator supplies the one-derivative forcing lift described below.
\item The lifted lapse-wave energy in affine joint order $12$.
The lifts are the four \emph{derived} coefficients
$f_0,f_1,f_2,f_3$, not independently supplied boundary data.
\end{enumerate}
Time jets throughout are those of the full triangular equations, including
known source derivatives. They are not arbitrary additional functions.
The generalized harmonic and scalar-curvature gauges of the auxiliary
particular solution are fixed once.

The compatibility conditions are: the determinant/trace algebraic Z4c
conditions; the differential reduction identities; geometric identification
of $H,M,T,z$ and the initial rescaled Weyl; the sourced Gauss--Codazzi and
conformal curvature identities; elliptic memory/trace compatibility;
and the finite affine time and lift corner identities.
Initial regular rescaled Weyl and Schouten data are explicit peeling
requirements. They do not follow from $H,M,T,z$ weights alone.
The independent regular $s$ quotient is likewise included.

Define $\mathcal E$ as the completion in this finite sum norm of compatible,
interior-smooth conormally regular Cauchy tuples, allowing the natural
$\Omega^{176/31}$ constraint branch. No conformal $C^\infty$ restriction
or unproved infinite ordinary-normal regularity is imposed.
The source jets are determined by the upstream tuple.
The identities are closed distributional/trace conditions in the stated
finite norm; the augmented evolution preserves them.

Let $\mathfrak R$ be the explicit physical-plus-gauge reconstruction of
all raw Z4c fields. Set
\[
 \boxed{
 \mathcal X=\mathcal E/\ker\mathfrak R,\qquad
 \|U\|_{\mathcal X}=
   \inf_{\mathfrak Re=U}\|e\|_{\mathcal E}.
 }
\]
This is an initial-state norm, using the stated finite Cauchy tuple and
elliptic solves. It is not defined by a future solution supremum.
A bounded projection onto the kernel is unnecessary.

\subsection{Coupled linear theorem and proof}
For every $U_0\in\mathcal X$ there is a unique reconstructed solution
$U\in C([0,\infty);\mathcal X)$ of the full linearized formulation above,
and for every finite $T$
\[
\boxed{
 \sup_{0\leq t\leq T}\|U(t)\|_{\mathcal X}
 \leq C_T\|U_0\|_{\mathcal X}.
}
\]
The following proof specifies the reconstruction and the estimates that
make the fixed-phase assertion meaningful.

\paragraph{Subsidiary evolution and falloffs.}
The common principal/residue energy, capped construction and homogeneous
reduction-error transports give the subsidiary evolution.
The repaired joint-conormal commutator and simultaneous weight bootstrap
preserve
\[
 H,T=\Omega^4(\hbox{natural }H^6\hbox{ fields}),\qquad
 M,z=\Omega^3(\hbox{natural }H^6\hbox{ fields}).
\]
The three outgoing leading amplitudes obey
\[
 H_4=-2T_4/(25a),\quad
 M_{3n}=302T_4/(75a),\quad
 Z_{3n}=-2T_4/75,\quad M_{3A}=-Z_{3A}/a.
\]
Constraint violations are retained. No Hamiltonian, momentum, scalar-Z4
or connection projection is used.

\paragraph{Physical curvature and a particular metric.}
These fields produce the conserved Ricci source
$\widetilde R_{ab}=\Omega^2r_{ab}$ and the regular Cotton source $J$.
The full ten-component Bianchi energy has nonpositive scri flux.
The traceful sourced conformal waves and the interior-equivalence argument
construct a regular particular physical metric.
The bounded coordinate/conformal adjustment puts its conformal factor
back to the prescribed $\Omega$. The first-order conormal commutator
also propagates its finite joint phase.
Gauss/Codazzi and $\widetilde R_{ij}=-Q^{\rm phys}_{ij}/N$
then identify its geometric constraints and ADM evolution with the
required off-constraint Z4c source.

\paragraph{Elliptic displacement and boundary clock.}
Solve
\[
 L\eta=V_{\rm memory}-V(h_P).
\]
Its boundary generator is the exact spherical Lam\'e DN operator with
the added clock forcing $u_P$ and tangential forcing $-(v_P)_A$.
The polar high modes have a uniform positive symmetrizer; all
$\ell\ge2$ modes decay, and the finite neutral modes have only polynomial
growth. Write its high-frequency inverse as $K$. The substitution
$b=\bar b+KF(P)$ changes the boundary forcing to
$-KF_t(P)$, gaining its required derivative.
The low modes are finite dimensional.
Thus the order-$20$, $H^6$ elliptic phase is bounded by the order-$32$
particular phase and its compatible initial data.

\paragraph{Forced lapse and explicit pole cancellation.}
For $H_\eta=\eta\cdot G$ and
$f=(A\tau-aH_\eta)/\Omega$, the exact pure-gauge identity is
\[
 \mathcal L_u(U_g)
 =\Omega\left[P_\nu\phi-\frac{2CAH_\eta}{a\Omega^3}\right],
 \qquad \phi=A\tau/\Omega.
\]
If $\mathcal L_u(U_P)=R_P+S_P/\Omega$, the source for $f$ has only
quadratic and linear poles. Its first lift is
$f_0=f_{0,\rm hom}-a\mu_P/2$, which gives
$\tau_n=\mu_P/(2A)$ and cancels the total normal metric residue.
The next three lifts cancel successive known source coefficients.
Writing $f_L=\sum_{j=0}^3\Omega^jEf_j$,
\[
 P_\nu(f-f_L)=F_3,\qquad F_3/\Omega^2\in L^2.
\]
This uses the anisotropic fourfold normal Hardy inequality on the full
numerator, with its first four normal jets zero.
The bulk normal requirements are only $S_P\in H^4$, $R_P\in H^4$,
and $\eta\in H^6$. A safe \emph{eight} additional joint orders control
boundary coefficient extensions, so $12+8=20$ fits the phase.
There is no demand for ordinary $H^7$ constraint regularity.

\paragraph{Lapse-wave existence and weighted embedding.}
After rescaling by $A$, the positive potential wave has a paired skew
$1/\Omega$ multiplier. Paired capping, weak compactness and the
$O(\Omega^2)$ glancing cutoff give its unique weak energy evolution.
The affine time/conormal hierarchy propagates the order-$12$ phase.
From the compatible equations
$R_t=c_0w/\Omega$ and the $w$ equation one recovers the successive weights
$\Omega^{-p}R$ at a cost of $2p$ joint orders and source weights through
$F_3/\Omega^{p-1}$. For $p\le3$ and five additional joint derivatives,
order $11$ suffices. This proves bounded raw reconstruction into
\[
 h,c,u,v\in H^4,\qquad a,q\in H^3,\qquad q/\Omega\in H^2,
\]
with the retained constraint and connection fields in their stated
regular classes. The extra order controls the needed time derivatives.

\paragraph{One fixed raw-state space.}
The compatible extended hierarchy is a bounded triangular evolution on
the complete phase $\mathcal E$. Its reconstruction is continuous into
distributions, hence its kernel is closed.
A zero raw initial state has zero independent subsidiary and Weyl data.
Their uniqueness makes the reconstructed physical metric flat to first
order, hence pure physical gauge on the simply connected domain.
The explicit hyperbolic KID equations
\[
 D^p_{(i}\eta_{j)}=-lF\gamma^p_{ij},\qquad
 D^p_iD^p_jF=l^2F\gamma^p_{ij}
\]
give precisely the ten Minkowski Killing initial data.
Subtracting that physically trivial Killing field gives zero lapse-wave
and displacement Cauchy data. Homogeneous DN and positive-wave uniqueness
then make the reconstructed raw solution zero.
Thus $\ker\mathfrak R$ is invariant, without assuming the desired raw
uniqueness theorem. The quotient carries the strongly continuous
evolution and the displayed same-space estimate.

\subsection{Actual scri closure and the limits of the theorem}
The preceding estimates imply, for all times,
\[
 T|_{\scri}=z|_{\scri}=H|_{\scri}=M|_{\scri}=0,\qquad
 u|_{\scri}=v|_{\scri}=0,\qquad
 (h_{nn}-c)|_{\scri}=q|_{\scri}=0.
\]
For each combined pole numerator $S_U$, define its regular quotient
by the reconstructed identity
\[
 Q_U=(U_t-R_U)=S_U/\Omega\quad(\Omega>0).
\]
The identity and weighted recovery give $Q_U\in H^2$.
Its trace, rather than separate divergent summands, is the scri RHS.
Since $\partial_n\Omega=-1/a$ at the reference sphere, the explicit
boundary update for the geometric and Z4c fields is
\[
 \boxed{\quad U_t|_{\scri}=(R_U-a\partial_nS_U)|_{\scri}.\quad}
\]
The lapse uses $u_t|_{\scri}=0$ and the elliptic shift uses its stated
Dirichlet values. The $S_U$ are the spatial pole numerators in the
companion full linearization. They contain no unknown $U_t$.
The derivative formula holds in the Sobolev trace sense, including the
connection block; the proved quotient has the stronger $H^{3/2}$ trace.
Thus the reconstruction identity proves regularity but is not needed
to evaluate an unknown time derivative in the boundary update.
This defines the unique regular limiting boundary equation and preserves
its finite corner hierarchy. No incoming constraint or physical radiation
function is supplied at scri.

The conclusion is a full coupled \emph{linear} evolution in the explicit
compatible nonlocal quotient graph phase. It is not an assertion that
arbitrary raw $H^4/H^3$ data have a norm-equivalent estimate.
It proves neither nonlinear binary-black-hole evolution nor stability
of a particular discretization. Those statements are outside the
linearized scope selected here.


% Adversarial finite-phase audit; input fragment.
\section{Independent finite-phase completion audit}
\label{sec:linear-completion-certificate}

This certificate concerns the repaired \emph{linearized} system and its
explicit nonlocal graph phase.  It does not certify the former
time/rotation-only subsidiary embedding, which has been withdrawn, or an
ordinary raw-$H^4/H^3$ theorem for arbitrary data.

\paragraph{The phase actually audited.}
Let $\mathcal Z$ consist of collar radial dilation
$\Omega\partial_\Omega$, covariant rotations, and overlapping cutoff
Cartesian derivatives in the interior.  All words in these generators are
included.  The extended phase is the completion of compatible Cauchy tuples
in the sum of the following norms:
\begin{enumerate}
\item the subsidiary joint norm $\mathcal G_{128}$, with conormal
regularity present initially;
\item the regular Bianchi field joint norm through order $32$ in $H^4$,
and the full CEFE mixed phase $\mathcal E_{4,32}$, including its
$\omega\in H^6$, $s\in H^5$ components and the listed velocities;
\item the prescribed Dirac memory in the corresponding finite joint
$H^4$ norm, and the elliptic spatial gauge joint norm through order $20$
in $H^6$, using its affine DN time jets;
\item the four-lift gauge-wave first-order energy norm through joint
order $12$, with affine jets
$U_{p+1}=\mathcal A_\nu U_p+\partial_t^pF_3(0)$.
\end{enumerate}
The forcing jets are determined by the upstream Cauchy tuple, not supplied
as independent future histories.  All algebraic and geometric definitions,
initial Weyl Gauss--Codazzi identities, CEFE/gauge identities, elliptic
boundary compatibility, derivative reductions and finite lift identities
are imposed as closed linear conditions.  The core allows the actual
fractional branch $176/31$ and is not restricted to conformally smooth
boundary data.  Interior derivatives are included by the generator
definition; no unlisted high interior regularity is invoked.

\paragraph{Requirements checked.}
\begin{enumerate}
\item \emph{Formulation and source disclosure.}  The system is the stated
tensor $C_{\rm Z4c}=0$ linearization with the explicit scalar correction
$3AwT/\Omega$, $a\kappa_1=4$, $\kappa_2=-99/100$, coordinate-time
metric-Dirac memory, the original reference lapse with zero additional
restoring parameter, and $u|_{\scri}=0$ plus the stated scri shift.
The audited principal coefficient extension has stationary smooth
\emph{radial} $\nu\ge0$, $F=A^2+\nu\Omega^2$.
The scalar correction changes the off-constraint equations and is disclosed;
no CMC projection or removal of differential constraint variables is used.

\item \emph{Constraint evolution and falloffs.}  The twenty-component common
energy, weak construction, derivative-reduction transport and the repaired
joint-generator propagation supply a genuine invariant subsidiary phase.
The simultaneous fast/slow weight induction starts with the needed
conormal norms already present, so it does not infer an unknown
$\mathcal DR_g$ from its own time derivative.  Forty joint orders suffice
for the natural $H^6$ falloffs; the larger budget supplies the commuted
physical sources.  Nonzero differential constraint errors are retained.

\item \emph{Physical particular metric.}  The Ricci source obeys the exact
physical conservation identity, including its trace.  The modified-$s$
traceful sourced CEFE have regular coefficients and a mixed wave estimate.
Interior comparison with conserved sourced linear physical Einstein,
Bianchi uniqueness and reduced-wave uniqueness establish the full CEFE
identities without importing a trace-free-matter theorem.  The bounded
fixed-$\Omega$ transformation yields a particular metric in $H^5$ with
curvature in $H^4$.  Initial curvature and Weyl regularity are explicit
compatibility data, not consequences claimed from bare raw $H^4$.

\item \emph{All angular sectors and boundary uniformity.}  The Weyl and
subsidiary systems keep their scalar/vector sectors.  The global ball
Lam\'e DN spectrum treats every polar and axial harmonic; its high-mode
symmetrizers are uniformly comparable to the identity.  The finite
zero/low-mode Jordan block has only finite-time polynomial growth.
There is therefore no hidden high-angular-frequency growth in the DN
estimate.  The one-order DN forcing lift is supplied by upstream time
regularity, and ordinary elliptic estimates give the interior $\eta$.

\item \emph{Forced lapse and normal recovery.}  The four coefficients
$f_0,\ldots,f_3$ are derived from the triangular nonzero $V_0|_{\scri}$
recursion.  An anisotropic Hardy-four estimate gives
$F_3/\Omega^2\in L^2$ from four normal derivatives of the particular pole,
five metric/lapse derivatives and six spatial-gauge derivatives, with
separate boundary time/angular norms.  A conservative eight additional
joint upstream orders fits $12+8=20$.
The skew-potential wave energy propagates the affine hierarchy.
Algebraic weight recovery through $p=3$ then gives the worst raw term
$\Omega^{-3}\mathcal D^5R$ with cost $6+5=11<12$.
Hence raw metric/lapse/shift are $H^4$, curvature is $H^3$,
$q/\Omega\in H^2$, and every analytically combined pole RHS has $H^2$
regularity and its corresponding finite scri trace.  This is not an
estimate of divergent summands separately.

\item \emph{Same-space realization.}  The finite cascade estimates preserve
the complete extended phase and are continuous at time zero.  Let
$\mathfrak R$ be the bounded raw reconstruction and
$K=\ker\mathfrak R$.  Closedness follows from boundedness into distributions.
Kernel invariance follows from subsidiary/Weyl uniqueness, flat physical
gauge rigidity, the ten-dimensional Minkowski KID calculation and the
homogeneous gauge uniqueness theorem.  Conformal Korn and the exact
$c=-2\operatorname{div}\eta/3-2lf$ identity place this kernel gauge in the
required energy class.  Thus
\[
 \mathcal X=\mathcal E/K,\qquad
 \|U\|_{\mathcal X}
  =\inf_{\mathfrak Re=U}\|e\|_{\mathcal E}
\]
is a complete raw-state phase with a strongly continuous linear evolution
and a finite-time estimate in this \emph{same} norm.  No bounded canonical
raw-to-CEFE inverse or ordinary Sobolev norm equivalence is asserted.
\end{enumerate}

The substantive issue found during this audit was the former
time/rotation-only subsidiary conormal claim.  Its explicit initial-domain
repair is incorporated in both copies of the existence and normal-bootstrap
sections.  With that repair, interior generators, radial $\nu$, and the safe
finite source count above, no further contradictory mode or derivative-count
obstruction was found in the stated linear quotient-graph realization.
No nonlinear, black-hole, discretization or arbitrary-raw-Sobolev theorem
follows from this certificate.

% Independent final audit by the scri-linearization agent.
\section{Independent audit of the final source and quotient-phase closure}
\label{sec:independent-final-source-phase-audit}

The files audited are the current assembled linear theorem, sourced CEFE
reconstruction, phase-quotient audit, original-gauge ball reduction,
subsidiary conormal-phase repair, and completion certificate.  The audit
checks their interfaces and the source, compatibility, derivative-count,
and kernel arguments.  It does not repeat the independently recorded
rational common-energy computation.  The conclusion is limited to the
explicit compatible finite quotient graph phase in those files.

\paragraph{Traceful source and conformal identities.}
The physical conservation normalization is correct.  If
$\widetilde R_{ab}=\Omega^2r_{ab}$ and
$\ell=(r-g\operatorname{tr}r/6)/2$, then
$r-g\operatorname{tr}r/2=2\ell-2g\operatorname{tr}\ell$ and
\[
 \widetilde\nabla^a\bigl[\Omega^2(r_{ab}
 -g_{ab}\operatorname{tr}r/2)\bigr]
 =2\Omega^4(\nabla^a\ell_{ab}-\nabla_b\operatorname{tr}\ell)
 -6\Omega^3\Sigma_b\operatorname{tr}\ell.
\]
Thus the displayed traceful conservation equation, including its factor
three, is the physical contracted Bianchi identity.  It is essential that
this follows from the closed subsidiary $H,M$ evolution; arbitrary smooth
external $r$ would not suffice.  Its boundary trace restriction agrees with
the leading subsidiary coefficient relation and is not an extra demand
that the bulk differential constraints vanish.

The sourced Hessian, scalar, curl, and divergence equations have the stated
normalizations.  In particular, differentiating
$2\Omega s-\Sigma^2=0$ and substituting the Hessian gives exactly
$\nabla s=-L\Sigma+\Omega^2\ell\Sigma$.
Tracing the Hessian gives
$\Box\Omega=4s-R\Omega/6+\Omega^3\operatorname{tr}\ell$.
The Cotton connection calculation gives the stated regular source $J$;
no trace-free-matter hypothesis was silently used.

\paragraph{Particular-metric construction and initial equivalence.}
The mixed wave orders are consistent: $\delta L$ at $H^4$ needs its wave
forcing at $H^3$; the second metric derivatives in that forcing have this
order when $\delta g\in H^5$.  Metric forcing is $H^4$, scalar-$s$
forcing is $H^4$, and the conformal-factor forcing is $H^5$ in the common
contracted-Christoffel gauge.  This last gauge reduction is necessary for
$\omega\in H^6$ and the subsequent bounded fixed-$\Omega$ adjustment.

The interior comparison with conserved sourced physical Einstein is a
legitimate way to prove the sourced CEFE zero quantities, provided the
initial tuple has the full sourced curvature/Hessian/gauge compatibility
listed in the phase.  Such compatibility fixes the Schouten and scalar
normal jets; they are not arbitrary auxiliary Cauchy functions.  The four
physical gauge values and four sourced Einstein constraints provide the
zero Lorenz-gauge data.  Compatible initial Weyl data identify the genuine
and independently evolved Bianchi fields.  The backward domain of an
interior point does not reach scri in finite time because the outward
backward glancing speed is $O(\Omega^2)$.  Hence this comparison does not
import unspecified exterior data or an unproved traceful-matter CEFE
subsidiary theorem.

There is no requirement $H=M=0$ in this argument: the prescribed Einstein
source has $G_{nn}=H/2$, $G_{ni}=-M_i$, and the specified spatial stress.
Initial data compactly supported away from scri give a nontrivial
infinite-dimensional off-constraint compatible class; the phase is not
merely vacuum data renamed as a sourced system.

\paragraph{Finite source and graph budgets.}
The repaired phase includes conormal derivatives initially and cutoff
Cartesian generators in the interior.  This removes both the earlier
time-only normal-embedding circularity and the missing interior Sobolev
orders.  Radial conormal commutation cancels the apparent
$1/\Omega$ coefficient derivative as displayed in the repair; it does not
assume cap estimates that would fail uniformly.

The interfaces $128\to32\to20\to12$ are consistent with the supplied
finite source estimates.  In particular the four-lift Hardy argument uses
four ordinary normal derivatives of the complete pole numerator, hence
only four ordinary normal derivatives of $S_P$, five of metric/lapse,
and six of the displacement.  The additional time/angular trace
derivatives are charged to the joint budget.  Reserving eight extra joint
orders, rather than the earlier optimistic six, fits $12+8=20$.
No ordinary seventh derivative of the fractional subsidiary branch is
invoked.  The algebraic wave recovery uses at most $2p+m=11$ joint orders
for $p\le3$, $m\le5$, with the supplied source weights
$F_3/\Omega^2$.  These bounds control the asserted raw metric $H^4$,
curvature $H^3$, and combined pole RHS $H^2$ classes.

The particular pole and lapse residual estimates depend on actual
compatible time jets.  Replacing them by uncoupled powers of the
homogeneous gauge generator would be incorrect.  The current affine
hierarchy explicitly avoids that replacement.

\paragraph{Closed phase and kernel.}
The quotient lemma is mathematically sufficient for a same-space result;
a canonical bounded raw-to-CEFE inverse is not needed.  Each compatibility
condition must be read in the declared extended norm: multiplication by
$\Omega$, geometric differential identities, the finite trace jets, and
the affine generator identities are continuous distributional/trace
conditions there.  Completing the compatible core therefore preserves
them as closed conditions.  The separately proved cascade estimates
preserve this completed phase, not merely a weaker raw output norm.

A zero raw initial state forces zero subsidiary and initial Weyl data by
their geometric identification.  Their uniqueness makes the total
reconstructed physical metric linearized-flat.  On the simply connected
domain it is pure physical gauge.  Vanishing initial metric and extrinsic
curvature give the stated hyperbolic KID equations.  The Hessian equation
has four solutions, and the residual spatial Killing field has six; these
are precisely the ten physically trivial Minkowski Killing generators.
Subtracting that global Killing field leaves zero displacement and normal
gauge Cauchy values.  Vanishing raw shift/lapse then gives zero initial
time data.  The established homogeneous gauge uniqueness proves kernel
invariance rather than assuming the desired raw uniqueness theorem.
The conformal Korn estimate and the pole-free identity for $f$ place this
extracted potential in the needed gauge uniqueness class.

\paragraph{Conclusion of this audit.}
No additional concrete source inconsistency, forbidden bulk constraint
projection, contradictory boundary mode, or derivative-count obstruction
was found in the repaired finite quotient-phase argument.  The earlier
counterexamples concern different lapse closures or missing clock data
and do not contradict this specified original-reference, clock-closed
formulation.  The conclusion is a linear theorem on
$\mathcal E/\ker\mathfrak R$, assuming the previously audited subsidiary,
DN, and weak-wave lemmas recorded in the accompanying source files.
It is not norm equivalence with arbitrary raw $H^4/H^3$ data, not a
nonlinear theorem, and not numerical stability of an implementation.


\section*{Final completion audit for the selected linear scope}
\begin{center}
\begin{tabular}{p{.65\linewidth}p{.23\linewidth}}
Explicit formulation and mapping to the reference equations & established\\
Free differential constraint modes and invariant falloffs & established in the joint phase\\
All angular sectors and global ball elliptic boundary estimate & established\\
Traceful sourced physical metric and full Z4c reconstruction & established\\
Four derived lifts and regular actual-scri pole traces & established\\
Existence, uniqueness and continuous dependence in one stated norm & quotient graph-phase theorem\\
Independent source, derivative-count and phase audits & completed with repairs recorded
\end{tabular}
\end{center}
The selected full linearized objective is covered in the stated compatible phase. No production evolution equations have been changed, and no numerical simulation has been used as evidence of the theorem. Nonlinear evolution, arbitrary raw Sobolev norm equivalence and discretization stability are not asserted.
\end{document}
